% -------------------------------------------------------------
% NOTE ON THE DETAILED AND SHORT VERSIONS:
% -------------------------------------------------------------
% This paper comes in two versions, a detailed and a short one.
% The short version should be more than sufficient for any
% reasonable use; the detailed one was written purely to
% convince the author of its correctness.
% To switch between the two versions, find the line containing
% "\newenvironment{noncompile}{}{}" in this LaTeX file.
% Look at the two lines right beneath this line.
% To compile the detailed version, they should be as follows:
%   \includecomment{verlong}
%   \excludecomment{vershort}
% To compile the short version, they should be as follows:
%   \excludecomment{verlong}
%   \includecomment{vershort}
% As a rule, the line
%   \excludecomment{noncompile}
% should stay as it is.
% -------------------------------------------------------------
% NOTES ON SOME HACKS USED IN THIS FILE:
% -------------------------------------------------------------
% One of my pet peeves with amsthm is its use of italics in the theorem and
% proposition environments; this makes math and text indistinguishable in said
% enviroments. To avoid this, I redefine the enviroments to use the standard
% font and to use a hanging indent, along with a bold vertical bar to its
% left, to distinguish these environments from surrounding text. (Along with
% the advantage of distinguishing math from text, this also allows nesting
% several such environments inside each other, like a definition inside a
% remark. I'm not sure how good of an idea this is, though. There are also
% downsides related to the hanging indentation, such as footnotes out of it
% being painful to do right.) This is done starting from the line
%   \theoremstyle{definition}
% and until the line
%   {\end{leftbar}\end{exmp}}
% Draw a horizontal line from left to right margin.

\documentclass[numbers=enddot,12pt,final,onecolumn]{scrartcl}%
\usepackage[headsepline,footsepline,manualmark]{scrlayer-scrpage}
\usepackage{amssymb}
\usepackage{amsmath}
\usepackage{amsthm}
\usepackage{framed}
\usepackage{comment}
\usepackage{color}
\usepackage[breaklinks=True]{hyperref}
\usepackage{hyperxmp}
\usepackage{lmodern} % Established rendering compatibility for absent Palatino metrics.
\usepackage[T1]{fontenc}
% Unused license-logo package omitted; original CC BY4.0 declaration retained in item license/evidence.
\usepackage{tikz}
\usepackage{needspace}
\usepackage{tabls}
%TCIDATA{OutputFilter=latex2.dll}
%TCIDATA{Version=5.50.0.2960}
%TCIDATA{LastRevised=Monday, April 26, 2021 13:56:16}
%TCIDATA{SuppressPackageManagement}
%TCIDATA{<META NAME="GraphicsSave" CONTENT="32">}
%TCIDATA{<META NAME="SaveForMode" CONTENT="1">}
%TCIDATA{BibliographyScheme=Manual}
%TCIDATA{Language=American English}
%BeginMSIPreambleData
\providecommand{\U}[1]{\protect\rule{.1in}{.1in}}
%EndMSIPreambleData
\usetikzlibrary{arrows}
\newcounter{exer}
\theoremstyle{definition}
\newtheorem{theo}{Theorem}[section]
\newenvironment{theorem}[1][]
{\begin{theo}[#1]\begin{leftbar}}
{\end{leftbar}\end{theo}}
\newtheorem{lem}[theo]{Lemma}
\newenvironment{lemma}[1][]
{\begin{lem}[#1]\begin{leftbar}}
{\end{leftbar}\end{lem}}
\newtheorem{prop}[theo]{Proposition}
\newenvironment{proposition}[1][]
{\begin{prop}[#1]\begin{leftbar}}
{\end{leftbar}\end{prop}}
\newtheorem{defi}[theo]{Definition}
\newenvironment{definition}[1][]
{\begin{defi}[#1]\begin{leftbar}}
{\end{leftbar}\end{defi}}
\newtheorem{remk}[theo]{Remark}
\newenvironment{remark}[1][]
{\begin{remk}[#1]\begin{leftbar}}
{\end{leftbar}\end{remk}}
\newtheorem{coro}[theo]{Corollary}
\newenvironment{corollary}[1][]
{\begin{coro}[#1]\begin{leftbar}}
{\end{leftbar}\end{coro}}
\newtheorem{conv}[theo]{Convention}
\newenvironment{convention}[1][]
{\begin{conv}[#1]\begin{leftbar}}
{\end{leftbar}\end{conv}}
\newtheorem{quest}[theo]{Question}
\newenvironment{question}[1][]
{\begin{quest}[#1]\begin{leftbar}}
{\end{leftbar}\end{quest}}
\newtheorem{warn}[theo]{Warning}
\newenvironment{conclusion}[1][]
{\begin{warn}[#1]\begin{leftbar}}
{\end{leftbar}\end{warn}}
\newtheorem{conj}[theo]{Conjecture}
\newenvironment{conjecture}[1][]
{\begin{conj}[#1]\begin{leftbar}}
{\end{leftbar}\end{conj}}
\newtheorem{exam}[theo]{Example}
\newenvironment{example}[1][]
{\begin{exam}[#1]\begin{leftbar}}
{\end{leftbar}\end{exam}}
\newtheorem{exmp}[exer]{Exercise}
\newenvironment{exercise}[1][]
{\begin{exmp}[#1]\begin{leftbar}}
{\end{leftbar}\end{exmp}}
\newenvironment{statement}{\begin{quote}}{\end{quote}}
\newenvironment{fineprint}{\begin{small}}{\end{small}}
\iffalse
\newenvironment{proof}[1][Proof]{\noindent\textbf{#1.} }{\ \rule{0.5em}{0.5em}}
\newenvironment{convention}[1][Convention]{\noindent\textbf{#1.} }{\ \rule{0.5em}{0.5em}}
\newenvironment{question}[1][Question]{\noindent\textbf{#1.} }{\ \rule{0.5em}{0.5em}}
\fi
\let\sumnonlimits\sum
\let\prodnonlimits\prod
\let\cupnonlimits\bigcup
\let\capnonlimits\bigcap
\renewcommand{\sum}{\sumnonlimits\limits}
\renewcommand{\prod}{\prodnonlimits\limits}
\renewcommand{\bigcup}{\cupnonlimits\limits}
\renewcommand{\bigcap}{\capnonlimits\limits}
\setlength\tablinesep{3pt}
\setlength\arraylinesep{3pt}
\setlength\extrarulesep{3pt}
\setlength\textheight{22.5cm}
\setlength\textwidth{14.8cm}
\newenvironment{verlong}{}{}
\newenvironment{vershort}{}{}
\newenvironment{noncompile}{}{}
\includecomment{verlong}
\excludecomment{vershort}
\excludecomment{noncompile}
\newcommand{\defn}[1]{{\color{darkred}\emph{#1}}}
\newcommand{\CC}{\mathbb{C}}
\newcommand{\RR}{\mathbb{R}}
\newcommand{\QQ}{\mathbb{Q}}
\newcommand{\NN}{\mathbb{N}}
\newcommand{\ZZ}{\mathbb{Z}}
\newcommand{\KK}{\mathbb{K}}
\newcommand{\set}[1]{\left\{ #1 \right\}}
\newcommand{\abs}[1]{\left| #1 \right|}
\newcommand{\tup}[1]{\left( #1 \right)}
\newcommand{\ive}[1]{\left[ #1 \right]}
\renewcommand{\leq}{\leqslant}
\renewcommand{\geq}{\geqslant}
\newcommand\horiline{\noindent\makebox[\linewidth]{\rule{\textwidth}{1pt}} \newline}
\newcommand\arxiv[1]{\href{http://www.arxiv.org/abs/#1}{\texttt{arXiv:#1}}}
\iffalse
\NOEXPAND\horiline{\noindent\makebox[\linewidth]{\rule{\textwidth}{1pt}} \newline}
\fi
\ihead{The Elser nuclei sum revisited}
\ohead{page \thepage}
\cfoot{}
\begin{document}
\begin{titlepage} % Suppresses displaying the page number on the title page and the subsequent page counts as page 1
	
	\raggedleft % Right align the title page
	
	\rule{1pt}{\textheight} % Vertical line
	\hspace{0.05\textwidth} % Whitespace between the vertical line and title page text
	\parbox[b]{0.85\textwidth}{ % Paragraph box for holding the title page text, adjust the width to move the title page left or right on the page
\vspace{0.05\textheight}
{\Huge\textbf{The Elser nuclei sum}}\\[0.5\baselineskip]
{\Huge\textbf{revisited}}\\[0.8\baselineskip]
%{\Large\textbf{Volume 1 of 2}}\\[4\baselineskip] % Subtitle or further description
{\Large\textit{Version of April 26, 2021 (with minor corrections August 7, 2026)}}\\[0.5\baselineskip]
(\arxiv{2009.11527v8} + update)\\[4\baselineskip]
{\Large\textsc{Darij Grinberg}}\\[0.4\baselineskip]
\textit{Drexel University}\\[2\baselineskip]
\vspace{0.08\textheight}
\[
\begin{tikzpicture}%
[-,auto,node distance=3cm, thick,main node/.style={circle,fill=blue!20,draw}%
]
\node[main node] (1) {$v$};
\node[main node] [below left of=1] (2) {};
\node[main node] [below right of=1] (4) {};
\node[main node] [above of=1] (3) {};
\path[every node/.style={font=\sffamily\small}]
(1) -- (3) node[midway, right] {$c$}
(3) -- (2) node[midway, left] {$b$}
(2) -- (4) node[midway, below] {$f$}
(4) -- (3) node[midway, right] {$d$}
(2) -- (1) node[midway, below] {$a$}
(1) -- (4) node[midway, below] {$e$};
\path[every node/.style={font=\sffamily\small}] (1) edge
(3) (3) edge (2) (2) edge
(4) (4) edge (3) (2) edge
(1) (1) edge (4);
\end{tikzpicture}
\qquad\qquad
\raisebox{15pt}{
\begin{tikzpicture}[thick, main node/.style={circle,fill=yellow!20,draw}]
\filldraw[draw=black, fill=red!70] (-1, -1) -- (0, 1) -- (1, -1) -- cycle;
\node[main node] (a) at (0, 2.4) {$a$};
\node[main node] (b) at (-1, -1) {$b$};
\node[main node] (c) at (2, -2) {$c$};
\node[main node] (d) at (0, 1) {$d$};
\node[main node] (e) at (-2, -2) {$e$};
\node[main node] (f) at (1, -1) {$f$};
\draw (a) edge (d);
\draw (b) edge (e);
\draw (c) edge (f);
\end{tikzpicture}
}
\]
\vspace{0.09\textheight} % Whitespace between the title block and the publisher
{\noindent updated version of Oberwolfach Preprint OWP-2021-05}\\[\baselineskip] % Publisher and logo
}
\end{titlepage}


\setcounter{page}{2}

\horiline


\textbf{The Elser nuclei sum revisited}.

\bigskip

Version of 26 April 2021, \newline\phantom{u} \qquad equivalent to
\href{http://www.arxiv.org/abs/2009.11527v8}{\texttt{arXiv:2009.11527v8}}.

\bigskip

\textsc{Darij Grinberg}: \newline\phantom{u}\qquad Drexel University
\newline\phantom{u}\qquad Korman Center, Room 263 \newline\phantom{u}\qquad15
S 33rd Street \newline\phantom{u}\qquad Philadelphia PA, 19104 \newline%
\phantom{u}\qquad USA \newline\phantom{u}\qquad(temporary:) \newline%
\phantom{u}\qquad Mathematisches Forschungsinstitut Oberwolfach,\newline%
\phantom{u}\qquad Schwarzwaldstrasse 9--11\newline\phantom{u}\qquad77709
Oberwolfach\newline\phantom{u}\qquad Germany \newline\phantom{u}\qquad
\url{darijgrinberg@gmail.com}\newline\phantom{u}\qquad\url{https://www.cip.ifi.lmu.de/~grinberg/}

\bigskip\textbf{Abstract.} Fix a finite undirected graph $\Gamma$ and a vertex
$v$ of $\Gamma$. Let $E$ be the set of edges of $\Gamma$. We call a subset $F$
of $E$ \textit{pandemic} if each edge of $\Gamma$ has at least one endpoint
that can be connected to $v$ by an $F$-path (i.e., a path using edges from $F$
only). In 1984, Elser showed that the sum of $\left(  -1\right)  ^{\left\vert
F\right\vert }$ over all pandemic subsets $F$ of $E$ is $0$ if $E\neq
\varnothing$. We give a simple proof of this result via a sign-reversing
involution, and discuss variants, generalizations and refinements, revealing
connections to abstract convexity (the notion of an antimatroid) and discrete
Morse theory.

\bigskip\textbf{Keywords:} graph theory, nuclei, simplicial complex, discrete
Morse theory, alternating sum, enumerative combinatorics, inclusion/exclusion,
convexity, antimatroids. \vspace{1.0038pc}

\horiline


\noindent Future versions of this text will be available from the
author's website: \newline%
\url{http://www.cip.ifi.lmu.de/~grinberg/algebra/elsersum.pdf} \newline%
\url{http://www.cip.ifi.lmu.de/~grinberg/algebra/elsersum-long.pdf} (detailed version).

\vspace{1pc}

\horiline


Typeset with \LaTeX.

\doclicenseThis


\section*{****}

In \cite{Elser84}, Veit Elser studied the probabilities of clusters forming
when $n$ points are sampled randomly in a $d$-dimensional volume. In the
process, he found a purely graph-theoretical lemma \cite[Lemma 1]{Elser84},
which served a crucial role in his work. For decades, the lemma stayed hidden
from the eyes of combinatorialists in a physics journal, until it resurfaced
in recent work \cite{DHLetc19} by Dorpalen-Barry, Hettle, Livingston, Martin,
Nasr, Vega and Whitlatch. In this note, I will show a simpler proof of the
lemma using a sign-reversing involution. The proof also suggests multiple
venues of generalization that I will explore in the later sections; one
extends the lemma to a statement about arbitrary antimatroids (and even a
wider setting). Finally, I will strengthen the lemma to a Morse-theoretical
result, stating the collapsibility of a certain simplicial complex. Some open
questions will be posed.

\subsection*{Note to the reader}

The pictures on the title page illustrate the simplicial complex $\mathcal{A}$
from Proposition \ref{prop.elser-sc} on an example. The left picture is a
graph $\Gamma$ (with the vertex labelled $v$ playing the role of $v$), whereas
the right picture shows the corresponding simplicial complex $\mathcal{A}$ for
$G=E$ (that is, the simplicial complex whose faces are the subsets of $E$ that
are not pandemic).

Much of this text has been conceived and written during a stay at the
Mathematisches Forschungsinstitut Oberwolfach in 2020. This research was
supported through the programme \textquotedblleft Oberwolfach Leibniz
Fellows\textquotedblright\ by the Mathematisches Forschungsinstitut
Oberwolfach in 2020.

An older version of this text has been published as Oberwolfach Preprint OWP-2021-05.

\subsection*{Remark on alternative versions}

\begin{vershort}
This paper also has a detailed version \cite{verlong}, which elaborates on the proofs.
\end{vershort}

\begin{verlong}
You are reading the detailed version of this paper. For the standard version
(which is shorter by virtue of omitting some details in proofs), see
\cite{vershort}.
\end{verlong}

\section{\label{sec.elser}Elser's result}

Let us first introduce our setting, which is slightly more general (and
perhaps also simpler) than that used in \cite{Elser84}. (In Section
\ref{sec.abstract}, we will move to a more general setup.)

We fix an arbitrary graph $\Gamma$ with vertex set $V$ and edge set $E$. Here,
\textquotedblleft graph\textquotedblright\ means \textquotedblleft finite
undirected multigraph\textquotedblright\ -- i.e., it can have self-loops and
parallel edges, but it has finitely many vertices and edges, and its edges are undirected.

We fix a vertex $v\in V$.

If $F\subseteq E$, then an $F$\emph{-path} shall mean a path of $\Gamma$ such
that all edges of the path belong to $F$.

If $e\in E$ is any edge and $F\subseteq E$ is any subset, then we say that $F$
\emph{infects }$e$ if there exists an $F$-path from $v$ to some endpoint of
$e$. (The terminology is inspired by the idea of an infectious disease
starting in the vertex $v$ and being transmitted along edges.)\footnote{Note
that if an edge $e$ contains the vertex $v$, then any subset $F$ of $E$ (even
the empty one) infects $e$, since there is a trivial (edgeless) $F$-path from
$v$ to $v$.}

A subset $F\subseteq E$ is said to be \emph{pandemic} if it infects each edge
$e\in E$.

\begin{example}
\label{exa.example1} Let $\Gamma$ be the following graph:%
\[%
%TCIMACRO{\TeXButton{graph with 6 vertices and 8 edges}{\begin{tikzpicture}%
%[-,>=stealth',shorten >=1pt,auto,node distance=3cm, thick,main node/.style={circle,fill=blue!20,draw}%
%]
%\node[main node] (1) {$v$};
%\node[main node] [above of=1] (2) {$p$};
%\node[main node] [right of=1] (3) {$w$};
%\node[main node] [above of=3] (4) {$q$};
%\node[main node] [right of=3] (5) {$t$};
%\node[main node] [above of=5] (6) {$r$};
%\path[every node/.style={font=\sffamily\small}] (1) edge node {$1$}
%(2) (2) edge node {$2$} (4) (4) edge node {$3$}
%(6) (6) edge [bend left] node {$4$} (5) (5) edge node {$5$}
%(3) (3) edge node {$6$} (1) (5) edge [bend left] node {$7$}
%(6) (4) edge node {$8$} (3);
%\end{tikzpicture}}}%
%BeginExpansion
\begin{tikzpicture}%
[-,>=stealth',shorten >=1pt,auto,node distance=3cm, thick,main node/.style={circle,fill=blue!20,draw}%
]
\node[main node] (1) {$v$};
\node[main node] [above of=1] (2) {$p$};
\node[main node] [right of=1] (3) {$w$};
\node[main node] [above of=3] (4) {$q$};
\node[main node] [right of=3] (5) {$t$};
\node[main node] [above of=5] (6) {$r$};
\path[every node/.style={font=\sffamily\small}] (1) edge node {$1$}
(2) (2) edge node {$2$} (4) (4) edge node {$3$}
(6) (6) edge [bend left] node {$4$} (5) (5) edge node {$5$}
(3) (3) edge node {$6$} (1) (5) edge [bend left] node {$7$}
(6) (4) edge node {$8$} (3);
\end{tikzpicture}%
%EndExpansion
\]
(where the vertex $v$ is the vertex labelled $v$). Then, for example, the set
$\left\{  1,2\right\}  \subseteq E$ infects edges $1,2,3,6,8$ (but none of the
other edges). The set $\left\{  1,2,5\right\}  $ infects the same edges as
$\left\{  1,2\right\}  $ (indeed, the additional edge $5$ does not increase
its infectiousness, since it is not on any $\left\{  1,2,5\right\}  $-path
from $v$). The set $\left\{  1,2,3\right\}  $ infects every edge other than
$5$. The set $\left\{  1,2,3,4\right\}  $ infects each edge, and thus is pandemic.
\end{example}

Now, we can state our version of \cite[Lemma 1]{Elser84}:

\begin{theorem}
\label{thm.elser0}Assume that $E\neq\varnothing$. Then,%
\begin{equation}
\sum_{\substack{F\subseteq E\text{ is}\\\text{pandemic}}}\left(  -1\right)
^{\left\vert F\right\vert }=0. \label{eq.thm.elser0.eq}%
\end{equation}

\end{theorem}

\begin{example}
\label{exa.example2} Let $\Gamma$ be the following graph:%
\[%
%TCIMACRO{\TeXButton{C4 graph}{\begin{tikzpicture}%
%[-,>=stealth',shorten >=1pt,auto,node distance=3cm, thick,main node/.style={circle,fill=blue!20,draw}%
%]
%\node[main node] (1) {$v$};
%\node[main node] [above of=1] (2) {$p$};
%\node[main node] [right of=2] (3) {$q$};
%\node[main node] [right of=1] (4) {$w$};
%\path[every node/.style={font=\sffamily\small}] (1) edge node {$1$}
%(2) (2) edge node {$2$} (3) (3) edge node {$3$} (4) (4) edge node {$4$} (1);
%\end{tikzpicture}}}%
%BeginExpansion
\begin{tikzpicture}%
[-,>=stealth',shorten >=1pt,auto,node distance=3cm, thick,main node/.style={circle,fill=blue!20,draw}%
]
\node[main node] (1) {$v$};
\node[main node] [above of=1] (2) {$p$};
\node[main node] [right of=2] (3) {$q$};
\node[main node] [right of=1] (4) {$w$};
\path[every node/.style={font=\sffamily\small}] (1) edge node {$1$}
(2) (2) edge node {$2$} (3) (3) edge node {$3$} (4) (4) edge node {$4$} (1);
\end{tikzpicture}%
%EndExpansion
\]
(where the vertex $v$ is the vertex labelled $v$). Then, the pandemic subsets
of $E$ are the sets%
\[
\left\{  1,2\right\}  ,\ \ \left\{  1,4\right\}  ,\ \ \left\{  3,4\right\}
,\ \ \left\{  1,2,3\right\}  ,\ \ \left\{  1,3,4\right\}  ,\ \ \left\{
1,2,4\right\}  ,\ \ \left\{  2,3,4\right\}  ,\ \ \left\{  1,2,3,4\right\}  .
\]
The sizes of these subsets are $2,2,2,3,3,3,3,4$, respectively. Hence,
(\ref{eq.thm.elser0.eq}) says that%
\[
\left(  -1\right)  ^{2}+\left(  -1\right)  ^{2}+\left(  -1\right)
^{2}+\left(  -1\right)  ^{3}+\left(  -1\right)  ^{3}+\left(  -1\right)
^{3}+\left(  -1\right)  ^{3}+\left(  -1\right)  ^{4}=0.
\]

\end{example}

We note that the equality (\ref{eq.thm.elser0.eq}) can be restated as
\textquotedblleft there are equally many pandemic subsets $F\subseteq E$ of
even size and pandemic subsets $F\subseteq E$ of odd size\textquotedblright.
Thus, in particular, the number of all pandemic subsets $F$ of $E$ is even
(when $E\neq\varnothing$).

\begin{remark}
Theorem \ref{thm.elser0} is a bit more general than \cite[Lemma 1]{Elser84}.
To see why, we assume that the graph $\Gamma$ is connected and simple (i.e.,
has no self-loops and parallel edges). Then, a \emph{nucleus} is defined in
\cite{Elser84} as a subgraph $N$ of $\Gamma$ with the properties that

\begin{enumerate}
\item the subgraph $N$ is connected, and

\item each edge of $\Gamma$ has at least one endpoint in $N$.
\end{enumerate}

\noindent Given a subgraph $N$ of $\Gamma$, we let $\operatorname*{E}\left(
N\right)  $ denote the set of all edges of $N$. Now, \cite[Lemma 1]{Elser84}
claims that if $E\neq\varnothing$, then%
\[
\sum_{\substack{N\text{ is a nucleus}\\\text{containing }v}}\left(  -1\right)
^{\left\vert \operatorname*{E}\left(  N\right)  \right\vert }=0.
\]


But this is equivalent to (\ref{eq.thm.elser0.eq}), because there is a
bijection%
\begin{align*}
\left\{  \text{nuclei containing }v\right\}   &  \rightarrow\left\{
\text{pandemic subsets }F\subseteq E\right\}  ,\\
N  &  \mapsto\operatorname*{E}\left(  N\right)  .
\end{align*}
We leave it to the reader to check this in detail; what needs to be checked
are the following three statements:

\begin{itemize}
\item If $N$ is a nucleus containing $v$, then $\operatorname*{E}\left(
N\right)  $ is a pandemic subset of $E$.

\item Every nucleus $N$ containing $v$ is uniquely determined by the set
$\operatorname*{E}\left(  N\right)  $. (Indeed, since a nucleus has to be
connected, each of its vertices must be an endpoint of one of its edges,
unless its only vertex is $v$.)

\item If $F$ is a pandemic subset of $E$, then there is a nucleus $N$
containing $v$ such that $\operatorname*{E}\left(  N\right)  =F$. (Indeed, $N$
can be defined as the subgraph of $\Gamma$ whose vertices are the endpoints of
all edges in $F$ as well as the vertex $v$, and whose edges are the edges in
$F$. To see that this subgraph $N$ is connected, it suffices to argue that
each of its vertices has a path to $v$; but this follows from the definition
of \textquotedblleft pandemic\textquotedblright, since each vertex of $N$
other than $v$ belongs to at least one edge in $F$.)
\end{itemize}

Thus, Theorem \ref{thm.elser0} is equivalent to \cite[Lemma 1]{Elser84} in the
case when $\Gamma$ is connected and simple.
\end{remark}

\begin{remark}
It might appear more natural to talk about a subset $F\subseteq E$ infecting a
vertex rather than an edge. (Namely, we can say that $F$ infects a vertex $w$
if there is an $F$-path from $v$ to $w$.) However, the analogue of Theorem
\ref{thm.elser0} in which pandemicity is defined via infecting all vertices is
not true. The graph of Example \ref{exa.example2} provides a counterexample.
\end{remark}

\section{The proof}

\subsection{Shades}

Our proof of Theorem \ref{thm.elser0} will rest on a few notions. The first is
that of a \emph{shade}:

\begin{definition}
\label{def.Shade}Let $F$ be a subset of $E$. Then, we define a subset
$\operatorname*{Shade}F$ of $E$ by%
\begin{equation}
\operatorname*{Shade}F=\left\{  e\in E\ \mid\ F\text{ infects }e\right\}  .
\label{eq.def.Shade.eq}%
\end{equation}
We refer to $\operatorname*{Shade}F$ as the \emph{shade} of $F$.
\end{definition}

Thus, the shade of a subset $F\subseteq E$ is the set of all edges of $\Gamma$
that are infected by $F$. (In more standard graph-theoretical lingo, this
means that $\operatorname*{Shade}F$ is the set of edges that contain at least
one vertex of the connected component containing $v$ of the graph $\left(
V,F\right)  $.)

\begin{example}
In Example \ref{exa.example1}, we have $\operatorname*{Shade}\left\{
1,2\right\}  =\left\{  1,2,3,6,8\right\}  $ and $\operatorname*{Shade}\left\{
1\right\}  =\left\{  1,2,6\right\}  $ and $\operatorname*{Shade}\left\{
8\right\}  =\left\{  1,6\right\}  $.
\end{example}

The following property of shades is rather obvious:

\begin{lemma}
\label{lem.Shade-monoton}Let $A$ and $B$ be two subsets of $E$ such that
$A\subseteq B$. Then, $\operatorname*{Shade}A\subseteq\operatorname*{Shade}B$.
\end{lemma}

\begin{vershort}
\begin{proof}
[Proof of Lemma \ref{lem.Shade-monoton}.]We must show that each $q\in
\operatorname*{Shade}A$ satisfies $q\in\operatorname*{Shade}B$. In other
words, we must prove that if $A$ infects some edge $q\in E$, then $B$ also
infects this edge $q$. But this is clear, since any $A$-path is a $B$-path.
\end{proof}
\end{vershort}

\begin{verlong}
\begin{proof}
[Proof of Lemma \ref{lem.Shade-monoton}.]Let $q\in\operatorname*{Shade}A$.
Thus, $q\in\operatorname*{Shade}A=\left\{  e\in E\ \mid\ A\text{ infects
}e\right\}  $ (by the definition of $\operatorname*{Shade}A$). In other words,
$q$ is an $e\in E$ such that $A$ infects $e$. In other words, $q$ is an
element of $E$ with the property that $A$ infects $q$.

We know that $A$ infects $q$. In other words, there exists a $A$-path from $v$
to some endpoint of $q$ (by the definition of \textquotedblleft
infects\textquotedblright). Hence, there exists a $B$-path from $v$ to some
endpoint of $q$ (since any $A$-path is automatically a $B$%
-path\footnote{because $A\subseteq B$}). In other words, $B$ infects $q$ (by
the definition of \textquotedblleft infects\textquotedblright). Thus, $q$ is
an $e\in E$ such that $B$ infects $e$ (since $q$ is an element of $E$). In
other words, $q\in\left\{  e\in E\ \mid\ B\text{ infects }e\right\}  $. In
other words, $q\in\operatorname*{Shade}B$ (since the definition of
$\operatorname*{Shade}B$ yields $\operatorname*{Shade}B=\left\{  e\in
E\ \mid\ B\text{ infects }e\right\}  $).

Forget that we fixed $q$. We thus have proved that $q\in\operatorname*{Shade}%
B$ for each $q\in\operatorname*{Shade}A$. In other words,
$\operatorname*{Shade}A\subseteq\operatorname*{Shade}B$. This proves Lemma
\ref{lem.Shade-monoton}.
\end{proof}
\end{verlong}

The major property of shades that we will need is the following:

\begin{lemma}
\label{lem.Shade-tog}Let $F$ be a subset of $E$. Let $u\in E$ be such that
$u\notin\operatorname*{Shade}F$. Then,%
\begin{equation}
\operatorname*{Shade}\left(  F\cup\left\{  u\right\}  \right)
=\operatorname*{Shade}F \label{eq.lem.Shade-tog.union}%
\end{equation}
and%
\begin{equation}
\operatorname*{Shade}\left(  F\setminus\left\{  u\right\}  \right)
=\operatorname*{Shade}F. \label{eq.lem.Shade-tog.diff}%
\end{equation}

\end{lemma}

\begin{proof}
[Proof of Lemma \ref{lem.Shade-tog}.]We shall prove
(\ref{eq.lem.Shade-tog.union}) and (\ref{eq.lem.Shade-tog.diff}) separately:

\begin{vershort}
[\textit{Proof of (\ref{eq.lem.Shade-tog.union}):} Let $q\in
\operatorname*{Shade}\left(  F\cup\left\{  u\right\}  \right)  $. We shall
show that $q\in\operatorname*{Shade}F$.

We have assumed that $q\in\operatorname*{Shade}\left(  F\cup\left\{
u\right\}  \right)  $. In other words, $q$ is an edge in $E$ with the property
that $F\cup\left\{  u\right\}  $ infects $q$ (by the definition of
$\operatorname*{Shade}\left(  F\cup\left\{  u\right\}  \right)  $).

We shall now show that $F$ infects $q$. Indeed, assume the contrary. Thus, $F$
does not infect $q$. In other words, there exists no $F$-path from $v$ to any
endpoint of $q$ (by the definition of \textquotedblleft
infects\textquotedblright).

We know that $F\cup\left\{  u\right\}  $ infects $q$. In other words, there
exists an $\left(  F\cup\left\{  u\right\}  \right)  $-path from $v$ to some
endpoint of $q$ (by the definition of \textquotedblleft
infects\textquotedblright). Let $\pi$ be this path. If this $\left(
F\cup\left\{  u\right\}  \right)  $-path $\pi$ did not contain the edge $u$,
then it would be an $F$-path, which would contradict the fact that there
exists no $F$-path from $v$ to any endpoint of $q$. Hence, this $\left(
F\cup\left\{  u\right\}  \right)  $-path $\pi$ must contain the edge $u$. By
removing $u$, we can thus cut this path $\pi$ into two segments: The first
segment is a path from $v$ to some endpoint of $u$, while the second segment
is a path from the other endpoint of $u$ to some endpoint of $q$. Both
segments are $F$-paths. Thus, in particular, the first segment is an $F$-path
from $v$ to some endpoint of $u$. Hence, there exists an $F$-path from $v$ to
some endpoint of $u$. In other words, $F$ infects $u$ (by the definition of
\textquotedblleft infects\textquotedblright). Hence, $u\in
\operatorname*{Shade}F$ (because of (\ref{eq.def.Shade.eq})). This contradicts
$u\notin\operatorname*{Shade}F$.

This contradiction shows that our assumption was false. Hence, we have proved
that $F$ infects $q$. In other words $q\in\operatorname*{Shade}F$.

Forget that we fixed $q$. We thus have shown that $q\in\operatorname*{Shade}F$
for each $q\in\operatorname*{Shade}\left(  F\cup\left\{  u\right\}  \right)
$. In other words, $\operatorname*{Shade}\left(  F\cup\left\{  u\right\}
\right)  \subseteq\operatorname*{Shade}F$. On the other hand, $F\subseteq
F\cup\left\{  u\right\}  $; therefore, $\operatorname*{Shade}F\subseteq
\operatorname*{Shade}\left(  F\cup\left\{  u\right\}  \right)  $ (by Lemma
\ref{lem.Shade-monoton}). Combining this with $\operatorname*{Shade}\left(
F\cup\left\{  u\right\}  \right)  \subseteq\operatorname*{Shade}F$, we obtain
$\operatorname*{Shade}\left(  F\cup\left\{  u\right\}  \right)
=\operatorname*{Shade}F$.\ This proves (\ref{eq.lem.Shade-tog.union}).]
\end{vershort}

\begin{verlong}
[\textit{Proof of (\ref{eq.lem.Shade-tog.union}):} Note that $F\cup\left\{
u\right\}  $ is a subset of $E$ (since $F$ is a subset of $E$, and since $u\in
E$). Thus, $\operatorname*{Shade}\left(  F\cup\left\{  u\right\}  \right)  $
is well-defined.

Let $q\in\operatorname*{Shade}\left(  F\cup\left\{  u\right\}  \right)  $. We
shall show that $q\in\operatorname*{Shade}F$.

We have $q\in\operatorname*{Shade}\left(  F\cup\left\{  u\right\}  \right)
=\left\{  e\in E\ \mid\ F\cup\left\{  u\right\}  \text{ infects }e\right\}  $
(by the definition of $\operatorname*{Shade}\left(  F\cup\left\{  u\right\}
\right)  $). In other words, $q$ is an $e\in E$ such that $F\cup\left\{
u\right\}  $ infects $e$. In other words, $q$ is an element of $E$ with the
property that $F\cup\left\{  u\right\}  $ infects $q$.

We shall now show that $F$ infects $q$. Indeed, assume the contrary. Thus, $F$
does not infect $q$. In other words, there exists no $F$-path from $v$ to any
endpoint of $q$ (by the definition of \textquotedblleft
infects\textquotedblright).

We know that $F\cup\left\{  u\right\}  $ infects $q$. In other words, there
exists an $\left(  F\cup\left\{  u\right\}  \right)  $-path from $v$ to some
endpoint of $q$ (by the definition of \textquotedblleft
infects\textquotedblright). Let $\pi$ be this $\left(  F\cup\left\{
u\right\}  \right)  $-path. If this $\left(  F\cup\left\{  u\right\}  \right)
$-path $\pi$ did not contain the edge $u$, then it would be an $F$-path, which
would contradict the fact that there exists no $F$-path from $v$ to any
endpoint of $q$. Hence, this $\left(  F\cup\left\{  u\right\}  \right)  $-path
$\pi$ must contain the edge $u$. By removing $u$, we can thus cut this path
$\pi$ into two segments: The first segment is a path from $v$ to some endpoint
of $u$, while the second segment is a path from the other endpoint of $u$ to
some endpoint of $q$. Both segments are $F$-paths (since they arise by
removing $u$ from an $\left(  F\cup\left\{  u\right\}  \right)  $-path). Thus,
in particular, the first segment is an $F$-path from $v$ to some endpoint of
$u$. Hence, there exists an $F$-path from $v$ to some endpoint of $u$. In
other words, $F$ infects $u$ (by the definition of \textquotedblleft
infects\textquotedblright). Hence, $u$ is an $e\in E$ such that $F$ infects
$e$ (since $u\in E$). In other words, $u\in\left\{  e\in E\ \mid\ F\text{
infects }e\right\}  $. This can be rewritten as $u\in\operatorname*{Shade}F$
(because of (\ref{eq.def.Shade.eq})). This contradicts $u\notin%
\operatorname*{Shade}F$.

This contradiction shows that our assumption was false. Hence, we have proved
that $F$ infects $q$. Thus, $q$ is an $e\in E$ such that $F$ infects $e$
(since $q\in E$). In other words, $q\in\left\{  e\in E\ \mid\ F\text{ infects
}e\right\}  $. This can be rewritten as $q\in\operatorname*{Shade}F$ (because
of (\ref{eq.def.Shade.eq})).

Forget that we fixed $q$. We thus have shown that $q\in\operatorname*{Shade}F$
for each $q\in\operatorname*{Shade}\left(  F\cup\left\{  u\right\}  \right)
$. In other words, $\operatorname*{Shade}\left(  F\cup\left\{  u\right\}
\right)  \subseteq\operatorname*{Shade}F$. On the other hand, $F\subseteq
F\cup\left\{  u\right\}  $; therefore, $\operatorname*{Shade}F\subseteq
\operatorname*{Shade}\left(  F\cup\left\{  u\right\}  \right)  $ (by Lemma
\ref{lem.Shade-monoton}, applied to $A=F$ and $B=F\cup\left\{  u\right\}  $).
Combining this with $\operatorname*{Shade}\left(  F\cup\left\{  u\right\}
\right)  \subseteq\operatorname*{Shade}F$, we obtain $\operatorname*{Shade}%
\left(  F\cup\left\{  u\right\}  \right)  =\operatorname*{Shade}F$.\ This
proves (\ref{eq.lem.Shade-tog.union}).]
\end{verlong}

\begin{vershort}
[\textit{Proof of (\ref{eq.lem.Shade-tog.diff}):} We must prove that
$\operatorname*{Shade}\left(  F\setminus\left\{  u\right\}  \right)
=\operatorname*{Shade}F$. This is obvious if $F\setminus\left\{  u\right\}
=F$. Thus, for the rest of this proof, we WLOG assume that $F\setminus\left\{
u\right\}  \neq F$. Hence, $u\in F$ and thus $\left(  F\setminus\left\{
u\right\}  \right)  \cup\left\{  u\right\}  =F$.

We have $F\setminus\left\{  u\right\}  \subseteq F$ and thus
$\operatorname*{Shade}\left(  F\setminus\left\{  u\right\}  \right)
\subseteq\operatorname*{Shade}F$ (by Lemma \ref{lem.Shade-monoton}). Hence,
from $u\notin\operatorname*{Shade}F$, we obtain $u\notin\operatorname*{Shade}%
\left(  F\setminus\left\{  u\right\}  \right)  $. Therefore,
(\ref{eq.lem.Shade-tog.union}) (applied to $F\setminus\left\{  u\right\}  $
instead of $F$) yields $\operatorname*{Shade}\left(  \left(  F\setminus
\left\{  u\right\}  \right)  \cup\left\{  u\right\}  \right)
=\operatorname*{Shade}\left(  F\setminus\left\{  u\right\}  \right)  $. Thus,
$\operatorname*{Shade}\left(  F\setminus\left\{  u\right\}  \right)
=\operatorname*{Shade}\underbrace{\left(  \left(  F\setminus\left\{
u\right\}  \right)  \cup\left\{  u\right\}  \right)  }_{=F}%
=\operatorname*{Shade}F$. This proves (\ref{eq.lem.Shade-tog.diff}).]
\end{vershort}

\begin{verlong}
[\textit{Proof of (\ref{eq.lem.Shade-tog.diff}):} Note that $F\setminus
\left\{  u\right\}  $ is a subset of $E$ (since $F\setminus\left\{  u\right\}
\subseteq F\subseteq E$). Thus, $\operatorname*{Shade}\left(  F\setminus
\left\{  u\right\}  \right)  $ is well-defined.

We must prove that $\operatorname*{Shade}\left(  F\setminus\left\{  u\right\}
\right)  =\operatorname*{Shade}F$. This is obvious if $F\setminus\left\{
u\right\}  =F$. Thus, for the rest of this proof, we WLOG assume that
$F\setminus\left\{  u\right\}  \neq F$. Hence, $u\in F$ (since otherwise, we
would have $u\notin F$ and thus $F\setminus\left\{  u\right\}  =F$, which
would contradict $F\setminus\left\{  u\right\}  \neq F$). Thus, $\left(
F\setminus\left\{  u\right\}  \right)  \cup\left\{  u\right\}  =F$.

We have $F\setminus\left\{  u\right\}  \subseteq F$ and thus
$\operatorname*{Shade}\left(  F\setminus\left\{  u\right\}  \right)
\subseteq\operatorname*{Shade}F$ (by Lemma \ref{lem.Shade-monoton}, applied to
$A=F\setminus\left\{  u\right\}  $ and $B=F$). Hence, from $u\notin%
\operatorname*{Shade}F$, we obtain $u\notin\operatorname*{Shade}\left(
F\setminus\left\{  u\right\}  \right)  $ (because otherwise, we would have
$u\in\operatorname*{Shade}\left(  F\setminus\left\{  u\right\}  \right)
\subseteq\operatorname*{Shade}F$, which would contradict $u\notin%
\operatorname*{Shade}F$). Therefore, (\ref{eq.lem.Shade-tog.union}) (applied
to $F\setminus\left\{  u\right\}  $ instead of $F$) yields
$\operatorname*{Shade}\left(  \left(  F\setminus\left\{  u\right\}  \right)
\cup\left\{  u\right\}  \right)  =\operatorname*{Shade}\left(  F\setminus
\left\{  u\right\}  \right)  $. Thus, $\operatorname*{Shade}\left(
F\setminus\left\{  u\right\}  \right)  =\operatorname*{Shade}%
\underbrace{\left(  \left(  F\setminus\left\{  u\right\}  \right)
\cup\left\{  u\right\}  \right)  }_{=F}=\operatorname*{Shade}F$. This proves
(\ref{eq.lem.Shade-tog.diff}).]
\end{verlong}

We have now proved both (\ref{eq.lem.Shade-tog.union}) and
(\ref{eq.lem.Shade-tog.diff}). Thus, Lemma \ref{lem.Shade-tog} is proved.
\end{proof}

\subsection{A slightly more general claim}

Lemma \ref{lem.Shade-tog} might not look very powerful, but it contains all we
need to prove Theorem \ref{thm.elser0}. Better yet, we shall prove the
following slightly more general version of Theorem \ref{thm.elser0}:

\begin{theorem}
\label{thm.elser-Shade1}Let $G$ be any subset of $E$. Assume that
$E\neq\varnothing$. Then,%
\[
\sum_{\substack{F\subseteq E;\\G\subseteq\operatorname*{Shade}F}}\left(
-1\right)  ^{\left\vert F\right\vert }=0.
\]

\end{theorem}

We will soon prove Theorem \ref{thm.elser-Shade1} and explain how Theorem
\ref{thm.elser0} follows from it. First, however, let us give an equivalent
(but slightly easier to prove) version of Theorem \ref{thm.elser-Shade1}:

\begin{theorem}
\label{thm.elser-Shade0}Let $G$ be any subset of $E$. Then,%
\[
\sum_{\substack{F\subseteq E;\\G\not \subseteq \operatorname*{Shade}F}}\left(
-1\right)  ^{\left\vert F\right\vert }=0.
\]

\end{theorem}

\begin{vershort}
\begin{proof}
[Proof of Theorem \ref{thm.elser-Shade0}.]Let
\begin{equation}
\mathcal{A}=\left\{  P\subseteq E\ \mid\ G\not \subseteq \operatorname*{Shade}%
P\right\}  . \label{pf.thm.elser-Shade0.short.A=}%
\end{equation}
Thus, $\mathcal{A}$ is a subset of the power set of $E$, and each
$F\in\mathcal{A}$ satisfies $G\not \subseteq \operatorname*{Shade}F$.

We equip the finite set $E$ with a total order (chosen arbitrarily, but fixed
henceforth). If $F\in\mathcal{A}$, then there exists a unique
\textbf{smallest} edge $e\in G\setminus\operatorname*{Shade}F$ (since
$F\in\mathcal{A}$ entails $G\not \subseteq \operatorname*{Shade}F$ and thus
$G\setminus\operatorname*{Shade}F\neq\varnothing$). This unique smallest edge
$e$ will be denoted by $\varepsilon\left(  F\right)  $.

We note that the edge $\varepsilon\left(  F\right)  $ (for a set
$F\in\mathcal{A}$) depends only on $\operatorname*{Shade}F$, but not on $F$
itself (because it was defined as the smallest edge $e\in G\setminus
\operatorname*{Shade}F$). Thus, if two sets $F_{1}\in\mathcal{A}$ and
$F_{2}\in\mathcal{A}$ satisfy $\operatorname*{Shade}\left(  F_{1}\right)
=\operatorname*{Shade}\left(  F_{2}\right)  $, then
\begin{equation}
\varepsilon\left(  F_{1}\right)  =\varepsilon\left(  F_{2}\right)  .
\label{pf.thm.elser-Shade0.short.eps=eps}%
\end{equation}


We also notice the following simple fact: If $F$ and $F^{\prime}$ are two
subsets of $E$ such that $F\in\mathcal{A}$ and $\operatorname*{Shade}\left(
F^{\prime}\right)  =\operatorname*{Shade}F$, then%
\begin{equation}
F^{\prime}\in\mathcal{A}. \label{pf.thm.elser-Shade0.short.Fpri}%
\end{equation}
(Indeed, $F\in\mathcal{A}$ means that $G\not \subseteq \operatorname*{Shade}%
F$; but because of $\operatorname*{Shade}\left(  F^{\prime}\right)
=\operatorname*{Shade}F$, this entails $G\not \subseteq \operatorname*{Shade}%
\left(  F^{\prime}\right)  $ as well, and therefore $F^{\prime}\in\mathcal{A}$.)

We now define two subsets $\mathcal{A}_{+}$ and $\mathcal{A}_{-}$ of
$\mathcal{A}$ by%
\[
\mathcal{A}_{+}=\left\{  P\in\mathcal{A}\ \mid\ \varepsilon\left(  P\right)
\in P\right\}  \ \ \ \ \ \ \ \ \ \ \text{and}\ \ \ \ \ \ \ \ \ \ \mathcal{A}%
_{-}=\left\{  P\in\mathcal{A}\ \mid\ \varepsilon\left(  P\right)  \notin
P\right\}  .
\]


Next, we claim the following:

\begin{statement}
\textit{Claim 1:} Let $F\in\mathcal{A}_{+}$. Set $F^{\prime}=F\setminus
\left\{  \varepsilon\left(  F\right)  \right\}  $. Then, $F^{\prime}%
\in\mathcal{A}_{-}$ and $F^{\prime}\cup\left\{  \varepsilon\left(  F^{\prime
}\right)  \right\}  =F$ and $\left(  -1\right)  ^{\left\vert F^{\prime
}\right\vert }=-\left(  -1\right)  ^{\left\vert F\right\vert }$.
\end{statement}

[\textit{Proof of Claim 1:} We have $F\in\mathcal{A}_{+}=\left\{
P\in\mathcal{A}\ \mid\ \varepsilon\left(  P\right)  \in P\right\}  $; in other
words, $F\in\mathcal{A}$ and $\varepsilon\left(  F\right)  \in F$. From
$F\in\mathcal{A}=\left\{  P\subseteq E\ \mid\ G\not \subseteq
\operatorname*{Shade}P\right\}  $, we obtain $F\subseteq E$ and
$G\not \subseteq \operatorname*{Shade}F$.

The definition of $F^{\prime}$ yields $F^{\prime}\subseteq F\subseteq E$.

Recall that $\varepsilon\left(  F\right)  $ is the smallest edge $e\in
G\setminus\operatorname*{Shade}F$ (by the definition of $\varepsilon\left(
F\right)  $). Hence, $\varepsilon\left(  F\right)  \in G\setminus
\operatorname*{Shade}F$. In other words, $\varepsilon\left(  F\right)  \in G$
and $\varepsilon\left(  F\right)  \notin\operatorname*{Shade}F$. Thus,
$\varepsilon\left(  F\right)  \in G\subseteq E$ and $\varepsilon\left(
F\right)  \notin\operatorname*{Shade}F$. Therefore,
(\ref{eq.lem.Shade-tog.diff}) (applied to $u=\varepsilon\left(  F\right)  $)
yields $\operatorname*{Shade}\left(  F\setminus\left\{  \varepsilon\left(
F\right)  \right\}  \right)  =\operatorname*{Shade}F$. This can be rewritten
as $\operatorname*{Shade}\left(  F^{\prime}\right)  =\operatorname*{Shade}F$
(since $F^{\prime}=F\setminus\left\{  \varepsilon\left(  F\right)  \right\}
$). Hence, (\ref{pf.thm.elser-Shade0.short.Fpri}) yields $F^{\prime}%
\in\mathcal{A}$. In light of the preceding two sentences,
(\ref{pf.thm.elser-Shade0.short.eps=eps}) (applied to $F^{\prime}$ and $F$
instead of $F_{1}$ and $F_{2}$) yields $\varepsilon\left(  F^{\prime}\right)
=\varepsilon\left(  F\right)  $. However, $\varepsilon\left(  F\right)  \notin
F^{\prime}$ (since $F^{\prime}=F\setminus\left\{  \varepsilon\left(  F\right)
\right\}  $). In other words, $\varepsilon\left(  F^{\prime}\right)  \notin
F^{\prime}$ (since $\varepsilon\left(  F^{\prime}\right)  =\varepsilon\left(
F\right)  $). Hence, $F^{\prime}\in\left\{  P\in\mathcal{A}\ \mid
\ \varepsilon\left(  P\right)  \notin P\right\}  $ (since $F^{\prime}%
\in\mathcal{A}$). In other words, $F^{\prime}\in\mathcal{A}_{-}$ (since
$\mathcal{A}_{-}=\left\{  P\in\mathcal{A}\ \mid\ \varepsilon\left(  P\right)
\notin P\right\}  $).

Moreover, from $\varepsilon\left(  F^{\prime}\right)  =\varepsilon\left(
F\right)  $, we obtain $F^{\prime}\cup\left\{  \varepsilon\left(  F^{\prime
}\right)  \right\}  =F^{\prime}\cup\left\{  \varepsilon\left(  F\right)
\right\}  =F$ (since $F^{\prime}=F\setminus\left\{  \varepsilon\left(
F\right)  \right\}  $ and $\varepsilon\left(  F\right)  \in F$).

Finally, the set $F^{\prime}=F\setminus\left\{  \varepsilon\left(  F\right)
\right\}  $ has exactly one less element than the set $F$ (since
$\varepsilon\left(  F\right)  \in F$). That is, $\left\vert F^{\prime
}\right\vert =\left\vert F\right\vert -1$. Hence, $\left(  -1\right)
^{\left\vert F^{\prime}\right\vert }=-\left(  -1\right)  ^{\left\vert
F\right\vert }$. This completes the proof of Claim 1.]

\begin{statement}
\textit{Claim 2:} Let $F\in\mathcal{A}_{-}$. Set $F^{\prime}=F\cup\left\{
\varepsilon\left(  F\right)  \right\}  $. Then, $F^{\prime}\in\mathcal{A}_{+}$
and $F^{\prime}\setminus\left\{  \varepsilon\left(  F^{\prime}\right)
\right\}  =F$.
\end{statement}

[\textit{Proof of Claim 2:} We have $F\in\mathcal{A}_{-}=\left\{
P\in\mathcal{A}\ \mid\ \varepsilon\left(  P\right)  \notin P\right\}  $; in
other words, $F\in\mathcal{A}$ and $\varepsilon\left(  F\right)  \notin F$.
From $F\in\mathcal{A}=\left\{  P\subseteq E\ \mid\ G\not \subseteq
\operatorname*{Shade}P\right\}  $, we obtain $F\subseteq E$ and
$G\not \subseteq \operatorname*{Shade}F$.

As in the proof of Claim 1, we can see that $\varepsilon\left(  F\right)  \in
G$ and $\varepsilon\left(  F\right)  \notin\operatorname*{Shade}F$. Thus,
$\varepsilon\left(  F\right)  \in G\subseteq E$ and $\varepsilon\left(
F\right)  \notin\operatorname*{Shade}F$. Now, $F^{\prime}=F\cup\left\{
\varepsilon\left(  F\right)  \right\}  \subseteq E$ (since $F\subseteq E$ and
$\varepsilon\left(  F\right)  \in E$). Furthermore,
(\ref{eq.lem.Shade-tog.union}) (applied to $u=\varepsilon\left(  F\right)  $)
yields $\operatorname*{Shade}\left(  F\cup\left\{  \varepsilon\left(
F\right)  \right\}  \right)  =\operatorname*{Shade}F$ (since $\varepsilon
\left(  F\right)  \notin\operatorname*{Shade}F$). This can be rewritten as
$\operatorname*{Shade}\left(  F^{\prime}\right)  =\operatorname*{Shade}F$
(since $F^{\prime}=F\cup\left\{  \varepsilon\left(  F\right)  \right\}  $).
Hence, (\ref{pf.thm.elser-Shade0.short.Fpri}) yields $F^{\prime}\in
\mathcal{A}$. In light of the preceding two sentences,
(\ref{pf.thm.elser-Shade0.short.eps=eps}) (applied to $F^{\prime}$ and $F$
instead of $F_{1}$ and $F_{2}$) yields $\varepsilon\left(  F^{\prime}\right)
=\varepsilon\left(  F\right)  \in\left\{  \varepsilon\left(  F\right)
\right\}  \subseteq F\cup\left\{  \varepsilon\left(  F\right)  \right\}
=F^{\prime}$. Thus, $F^{\prime}\in\left\{  P\in\mathcal{A}\ \mid
\ \varepsilon\left(  P\right)  \in P\right\}  $ (since $F^{\prime}%
\in\mathcal{A}$). In other words, $F^{\prime}\in\mathcal{A}_{+}$ (since
$\mathcal{A}_{+}=\left\{  P\in\mathcal{A}\ \mid\ \varepsilon\left(  P\right)
\in P\right\}  $).

Moreover, from $\varepsilon\left(  F^{\prime}\right)  =\varepsilon\left(
F\right)  $, we obtain $F^{\prime}\setminus\left\{  \varepsilon\left(
F^{\prime}\right)  \right\}  =F^{\prime}\setminus\left\{  \varepsilon\left(
F\right)  \right\}  =F$ (since $F^{\prime}=F\cup\left\{  \varepsilon\left(
F\right)  \right\}  $ and $\varepsilon\left(  F\right)  \notin F$). This
completes the proof of Claim 2.]

Each $F\in\mathcal{A}_{+}$ satisfies $F\setminus\left\{  \varepsilon\left(
F\right)  \right\}  \in\mathcal{A}_{-}$ (by Claim 1, applied to $F^{\prime
}=F\setminus\left\{  \varepsilon\left(  F\right)  \right\}  $). Thus, we can
define a map%
\begin{align*}
\Phi:\mathcal{A}_{+}  &  \rightarrow\mathcal{A}_{-},\\
F  &  \mapsto F\setminus\left\{  \varepsilon\left(  F\right)  \right\}  .
\end{align*}


Each $F\in\mathcal{A}_{-}$ satisfies $F\cup\left\{  \varepsilon\left(
F\right)  \right\}  \in\mathcal{A}_{+}$ (by Claim 2, applied to $F^{\prime
}=F\cup\left\{  \varepsilon\left(  F\right)  \right\}  $). Thus, we can define
a map%
\begin{align*}
\Psi:\mathcal{A}_{-}  &  \rightarrow\mathcal{A}_{+},\\
F  &  \mapsto F\cup\left\{  \varepsilon\left(  F\right)  \right\}  .
\end{align*}


We have $\Phi\circ\Psi=\operatorname*{id}$ (this follows from the
\textquotedblleft$F^{\prime}\setminus\left\{  \varepsilon\left(  F^{\prime
}\right)  \right\}  =F$\textquotedblright\ part of Claim 2) and $\Psi\circ
\Phi=\operatorname*{id}$ (this follows from the \textquotedblleft$F^{\prime
}\cup\left\{  \varepsilon\left(  F^{\prime}\right)  \right\}  =F$%
\textquotedblright\ part of Claim 1). Thus, the maps $\Phi$ and $\Psi$ are
mutually inverse. Hence, the map $\Phi$ is invertible, thus a bijection.

Moreover, each $F\in\mathcal{A}_{+}$ satisfies%
\begin{equation}
\left(  -1\right)  ^{\left\vert \Phi\left(  F\right)  \right\vert }=-\left(
-1\right)  ^{\left\vert F\right\vert }.
\label{pf.thm.elser-Shade0.short.signrev}%
\end{equation}
(Indeed, this is just the \textquotedblleft$\left(  -1\right)  ^{\left\vert
F^{\prime}\right\vert }=-\left(  -1\right)  ^{\left\vert F\right\vert }%
$\textquotedblright\ part of Claim 1.)

Now, the summation sign \textquotedblleft$\sum_{\substack{F\subseteq
E;\\G\not \subseteq \operatorname*{Shade}F}}$\textquotedblright\ is equivalent
to the summation sign \textquotedblleft$\sum_{F\in\mathcal{A}}$%
\textquotedblright\ (since the set of all subsets $F$ of $E$ satisfying
$G\not \subseteq \operatorname*{Shade}F$ is precisely $\mathcal{A}$). Thus,%
\begin{align*}
&  \sum_{\substack{F\subseteq E;\\G\not \subseteq \operatorname*{Shade}%
F}}\left(  -1\right)  ^{\left\vert F\right\vert }\\
&  =\sum_{F\in\mathcal{A}}\left(  -1\right)  ^{\left\vert F\right\vert
}=\underbrace{\sum_{\substack{F\in\mathcal{A};\\\varepsilon\left(  F\right)
\in F}}}_{\substack{=\sum_{F\in\mathcal{A}_{+}}\\\text{(since }\left\{
P\in\mathcal{A}\ \mid\ \varepsilon\left(  P\right)  \in P\right\}
=\mathcal{A}_{+}\text{)}}}\left(  -1\right)  ^{\left\vert F\right\vert
}+\underbrace{\sum_{\substack{F\in\mathcal{A};\\\varepsilon\left(  F\right)
\notin F}}}_{\substack{=\sum_{F\in\mathcal{A}_{-}}\\\text{(since }\left\{
P\in\mathcal{A}\ \mid\ \varepsilon\left(  P\right)  \notin P\right\}
=\mathcal{A}_{-}\text{)}}}\left(  -1\right)  ^{\left\vert F\right\vert }\\
&  =\sum_{F\in\mathcal{A}_{+}}\left(  -1\right)  ^{\left\vert F\right\vert
}+\sum_{F\in\mathcal{A}_{-}}\left(  -1\right)  ^{\left\vert F\right\vert
}=\sum_{F\in\mathcal{A}_{+}}\left(  -1\right)  ^{\left\vert F\right\vert
}+\sum_{F\in\mathcal{A}_{+}}\underbrace{\left(  -1\right)  ^{\left\vert
\Phi\left(  F\right)  \right\vert }}_{\substack{=-\left(  -1\right)
^{\left\vert F\right\vert }\\\text{(by
(\ref{pf.thm.elser-Shade0.short.signrev}))}}}\\
&  \ \ \ \ \ \ \ \ \ \ \left(
\begin{array}
[c]{c}%
\text{here, we have substituted }\Phi\left(  F\right)  \text{ for }F\text{ in
the second sum,}\\
\text{since the map }\Phi:\mathcal{A}_{+}\rightarrow\mathcal{A}_{-}\text{ is a
bijection}%
\end{array}
\right) \\
&  =\sum_{F\in\mathcal{A}_{+}}\left(  -1\right)  ^{\left\vert F\right\vert
}+\sum_{F\in\mathcal{A}_{+}}\left(  -\left(  -1\right)  ^{\left\vert
F\right\vert }\right)  =\sum_{F\in\mathcal{A}_{+}}\left(  -1\right)
^{\left\vert F\right\vert }-\sum_{F\in\mathcal{A}_{+}}\left(  -1\right)
^{\left\vert F\right\vert }=0.
\end{align*}
This proves Theorem \ref{thm.elser-Shade0}.
\end{proof}
\end{vershort}

\begin{verlong}
\begin{proof}
[Proof of Theorem \ref{thm.elser-Shade0}.]Let
\begin{equation}
\mathcal{A}=\left\{  P\subseteq E\ \mid\ G\not \subseteq \operatorname*{Shade}%
P\right\}  . \label{pf.thm.elser-Shade0.A=}%
\end{equation}
Thus, $\mathcal{A}$ is a subset of the power set of $E$.

We equip the finite set $E$ with a total order (chosen arbitrarily, but fixed
henceforth). If $F\in\mathcal{A}$, then there exists a unique
\textbf{smallest} edge $e\in G\setminus\operatorname*{Shade}F$%
\ \ \ \ \footnote{\textit{Proof.} Let $F\in\mathcal{A}$. Thus, $F\in
\mathcal{A}=\left\{  P\subseteq E\ \mid\ G\not \subseteq \operatorname*{Shade}%
P\right\}  $. In other words, $F$ is a $P\subseteq E$ satisfying
$G\not \subseteq \operatorname*{Shade}P$. In other words, $F$ is a subset of
$E$ and satisfies $G\not \subseteq \operatorname*{Shade}F$. From
$G\not \subseteq \operatorname*{Shade}F$, we see that $G\setminus
\operatorname*{Shade}F\neq\varnothing$. In other words, there exists at least
one edge $e\in G\setminus\operatorname*{Shade}F$. Hence, there exists a unique
\textbf{smallest} edge $e\in G\setminus\operatorname*{Shade}F$ (because the
set $G\setminus\operatorname*{Shade}F$ is finite (being a subset of $E$) and
totally ordered (being a subset of $E$)). Qed.}. This unique smallest edge $e$
will be denoted by $\varepsilon\left(  F\right)  $.

We note that the edge $\varepsilon\left(  F\right)  $ (for a set
$F\in\mathcal{A}$) depends only on $\operatorname*{Shade}F$, but not on $F$
itself (because it was defined as the smallest edge $e\in G\setminus
\operatorname*{Shade}F$). Thus, if two sets $F_{1}\in\mathcal{A}$ and
$F_{2}\in\mathcal{A}$ satisfy $\operatorname*{Shade}\left(  F_{1}\right)
=\operatorname*{Shade}\left(  F_{2}\right)  $, then
\begin{equation}
\varepsilon\left(  F_{1}\right)  =\varepsilon\left(  F_{2}\right)  .
\label{pf.thm.elser-Shade0.eps=eps}%
\end{equation}


We also notice the following simple fact: If $F$ and $F^{\prime}$ are two
subsets of $E$ such that $F\in\mathcal{A}$ and $\operatorname*{Shade}\left(
F^{\prime}\right)  =\operatorname*{Shade}F$, then%
\begin{equation}
F^{\prime}\in\mathcal{A} \label{pf.thm.elser-Shade0.Fpri}%
\end{equation}
\footnote{\textit{Proof of (\ref{pf.thm.elser-Shade0.Fpri}):} Let $F$ and
$F^{\prime}$ be two subsets of $E$ such that $F\in\mathcal{A}$ and
$\operatorname*{Shade}\left(  F^{\prime}\right)  =\operatorname*{Shade}F$. We
have $F\in\mathcal{A}=\left\{  P\subseteq E\ \mid\ G\not \subseteq
\operatorname*{Shade}P\right\}  $. In other words, $F$ is a $P\subseteq E$
satisfying $G\not \subseteq \operatorname*{Shade}P$. In other words, $F$ is a
subset of $E$ and satisfies $G\not \subseteq \operatorname*{Shade}F$. Thus,
$G\not \subseteq \operatorname*{Shade}F=\operatorname*{Shade}\left(
F^{\prime}\right)  $ (since $\operatorname*{Shade}\left(  F^{\prime}\right)
=\operatorname*{Shade}F$). Now, $F^{\prime}$ is a subset of $E$ and satisfies
$G\not \subseteq \operatorname*{Shade}\left(  F^{\prime}\right)  $. In other
words, $F^{\prime}$ is a $P\subseteq E$ satisfying $G\not \subseteq
\operatorname*{Shade}P$. In other words, $F^{\prime}\in\left\{  P\subseteq
E\ \mid\ G\not \subseteq \operatorname*{Shade}P\right\}  $. This can be
rewritten as $F^{\prime}\in\mathcal{A}$ (because of
(\ref{pf.thm.elser-Shade0.A=})). This proves (\ref{pf.thm.elser-Shade0.Fpri}%
).}

We now define two subsets $\mathcal{A}_{+}$ and $\mathcal{A}_{-}$ of
$\mathcal{A}$ by%
\[
\mathcal{A}_{+}=\left\{  P\in\mathcal{A}\ \mid\ \varepsilon\left(  P\right)
\in P\right\}  \ \ \ \ \ \ \ \ \ \ \text{and}\ \ \ \ \ \ \ \ \ \ \mathcal{A}%
_{-}=\left\{  P\in\mathcal{A}\ \mid\ \varepsilon\left(  P\right)  \notin
P\right\}  .
\]


Each $F\in\mathcal{A}_{+}$ satisfies $F\setminus\left\{  \varepsilon\left(
F\right)  \right\}  \in\mathcal{A}_{-}$\ \ \ \ \footnote{\textit{Proof.} Let
$F\in\mathcal{A}_{+}$. We must show that $F\setminus\left\{  \varepsilon
\left(  F\right)  \right\}  \in\mathcal{A}_{-}$.
\par
Indeed, let us set $F^{\prime}=F\setminus\left\{  \varepsilon\left(  F\right)
\right\}  $. Then, $F\in\mathcal{A}_{+}=\left\{  P\in\mathcal{A}%
\ \mid\ \varepsilon\left(  P\right)  \in P\right\}  \subseteq\mathcal{A}%
=\left\{  P\subseteq E\ \mid\ G\not \subseteq \operatorname*{Shade}P\right\}
$. In other words, $F$ is a $P\subseteq E$ satisfying $G\not \subseteq
\operatorname*{Shade}P$. Hence, $F$ is a subset of $E$. Thus, $F^{\prime}$ is
a subset of $E$ (since $F^{\prime}=F\setminus\left\{  \varepsilon\left(
F\right)  \right\}  \subseteq F\subseteq E$).
\par
Recall that $\varepsilon\left(  F\right)  $ is the smallest edge $e\in
G\setminus\operatorname*{Shade}F$ (by the definition of $\varepsilon\left(
F\right)  $). Hence, $\varepsilon\left(  F\right)  \in G\setminus
\operatorname*{Shade}F$. In other words, $\varepsilon\left(  F\right)  \in G$
and $\varepsilon\left(  F\right)  \notin\operatorname*{Shade}F$. Thus,
$\varepsilon\left(  F\right)  \in G\subseteq E$ and $\varepsilon\left(
F\right)  \notin\operatorname*{Shade}F$. Therefore,
(\ref{eq.lem.Shade-tog.diff}) (applied to $u=\varepsilon\left(  F\right)  $)
yields $\operatorname*{Shade}\left(  F\setminus\left\{  \varepsilon\left(
F\right)  \right\}  \right)  =\operatorname*{Shade}F$. This can be rewritten
as $\operatorname*{Shade}\left(  F^{\prime}\right)  =\operatorname*{Shade}F$
(since $F^{\prime}=F\setminus\left\{  \varepsilon\left(  F\right)  \right\}
$). Hence, (\ref{pf.thm.elser-Shade0.Fpri}) yields $F^{\prime}\in\mathcal{A}$.
Therefore, (\ref{pf.thm.elser-Shade0.eps=eps}) (applied to $F^{\prime}$ and
$F$ instead of $F_{1}$ and $F_{2}$) yields $\varepsilon\left(  F^{\prime
}\right)  =\varepsilon\left(  F\right)  $ (since $F\in\mathcal{A}$ and
$F^{\prime}\in\mathcal{A}$ and $\operatorname*{Shade}\left(  F^{\prime
}\right)  =\operatorname*{Shade}F$). However, $\varepsilon\left(  F\right)
\notin F^{\prime}$ (since $F^{\prime}=F\setminus\left\{  \varepsilon\left(
F\right)  \right\}  $). In other words, $\varepsilon\left(  F^{\prime}\right)
\notin F^{\prime}$ (since $\varepsilon\left(  F^{\prime}\right)
=\varepsilon\left(  F\right)  $). Hence, $F^{\prime}$ is a $P\in\mathcal{A}$
satisfying $\varepsilon\left(  P\right)  \notin P$ (since $F^{\prime}%
\in\mathcal{A}$ and $\varepsilon\left(  F^{\prime}\right)  \notin F^{\prime}%
$). In other words, $F^{\prime}\in\left\{  P\in\mathcal{A}\ \mid
\ \varepsilon\left(  P\right)  \notin P\right\}  $. In other words,
$F^{\prime}\in\mathcal{A}_{-}$ (since $\mathcal{A}_{-}=\left\{  P\in
\mathcal{A}\ \mid\ \varepsilon\left(  P\right)  \notin P\right\}  $). In other
words, $F\setminus\left\{  \varepsilon\left(  F\right)  \right\}
\in\mathcal{A}_{-}$ (since $F^{\prime}=F\setminus\left\{  \varepsilon\left(
F\right)  \right\}  $). Qed.}. Thus, we can define a map%
\begin{align*}
\Phi:\mathcal{A}_{+}  &  \rightarrow\mathcal{A}_{-},\\
F  &  \mapsto F\setminus\left\{  \varepsilon\left(  F\right)  \right\}  .
\end{align*}


Each $F\in\mathcal{A}_{-}$ satisfies $F\cup\left\{  \varepsilon\left(
F\right)  \right\}  \in\mathcal{A}_{+}$\ \ \ \ \footnote{\textit{Proof.} Let
$F\in\mathcal{A}_{-}$. We must show that $F\cup\left\{  \varepsilon\left(
F\right)  \right\}  \in\mathcal{A}_{+}$.
\par
Indeed, let us set $F^{\prime}=F\cup\left\{  \varepsilon\left(  F\right)
\right\}  $. Recall that $\varepsilon\left(  F\right)  $ is the smallest edge
$e\in G\setminus\operatorname*{Shade}F$ (by the definition of $\varepsilon
\left(  F\right)  $). Hence, $\varepsilon\left(  F\right)  \in G\setminus
\operatorname*{Shade}F$. In other words, $\varepsilon\left(  F\right)  \in G$
and $\varepsilon\left(  F\right)  \notin\operatorname*{Shade}F$. Thus,
$\varepsilon\left(  F\right)  \in G\subseteq E$ and $\varepsilon\left(
F\right)  \notin\operatorname*{Shade}F$.
\par
We have $F\in\mathcal{A}_{-}=\left\{  P\in\mathcal{A}\ \mid\ \varepsilon
\left(  P\right)  \notin P\right\}  \subseteq\mathcal{A}=\left\{  P\subseteq
E\ \mid\ G\not \subseteq \operatorname*{Shade}P\right\}  $. In other words,
$F$ is a $P\subseteq E$ satisfying $G\not \subseteq \operatorname*{Shade}P$.
Thus, $F$ is a subset of $E$. Therefore, $F\cup\left\{  \varepsilon\left(
F\right)  \right\}  $ is a subset of $E$ as well (since $\varepsilon\left(
F\right)  \in E$). In other words, $F^{\prime}$ is a subset of $E$ (since
$F^{\prime}=F\cup\left\{  \varepsilon\left(  F\right)  \right\}  $).
\par
Moreover, (\ref{eq.lem.Shade-tog.union}) (applied to $u=\varepsilon\left(
F\right)  $) yields $\operatorname*{Shade}\left(  F\cup\left\{  \varepsilon
\left(  F\right)  \right\}  \right)  =\operatorname*{Shade}F$ (since
$\varepsilon\left(  F\right)  \notin\operatorname*{Shade}F$). This can be
rewritten as $\operatorname*{Shade}\left(  F^{\prime}\right)
=\operatorname*{Shade}F$ (since $F^{\prime}=F\cup\left\{  \varepsilon\left(
F\right)  \right\}  $). Hence, (\ref{pf.thm.elser-Shade0.Fpri}) yields
$F^{\prime}\in\mathcal{A}$. Also, (\ref{pf.thm.elser-Shade0.eps=eps}) (applied
to $F^{\prime}$ and $F$ instead of $F_{1}$ and $F_{2}$) yields $\varepsilon
\left(  F^{\prime}\right)  =\varepsilon\left(  F\right)  $ (since
$F\in\mathcal{A}$ and $F^{\prime}\in\mathcal{A}$ and $\operatorname*{Shade}%
\left(  F^{\prime}\right)  =\operatorname*{Shade}F$). Hence, $\varepsilon
\left(  F^{\prime}\right)  =\varepsilon\left(  F\right)  \in\left\{
\varepsilon\left(  F\right)  \right\}  \subseteq F\cup\left\{  \varepsilon
\left(  F\right)  \right\}  =F^{\prime}$ (since $F^{\prime}=F\cup\left\{
\varepsilon\left(  F\right)  \right\}  $). Therefore, $F^{\prime}$ is a
$P\in\mathcal{A}$ satisfying $\varepsilon\left(  P\right)  \in P$ (since
$F^{\prime}\in\mathcal{A}$ and $\varepsilon\left(  F^{\prime}\right)  \in
F^{\prime}$). In other words, $F^{\prime}\in\left\{  P\in\mathcal{A}%
\ \mid\ \varepsilon\left(  P\right)  \in P\right\}  $. In other words,
$F^{\prime}\in\mathcal{A}_{+}$ (since $\mathcal{A}_{+}=\left\{  P\in
\mathcal{A}\ \mid\ \varepsilon\left(  P\right)  \in P\right\}  $). In other
words, $F\cup\left\{  \varepsilon\left(  F\right)  \right\}  \in
\mathcal{A}_{+}$ (since $F^{\prime}=F\cup\left\{  \varepsilon\left(  F\right)
\right\}  $). Qed.}. Thus, we can define a map%
\begin{align*}
\Psi:\mathcal{A}_{-}  &  \rightarrow\mathcal{A}_{+},\\
F  &  \mapsto F\cup\left\{  \varepsilon\left(  F\right)  \right\}  .
\end{align*}


We have $\Phi\circ\Psi=\operatorname*{id}$\ \ \ \ \footnote{\textit{Proof.}
Let $F\in\mathcal{A}_{-}$. Thus, $F\in\mathcal{A}_{-}=\left\{  P\in
\mathcal{A}\ \mid\ \varepsilon\left(  P\right)  \notin P\right\}  $. In other
words, $F$ is a $P\in\mathcal{A}$ satisfying $\varepsilon\left(  P\right)
\notin P$. In other words, $F$ is an element of $\mathcal{A}$ and satisfies
$\varepsilon\left(  F\right)  \notin F$.
\par
We set $u=\varepsilon\left(  F\right)  $ and $F^{\prime}=F\cup\left\{
u\right\}  $. Thus, $u=\varepsilon\left(  F\right)  \notin F$ and $F^{\prime
}=F\cup\left\{  u\right\}  =F\cup\left\{  \varepsilon\left(  F\right)
\right\}  $ (since $u=\varepsilon\left(  F\right)  $).
\par
We have $F\in\mathcal{A}=\left\{  P\subseteq E\ \mid\ G\not \subseteq
\operatorname*{Shade}P\right\}  $. In other words, $F$ is a $P\subseteq E$
satisfying $G\not \subseteq \operatorname*{Shade}P$. Thus, $F$ is a subset of
$E$.
\par
Recall that $\varepsilon\left(  F\right)  $ is the smallest edge $e\in
G\setminus\operatorname*{Shade}F$ (by the definition of $\varepsilon\left(
F\right)  $). Hence, $\varepsilon\left(  F\right)  \in G\setminus
\operatorname*{Shade}F$. In other words, $u\in G\setminus\operatorname*{Shade}%
F$ (since $u=\varepsilon\left(  F\right)  $). In other words, $u\in G$ and
$u\notin\operatorname*{Shade}F$. Thus, $u\in G\subseteq E$ and $u\notin%
\operatorname*{Shade}F$. Therefore, $F\cup\left\{  u\right\}  $ is a subset of
$E$ (since $F$ is a subset of $E$, and since $u\in E$). In other words,
$F^{\prime}$ is a subset of $E$ (since $F^{\prime}=F\cup\left\{  u\right\}
$).
\par
Furthermore, (\ref{eq.lem.Shade-tog.union}) yields $\operatorname*{Shade}%
\left(  F\cup\left\{  u\right\}  \right)  =\operatorname*{Shade}F$ (since
$u\notin\operatorname*{Shade}F$). This can be rewritten as
$\operatorname*{Shade}\left(  F^{\prime}\right)  =\operatorname*{Shade}F$
(since $F^{\prime}=F\cup\left\{  u\right\}  $). Hence,
(\ref{pf.thm.elser-Shade0.Fpri}) yields $F^{\prime}\in\mathcal{A}$. Also,
(\ref{pf.thm.elser-Shade0.eps=eps}) (applied to $F^{\prime}$ and $F$ instead
of $F_{1}$ and $F_{2}$) yields $\varepsilon\left(  F^{\prime}\right)
=\varepsilon\left(  F\right)  $ (since $F\in\mathcal{A}$ and $F^{\prime}%
\in\mathcal{A}$ and $\operatorname*{Shade}\left(  F^{\prime}\right)
=\operatorname*{Shade}F$).
\par
The definition of $\Psi$ yields $\Psi\left(  F\right)  =F\cup\left\{
\varepsilon\left(  F\right)  \right\}  =F^{\prime}$ (since $F^{\prime}%
=F\cup\left\{  \varepsilon\left(  F\right)  \right\}  $), so that $F^{\prime
}=\Psi\left(  F\right)  \in\mathcal{A}_{+}$ (since $\Psi$ is a map from
$\mathcal{A}_{-}$ to $\mathcal{A}_{+}$). Hence, the definition of $\Phi$
yields
\begin{align*}
\Phi\left(  F^{\prime}\right)   &  =F^{\prime}\setminus\left\{  \varepsilon
\left(  F^{\prime}\right)  \right\}  =\underbrace{F^{\prime}}_{=F\cup\left\{
u\right\}  }\setminus\left\{  u\right\}  \ \ \ \ \ \ \ \ \ \ \left(
\text{since }\varepsilon\left(  F^{\prime}\right)  =\varepsilon\left(
F\right)  =u\right) \\
&  =\left(  F\cup\left\{  u\right\}  \right)  \setminus\left\{  u\right\}
=F\ \ \ \ \ \ \ \ \ \ \left(  \text{since }u\notin F\right)  .
\end{align*}
In view of $\Psi\left(  F\right)  =F^{\prime}$, this can be rewritten as
$\Phi\left(  \Psi\left(  F\right)  \right)  =F$. Thus, $\left(  \Phi\circ
\Psi\right)  \left(  F\right)  =\Phi\left(  \Psi\left(  F\right)  \right)
=F=\operatorname*{id}\left(  F\right)  $.
\par
Forget that we fixed $F$. We thus have shown that $\left(  \Phi\circ
\Psi\right)  \left(  F\right)  =\operatorname*{id}\left(  F\right)  $ for each
$F\in\mathcal{A}_{-}$. In other words, $\Phi\circ\Psi=\operatorname*{id}$.}
and $\Psi\circ\Phi=\operatorname*{id}$\ \ \ \ \footnote{\textit{Proof.} Let
$F\in\mathcal{A}_{+}$. Thus, $F\in\mathcal{A}_{+}=\left\{  P\in\mathcal{A}%
\ \mid\ \varepsilon\left(  P\right)  \in P\right\}  $. In other words, $F$ is
a $P\in\mathcal{A}$ satisfying $\varepsilon\left(  P\right)  \in P$. In other
words, $F$ is an element of $\mathcal{A}$ and satisfies $\varepsilon\left(
F\right)  \in F$.
\par
We set $u=\varepsilon\left(  F\right)  $ and $F^{\prime}=F\setminus\left\{
u\right\}  $. Hence, $u=\varepsilon\left(  F\right)  \in F$ and $F^{\prime
}=F\setminus\left\{  u\right\}  =F\setminus\left\{  \varepsilon\left(
F\right)  \right\}  $ (since $u=\varepsilon\left(  F\right)  $).
\par
We have $F\in\mathcal{A}_{+}=\left\{  P\in\mathcal{A}\ \mid\ \varepsilon
\left(  P\right)  \in P\right\}  \subseteq\mathcal{A}=\left\{  P\subseteq
E\ \mid\ G\not \subseteq \operatorname*{Shade}P\right\}  $. In other words,
$F$ is a $P\subseteq E$ satisfying $G\not \subseteq \operatorname*{Shade}P$.
Thus, $F$ is a subset of $E$. Therefore, $F\setminus\left\{  u\right\}  $ is a
subset of $E$ as well (since $F\setminus\left\{  u\right\}  \subseteq
F\subseteq E$). In other words, $F^{\prime}$ is a subset of $E$ (since
$F^{\prime}=F\setminus\left\{  u\right\}  $).
\par
Recall that $\varepsilon\left(  F\right)  $ is the smallest edge $e\in
G\setminus\operatorname*{Shade}F$ (by the definition of $\varepsilon\left(
F\right)  $). Hence, $\varepsilon\left(  F\right)  \in G\setminus
\operatorname*{Shade}F$. In other words, $u\in G\setminus\operatorname*{Shade}%
F$ (since $u=\varepsilon\left(  F\right)  $). In other words, $u\in G$ and
$u\notin\operatorname*{Shade}F$. Thus, $u\in G\subseteq E$ and $u\notin%
\operatorname*{Shade}F$. Therefore, (\ref{eq.lem.Shade-tog.diff}) yields
$\operatorname*{Shade}\left(  F\setminus\left\{  u\right\}  \right)
=\operatorname*{Shade}F$. This can be rewritten as $\operatorname*{Shade}%
\left(  F^{\prime}\right)  =\operatorname*{Shade}F$ (since $F^{\prime
}=F\setminus\left\{  u\right\}  $). Hence, (\ref{pf.thm.elser-Shade0.Fpri})
yields $F^{\prime}\in\mathcal{A}$. Also, (\ref{pf.thm.elser-Shade0.eps=eps})
(applied to $F^{\prime}$ and $F$ instead of $F_{1}$ and $F_{2}$) yields
$\varepsilon\left(  F^{\prime}\right)  =\varepsilon\left(  F\right)  $ (since
$F\in\mathcal{A}$ and $F^{\prime}\in\mathcal{A}$ and $\operatorname*{Shade}%
\left(  F^{\prime}\right)  =\operatorname*{Shade}F$).
\par
The definition of $\Phi$ yields $\Phi\left(  F\right)  =F\setminus\left\{
\varepsilon\left(  F\right)  \right\}  =F^{\prime}$ (since $F^{\prime
}=F\setminus\left\{  \varepsilon\left(  F\right)  \right\}  $), so that
$F^{\prime}=\Phi\left(  F\right)  \in\mathcal{A}_{-}$ (since $\Phi$ is a map
from $\mathcal{A}_{+}$ to $\mathcal{A}_{-}$). Hence, the definition of $\Psi$
yields
\begin{align*}
\Psi\left(  F^{\prime}\right)   &  =F^{\prime}\cup\left\{  \varepsilon\left(
F^{\prime}\right)  \right\}  =\underbrace{F^{\prime}}_{=F\setminus\left\{
u\right\}  }\cup\left\{  u\right\}  \ \ \ \ \ \ \ \ \ \ \left(  \text{since
}\varepsilon\left(  F^{\prime}\right)  =\varepsilon\left(  F\right)  =u\right)
\\
&  =\left(  F\setminus\left\{  u\right\}  \right)  \cup\left\{  u\right\}
=F\ \ \ \ \ \ \ \ \ \ \left(  \text{since }u\in F\right)  .
\end{align*}
In view of $\Phi\left(  F\right)  =F^{\prime}$, this can be rewritten as
$\Psi\left(  \Phi\left(  F\right)  \right)  =F$. Thus, $\left(  \Psi\circ
\Phi\right)  \left(  F\right)  =\Psi\left(  \Phi\left(  F\right)  \right)
=F=\operatorname*{id}\left(  F\right)  $.
\par
Forget that we fixed $F$. We thus have shown that $\left(  \Psi\circ
\Phi\right)  \left(  F\right)  =\operatorname*{id}\left(  F\right)  $ for each
$F\in\mathcal{A}_{+}$. In other words, $\Psi\circ\Phi=\operatorname*{id}$.}.
Thus, the maps $\Phi$ and $\Psi$ are mutually inverse. Hence, the map $\Phi$
is invertible, thus a bijection.

Moreover, each $F\in\mathcal{A}_{+}$ satisfies%
\begin{equation}
\left(  -1\right)  ^{\left\vert \Phi\left(  F\right)  \right\vert }=-\left(
-1\right)  ^{\left\vert F\right\vert }. \label{pf.thm.elser-Shade0.signrev}%
\end{equation}
\footnote{\textit{Proof of (\ref{pf.thm.elser-Shade0.signrev}):} Let
$F\in\mathcal{A}_{+}$. Thus, $F\in\mathcal{A}_{+}=\left\{  P\in\mathcal{A}%
\ \mid\ \varepsilon\left(  P\right)  \in P\right\}  $. In other words, $F$ is
a $P\in\mathcal{A}$ satisfying $\varepsilon\left(  P\right)  \in P$. In other
words, $F$ is an element of $\mathcal{A}$ and satisfies $\varepsilon\left(
F\right)  \in F$.
\par
The definition of $\Phi$ yields $\Phi\left(  F\right)  =F\setminus\left\{
\varepsilon\left(  F\right)  \right\}  $. Hence, $\left\vert \Phi\left(
F\right)  \right\vert =\left\vert F\setminus\left\{  \varepsilon\left(
F\right)  \right\}  \right\vert =\left\vert F\right\vert -1$ (since
$\varepsilon\left(  F\right)  \in F$). Therefore, $\left(  -1\right)
^{\left\vert \Phi\left(  F\right)  \right\vert }=\left(  -1\right)
^{\left\vert F\right\vert -1}=-\left(  -1\right)  ^{\left\vert F\right\vert }%
$. This proves (\ref{pf.thm.elser-Shade0.signrev}).}

Now,
\begin{align*}
&  \underbrace{\sum_{\substack{F\subseteq E;\\G\not \subseteq
\operatorname*{Shade}F}}}_{\substack{=\sum_{F\in\left\{  P\subseteq
E\ \mid\ G\not \subseteq \operatorname*{Shade}P\right\}  }\\=\sum
_{F\in\mathcal{A}}\\\text{(since }\left\{  P\subseteq E\ \mid
\ G\not \subseteq \operatorname*{Shade}P\right\}  =\mathcal{A}\text{)}%
}}\left(  -1\right)  ^{\left\vert F\right\vert }\\
&  =\sum_{F\in\mathcal{A}}\left(  -1\right)  ^{\left\vert F\right\vert
}=\underbrace{\sum_{\substack{F\in\mathcal{A};\\\varepsilon\left(  F\right)
\in F}}}_{\substack{=\sum_{F\in\left\{  P\in\mathcal{A}\ \mid\ \varepsilon
\left(  P\right)  \in P\right\}  }\\=\sum_{F\in\mathcal{A}_{+}}\\\text{(since
}\left\{  P\in\mathcal{A}\ \mid\ \varepsilon\left(  P\right)  \in P\right\}
=\mathcal{A}_{+}\text{)}}}\left(  -1\right)  ^{\left\vert F\right\vert
}+\underbrace{\sum_{\substack{F\in\mathcal{A};\\\varepsilon\left(  F\right)
\notin F}}}_{\substack{=\sum_{F\in\left\{  P\in\mathcal{A}\ \mid
\ \varepsilon\left(  P\right)  \notin P\right\}  }\\=\sum_{F\in\mathcal{A}%
_{-}}\\\text{(since }\left\{  P\in\mathcal{A}\ \mid\ \varepsilon\left(
P\right)  \notin P\right\}  =\mathcal{A}_{-}\text{)}}}\left(  -1\right)
^{\left\vert F\right\vert }\\
&  \ \ \ \ \ \ \ \ \ \ \left(
\begin{array}
[c]{c}%
\text{since each }F\in\mathcal{A}\text{ satisfies either }\varepsilon\left(
F\right)  \in F\text{ or }\varepsilon\left(  F\right)  \notin F\\
\text{(but not both at the same time)}%
\end{array}
\right) \\
&  =\sum_{F\in\mathcal{A}_{+}}\left(  -1\right)  ^{\left\vert F\right\vert
}+\sum_{F\in\mathcal{A}_{-}}\left(  -1\right)  ^{\left\vert F\right\vert
}=\sum_{F\in\mathcal{A}_{+}}\left(  -1\right)  ^{\left\vert F\right\vert
}+\sum_{F\in\mathcal{A}_{+}}\underbrace{\left(  -1\right)  ^{\left\vert
\Phi\left(  F\right)  \right\vert }}_{\substack{=-\left(  -1\right)
^{\left\vert F\right\vert }\\\text{(by (\ref{pf.thm.elser-Shade0.signrev}))}%
}}\\
&  \ \ \ \ \ \ \ \ \ \ \left(
\begin{array}
[c]{c}%
\text{here, we have substituted }\Phi\left(  F\right)  \text{ for }F\text{ in
the second sum,}\\
\text{since the map }\Phi:\mathcal{A}_{+}\rightarrow\mathcal{A}_{-}\text{ is a
bijection}%
\end{array}
\right) \\
&  =\sum_{F\in\mathcal{A}_{+}}\left(  -1\right)  ^{\left\vert F\right\vert
}+\underbrace{\sum_{F\in\mathcal{A}_{+}}\left(  -\left(  -1\right)
^{\left\vert F\right\vert }\right)  }_{=-\sum_{F\in\mathcal{A}_{+}}\left(
-1\right)  ^{\left\vert F\right\vert }}=\sum_{F\in\mathcal{A}_{+}}\left(
-1\right)  ^{\left\vert F\right\vert }-\sum_{F\in\mathcal{A}_{+}}\left(
-1\right)  ^{\left\vert F\right\vert }=0.
\end{align*}
This proves Theorem \ref{thm.elser-Shade0}.
\end{proof}
\end{verlong}

In order to derive Theorem \ref{thm.elser-Shade1} from Theorem
\ref{thm.elser-Shade0}, we need the following innocent lemma -- which is one
of the simplest facts in enumerative combinatorics:

\begin{lemma}
\label{lem.toggle}Let $U$ be a finite set with $U\neq\varnothing$. Then,%
\[
\sum_{F\subseteq U}\left(  -1\right)  ^{\left\vert F\right\vert }=0.
\]

\end{lemma}

Lemma \ref{lem.toggle} can easily be derived from the fact that $\sum
_{k=0}^{n}\left(  -1\right)  ^{k}\dbinom{n}{k}=0$ for any positive integer $n$
(as follows readily from the binomial theorem). However, keeping true to the
spirit of this paper, let us give a bijective proof for it:

\begin{proof}
[Proof of Lemma \ref{lem.toggle}.]This is a standard argument that underlies
many combinatorial proofs of alternating sum identities (see, for example,
\cite[proof of (2.4)]{Sagan20} or \cite[proof of (1)]{BenQui08}). For the sake
of completeness, let us nevertheless recall it.

We have $U\neq\varnothing$; hence, there exists some $u\in U$. Consider this
$u$.

\begin{vershort}
It is easy to see that the maps%
\begin{align*}
\Phi:\left\{  F\subseteq U\ \mid\ u\in F\right\}   &  \rightarrow\left\{
F\subseteq U\ \mid\ u\notin F\right\}  ,\\
F  &  \mapsto F\setminus\left\{  u\right\}
\end{align*}
and%
\begin{align*}
\Psi:\left\{  F\subseteq U\ \mid\ u\notin F\right\}   &  \rightarrow\left\{
F\subseteq U\ \mid\ u\in F\right\}  ,\\
F  &  \mapsto F\cup\left\{  u\right\}
\end{align*}
are well-defined.
\end{vershort}

\begin{verlong}
For every $F\subseteq U$ satisfying $u\in F$, we have $F\setminus\left\{
u\right\}  \subseteq U$\ \ \ \ \footnote{since $F\setminus\left\{  u\right\}
\subseteq F\subseteq U$} and $u\notin F\setminus\left\{  u\right\}  $. Hence,
the map%
\begin{align*}
\Phi:\left\{  F\subseteq U\ \mid\ u\in F\right\}   &  \rightarrow\left\{
F\subseteq U\ \mid\ u\notin F\right\}  ,\\
F  &  \mapsto F\setminus\left\{  u\right\}
\end{align*}
is well-defined. Consider this map $\Phi$.

For every $F\subseteq U$ satisfying $u\notin F$, we have $F\cup\left\{
u\right\}  \subseteq U$\ \ \ \ \footnote{since $F\subseteq U$ and $u\in U$}
and $u\in F\cup\left\{  u\right\}  $. Hence, the map%
\begin{align*}
\Psi:\left\{  F\subseteq U\ \mid\ u\notin F\right\}   &  \rightarrow\left\{
F\subseteq U\ \mid\ u\in F\right\}  ,\\
F  &  \mapsto F\cup\left\{  u\right\}
\end{align*}
is well-defined. Consider this map $\Psi$.
\end{verlong}

The maps $\Phi$ and $\Psi$ we just defined are clearly mutually inverse. Thus,
they are invertible, i.e., are bijections. Hence, in particular, $\Phi$ is a
bijection. Thus, we can substitute $\Phi\left(  F\right)  $ for $F$ in the sum
$\sum_{\substack{F\subseteq U;\\u\notin F}}\left(  -1\right)  ^{\left\vert
F\right\vert }$. We thus obtain%
\begin{align*}
\sum_{\substack{F\subseteq U;\\u\notin F}}\left(  -1\right)  ^{\left\vert
F\right\vert }  &  =\sum_{\substack{F\subseteq U;\\u\in F}}\underbrace{\left(
-1\right)  ^{\left\vert \Phi\left(  F\right)  \right\vert }}%
_{\substack{=\left(  -1\right)  ^{\left\vert F\setminus\left\{  u\right\}
\right\vert }\\\text{(since }\Phi\left(  F\right)  =F\setminus\left\{
u\right\}  \\\text{(by the definition of }\Phi\text{))}}}=\sum
_{\substack{F\subseteq U;\\u\in F}}\underbrace{\left(  -1\right)  ^{\left\vert
F\setminus\left\{  u\right\}  \right\vert }}_{\substack{=\left(  -1\right)
^{\left\vert F\right\vert -1}\\\text{(since }\left\vert F\setminus\left\{
u\right\}  \right\vert =\left\vert F\right\vert -1\\\text{(because }u\in
F\text{))}}}\\
&  =\sum_{\substack{F\subseteq U;\\u\in F}}\underbrace{\left(  -1\right)
^{\left\vert F\right\vert -1}}_{=-\left(  -1\right)  ^{\left\vert F\right\vert
}}=-\sum_{\substack{F\subseteq U;\\u\in F}}\left(  -1\right)  ^{\left\vert
F\right\vert }.
\end{align*}
However, each $F\subseteq U$ satisfies either $u\in F$ or $u\notin F$ (but not
both). Hence,%
\[
\sum_{F\subseteq U}\left(  -1\right)  ^{\left\vert F\right\vert }%
=\sum_{\substack{F\subseteq U;\\u\in F}}\left(  -1\right)  ^{\left\vert
F\right\vert }+\underbrace{\sum_{\substack{F\subseteq U;\\u\notin F}}\left(
-1\right)  ^{\left\vert F\right\vert }}_{=-\sum_{\substack{F\subseteq U;\\u\in
F}}\left(  -1\right)  ^{\left\vert F\right\vert }}=\sum_{\substack{F\subseteq
U;\\u\in F}}\left(  -1\right)  ^{\left\vert F\right\vert }-\sum
_{\substack{F\subseteq U;\\u\in F}}\left(  -1\right)  ^{\left\vert
F\right\vert }=0.
\]
This proves Lemma \ref{lem.toggle}.
\end{proof}

We can now easily derive Theorem \ref{thm.elser-Shade1} from Theorem
\ref{thm.elser-Shade0}:

\begin{proof}
[Proof of Theorem \ref{thm.elser-Shade1}.]Each subset $F$ of $E$ satisfies
either $G\subseteq\operatorname*{Shade}F$ or $G\not \subseteq
\operatorname*{Shade}F$ (but not both at the same time). Hence,%
\[
\sum_{F\subseteq E}\left(  -1\right)  ^{\left\vert F\right\vert }%
=\sum_{\substack{F\subseteq E;\\G\subseteq\operatorname*{Shade}F}}\left(
-1\right)  ^{\left\vert F\right\vert }+\underbrace{\sum_{\substack{F\subseteq
E;\\G\not \subseteq \operatorname*{Shade}F}}\left(  -1\right)  ^{\left\vert
F\right\vert }}_{\substack{=0\\\text{(by Theorem \ref{thm.elser-Shade0})}%
}}=\sum_{\substack{F\subseteq E;\\G\subseteq\operatorname*{Shade}F}}\left(
-1\right)  ^{\left\vert F\right\vert }.
\]
Therefore,%
\[
\sum_{\substack{F\subseteq E;\\G\subseteq\operatorname*{Shade}F}}\left(
-1\right)  ^{\left\vert F\right\vert }=\sum_{F\subseteq E}\left(  -1\right)
^{\left\vert F\right\vert }=0
\]
(by Lemma \ref{lem.toggle}, applied to $U=E$). This proves Theorem
\ref{thm.elser-Shade1}.
\end{proof}

\subsection{Proving Theorem \ref{thm.elser0}}

Theorem \ref{thm.elser0} is now a simple particular case of Theorem
\ref{thm.elser-Shade1}:

\begin{vershort}
\begin{proof}
[Proof of Theorem \ref{thm.elser0}.]Let $G$ be the set $E$. Thus, $G=E$.
Hence, for each subset $F$ of $E$, we have the following chain of logical
equivalences:%
\begin{align*}
\left(  G\subseteq\operatorname*{Shade}F\right)  \  &  \Longleftrightarrow
\ \left(  E\subseteq\operatorname*{Shade}F\right) \\
&  \Longleftrightarrow\ \left(  \text{each }u\in E\text{ satisfies }%
u\in\operatorname*{Shade}F\right) \\
&  \Longleftrightarrow\ \left(  \text{each }u\in E\text{ has the property that
}F\text{ infects }u\right) \\
&  \ \ \ \ \ \ \ \ \ \ \ \ \ \ \ \ \ \ \ \ \left(  \text{since }%
\operatorname*{Shade}F=\left\{  e\in E\ \mid\ F\text{ infects }e\right\}
\right) \\
&  \Longleftrightarrow\ \left(  F\text{ infects each }u\in E\right) \\
&  \Longleftrightarrow\ \left(  F\text{ is pandemic}\right)
\ \ \ \ \ \ \ \ \ \ \left(  \text{by the definition of \textquotedblleft
pandemic\textquotedblright}\right)  .
\end{align*}
Thus, the summation sign \textquotedblleft$\sum_{\substack{F\subseteq
E;\\G\subseteq\operatorname*{Shade}F}}$\textquotedblright\ can be rewritten as
\textquotedblleft$\sum_{\substack{F\subseteq E\text{ is}\\\text{pandemic}}%
}$\textquotedblright. Hence,%
\[
\sum_{\substack{F\subseteq E;\\G\subseteq\operatorname*{Shade}F}}\left(
-1\right)  ^{\left\vert F\right\vert }=\sum_{\substack{F\subseteq E\text{
is}\\\text{pandemic}}}\left(  -1\right)  ^{\left\vert F\right\vert }.
\]
But the left hand side of this equality is $0$ (by Theorem
\ref{thm.elser-Shade1}). Hence, its right hand side is $0$ as well. This
proves Theorem \ref{thm.elser0}.
\end{proof}
\end{vershort}

\begin{verlong}
\begin{proof}
[Proof of Theorem \ref{thm.elser0}.]Let $G$ be the set $E$. Thus, $G=E$.
Hence, for each subset $F$ of $E$, we have the following chain of logical
equivalences:%
\begin{align*}
\left(  G\subseteq\operatorname*{Shade}F\right)  \  &  \Longleftrightarrow
\ \left(  E\subseteq\operatorname*{Shade}F\right) \\
&  \Longleftrightarrow\ \left(  \text{each }u\in E\text{ satisfies }%
u\in\operatorname*{Shade}F\right) \\
&  \Longleftrightarrow\ \left(  \text{each }u\in E\text{ satisfies }%
u\in\left\{  e\in E\ \mid\ F\text{ infects }e\right\}  \right) \\
&  \ \ \ \ \ \ \ \ \ \ \ \ \ \ \ \ \ \ \ \ \left(  \text{since }%
\operatorname*{Shade}F=\left\{  e\in E\ \mid\ F\text{ infects }e\right\}
\right) \\
&  \Longleftrightarrow\ \left(  \text{each }u\in E\text{ has the property that
}F\text{ infects }u\right) \\
&  \Longleftrightarrow\ \left(  F\text{ infects each }u\in E\right) \\
&  \Longleftrightarrow\ \left(  F\text{ infects each }e\in E\right) \\
&  \Longleftrightarrow\ \left(  F\text{ is pandemic}\right)
\ \ \ \ \ \ \ \ \ \ \left(  \text{by the definition of \textquotedblleft
pandemic\textquotedblright}\right)  .
\end{align*}
Thus, the summation sign \textquotedblleft$\sum_{\substack{F\subseteq
E;\\G\subseteq\operatorname*{Shade}F}}$\textquotedblright\ can be rewritten as
\textquotedblleft$\sum_{\substack{F\subseteq E\text{ is}\\\text{pandemic}}%
}$\textquotedblright. Hence,%
\[
\sum_{\substack{F\subseteq E;\\G\subseteq\operatorname*{Shade}F}}\left(
-1\right)  ^{\left\vert F\right\vert }=\sum_{\substack{F\subseteq E\text{
is}\\\text{pandemic}}}\left(  -1\right)  ^{\left\vert F\right\vert }.
\]
Therefore,%
\[
\sum_{\substack{F\subseteq E\text{ is}\\\text{pandemic}}}\left(  -1\right)
^{\left\vert F\right\vert }=\sum_{\substack{F\subseteq E;\\G\subseteq
\operatorname*{Shade}F}}\left(  -1\right)  ^{\left\vert F\right\vert }=0
\]
(by Theorem \ref{thm.elser-Shade1}). This proves Theorem \ref{thm.elser0}.
\end{proof}
\end{verlong}

\section{Vertex infection and other variants}

In our study of graphs so far, we have barely ever mentioned vertices (even
though they are, of course, implicit in the notion of a path). It may appear
somewhat strange to talk about a subset infecting an edge, when the infection
is spread from vertex to vertex. One might thus wonder if there is also a
vertex counterpart of Theorem \ref{thm.elser0}. So let us define analogues of
our notions for vertices:

If $F\subseteq V$, then an $F$\emph{-vertex-path} shall mean a path of
$\Gamma$ such that all vertices of the path except (possibly) for its two
endpoints belong to $F$. (Thus, if a path has only one edge or none, then it
automatically is an $F$-vertex-path.)

If $w\in V\setminus\left\{  v\right\}  $ is any vertex and $F\subseteq
V\setminus\left\{  v\right\}  $ is any subset, then we say that $F$
\emph{vertex-infects }$w$ if there exists an $F$-vertex-path from $v$ to $w$.
(This is always true when $w$ is a neighbor of $v$.)

A subset $F\subseteq V\setminus\left\{  v\right\}  $ is said to be
\emph{vertex-pandemic} if it vertex-infects each vertex $w\in V\setminus
\left\{  v\right\}  $.

\begin{example}
Let $\Gamma$ be as in Example \ref{exa.example2}. Then, the path
$v\overset{1}{\longrightarrow}p\overset{2}{\longrightarrow}q$ is an
$F$-vertex-path for any subset $F\subseteq V$ that satisfies $p\in F$. The
subset $\left\{  p\right\}  $ of $V\setminus\left\{  v\right\}  $
vertex-infects each vertex (for example, $v\overset{1}{\longrightarrow
}p\overset{2}{\longrightarrow}q$ is a $\left\{  p\right\}  $-vertex-path from
$v$ to $q$, and $v\overset{4}{\longrightarrow}w$ is a $\left\{  p\right\}
$-vertex-path from $v$ to $w$), and thus is vertex-pandemic. The
vertex-pandemic subsets of $V\setminus\left\{  v\right\}  $ are the sets%
\[
\left\{  p\right\}  ,\ \ \left\{  w\right\}  ,\ \ \left\{  p,q\right\}
,\ \ \left\{  p,w\right\}  ,\ \ \left\{  q,w\right\}  ,\ \ \left\{
p,q,w\right\}  .
\]

\end{example}

We now have the following analogue of Theorem \ref{thm.elser0}:

\begin{theorem}
\label{thm.elser0-vert}Assume that $V\setminus\left\{  v\right\}
\neq\varnothing$. Then,%
\[
\sum_{\substack{F\subseteq V\setminus\left\{  v\right\}  \text{ is}%
\\\text{vertex-pandemic}}}\left(  -1\right)  ^{\left\vert F\right\vert }=0.
\]

\end{theorem}

\begin{proof}
[Proof of Theorem \ref{thm.elser0-vert}.]With just a few easy modifications,
our above proof of Theorem \ref{thm.elser0} can be repurposed as a proof of
Theorem \ref{thm.elser0-vert}. Namely:

\begin{itemize}
\item We need to replace \textquotedblleft edge\textquotedblright\ by
\textquotedblleft vertex\textquotedblright\ throughout the argument (including
Definition \ref{def.Shade}, Lemma \ref{lem.Shade-monoton}, Lemma
\ref{lem.Shade-tog}, Theorem \ref{thm.elser-Shade1} and Theorem
\ref{thm.elser-Shade0}), as well as replace $E$ by $V\setminus\left\{
v\right\}  $.

\item The words \textquotedblleft$F$-path\textquotedblright, \textquotedblleft
infects\textquotedblright\ and \textquotedblleft pandemic\textquotedblright%
\ have to be replaced by \textquotedblleft$F$-vertex-path\textquotedblright,
\textquotedblleft vertex-infects\textquotedblright\ and \textquotedblleft
vertex-pandemic\textquotedblright, respectively.

\item In the proofs of Lemma \ref{lem.Shade-monoton} and Lemma
\ref{lem.Shade-tog}, the words \textquotedblleft an endpoint
of\textquotedblright\ (as well as \textquotedblleft any endpoint
of\textquotedblright\ and \textquotedblleft some endpoint of\textquotedblright%
) need to be removed (since the notion of \textquotedblleft
vertex-infects\textquotedblright\ is defined not in terms of paths to an
endpoint of a given edge, but in terms of paths to a given vertex).

\item In the proof of Lemma \ref{lem.Shade-tog}, specifically in the proof of
(\ref{eq.lem.Shade-tog.union}), the path $\pi$ is now cut not by removing the
edge $u$, but by splitting the path $\pi$ at the vertex $u$.
\end{itemize}

\noindent The reader may check that these changes result in a valid proof of
Theorem \ref{thm.elser0-vert}.
\end{proof}

Another variant of Theorem \ref{thm.elser0} (and Theorem
\ref{thm.elser-Shade1} and Theorem \ref{thm.elser-Shade0}) is obtained by
replacing the undirected graph $\Gamma$ with a directed graph (while, of
course, replacing paths by directed paths). More generally, we can replace
$\Gamma$ by a \textquotedblleft hybrid\textquotedblright\ graph with some
directed and some undirected edges.\footnote{We understand that a directed
edge still has two endpoints: its source and its target.} No changes are
required to the above proofs. Yet another variation can be obtained by
replacing \textquotedblleft endpoint\textquotedblright\ by \textquotedblleft
source\textquotedblright\ (for directed edges). We cannot, however, replace
\textquotedblleft endpoint\textquotedblright\ by \textquotedblleft
target\textquotedblright.

\section{\label{sec.abstract}An abstract perspective}

Seeing how little graph theory we have used in proving Theorem
\ref{thm.elser0}, and how easily the same argument adapted to Theorem
\ref{thm.elser0-vert}, we get the impression that there might be some general
theory lurking behind it. What follows is an attempt at building this theory.

\subsection{Shade maps}

Let $\mathcal{P}\left(  E\right)  $ denote the power set of $E$. In Definition
\ref{def.Shade}, we have encoded the \textquotedblleft
infects\textquotedblright\ relation as a map $\operatorname*{Shade}%
:\mathcal{P}\left(  E\right)  \rightarrow\mathcal{P}\left(  E\right)  $
defined by%
\[
\operatorname*{Shade}F=\left\{  e\in E\ \mid\ F\text{ infects }e\right\}  .
\]
As we recall, Theorem \ref{thm.elser-Shade1} (a generalization of Theorem
\ref{thm.elser0}) states that
\begin{equation}
\sum_{\substack{F\subseteq E;\\G\subseteq\operatorname*{Shade}F}}\left(
-1\right)  ^{\left\vert F\right\vert }=0 \label{eq.Shade.goal}%
\end{equation}
for any $G\subseteq E$, under the assumption that $E\neq\varnothing$.

To generalize this, we forget about the graph $\Gamma$ and the map
$\operatorname*{Shade}$, and instead start with an \textbf{arbitrary} finite
set $E$. (This set $E$ corresponds to the set $E$ in Theorem \ref{thm.elser0}
and to the set $V\setminus\left\{  v\right\}  $ in Theorem
\ref{thm.elser0-vert}.) Let $\mathcal{P}\left(  E\right)  $ be the power set
of $E$. Let $\operatorname*{Shade}:\mathcal{P}\left(  E\right)  \rightarrow
\mathcal{P}\left(  E\right)  $ be an arbitrary map (meant to generalize the
map $\operatorname*{Shade}$ from the previous paragraph). We may now ask:

\begin{question}
What (combinatorial) properties must $\operatorname*{Shade}$ satisfy in order
for (\ref{eq.Shade.goal}) to hold for any $G\subseteq E$ under the assumption
that $E\neq\varnothing$ ?
\end{question}

A partial answer to this question can be given by analyzing our above proof of
Theorem \ref{thm.elser-Shade1} and extracting what was used:

\begin{definition}
\label{def.Shade.shademap}Let $E$ be a set. A \emph{shade map} on $E$ shall
mean a map $\operatorname*{Shade}:\mathcal{P}\left(  E\right)  \rightarrow
\mathcal{P}\left(  E\right)  $ that satisfies the following two axioms:

\begin{statement}
\textit{Axiom 1:} If $F\in\mathcal{P}\left(  E\right)  $ and $u\in
E\setminus\operatorname*{Shade}F$, then $\operatorname*{Shade}\left(
F\cup\left\{  u\right\}  \right)  =\operatorname*{Shade}F$.
\end{statement}

\begin{statement}
\textit{Axiom 2:} If $F\in\mathcal{P}\left(  E\right)  $ and $u\in
E\setminus\operatorname*{Shade}F$, then $\operatorname*{Shade}\left(
F\setminus\left\{  u\right\}  \right)  =\operatorname*{Shade}F$.
\end{statement}
\end{definition}

\begin{theorem}
\label{thm.elser.genE0}Let $E$ be a finite set. Let $\operatorname*{Shade}%
:\mathcal{P}\left(  E\right)  \rightarrow\mathcal{P}\left(  E\right)  $ be a
shade map on $E$.

Assume that $E\neq\varnothing$. Let $G$ be any subset of $E$. Then,%
\[
\sum_{\substack{F\subseteq E;\\G\subseteq\operatorname*{Shade}F}}\left(
-1\right)  ^{\left\vert F\right\vert }=0.
\]

\end{theorem}

\begin{proof}
[Proof sketch.]Again, the proof is analogous to our above proof of Theorem
\ref{thm.elser-Shade1}. (This time, in the proof of Lemma \ref{lem.Shade-tog},
the equalities (\ref{eq.lem.Shade-tog.union}) and (\ref{eq.lem.Shade-tog.diff}%
) follow directly from Axiom 1 and Axiom 2, respectively.)
\end{proof}

How do shade maps relate to known concepts in the combinatorics of set
families (such as topologies, clutters, matroids, or submodular functions)?
Are they just one of these known concepts in disguise? We shall answer two
versions of this question in the following subsections. Specifically:

\begin{itemize}
\item In Subsection \ref{subsect.abstract.amatr}, we will show that
inclusion-reversing shade maps on $E$ (i.e., shade maps $\operatorname*{Shade}%
$ that satisfy $\operatorname*{Shade}B\subseteq\operatorname*{Shade}A$
whenever $A\subseteq B$) are in bijection with \emph{antimatroidal
quasi-closure operators} (a slight variant of antimatroids) on $E$.

\item In Subsection \ref{subsect.abstract.bip}, we will show that arbitrary
shade maps are in bijection with \emph{Boolean interval partitions} of
$\mathcal{P}\left(  E\right)  $ (that is, set partitions of $\mathcal{P}%
\left(  E\right)  $ into intervals of the Boolean lattice $\mathcal{P}\left(
E\right)  $).
\end{itemize}

Before we come to these characterizations, we shall however make a few
elementary remarks on shade maps.

First, we observe that Axioms 1 and 2 in Definition \ref{def.Shade.shademap}
can be weakened to the following statements:

\begin{statement}
\textit{Axiom 1':} If $F\in\mathcal{P}\left(  E\right)  $ and $u\in
E\setminus\operatorname*{Shade}F$, then $\operatorname*{Shade}\left(
F\cup\left\{  u\right\}  \right)  \subseteq\operatorname*{Shade}F$.
\end{statement}

\begin{statement}
\textit{Axiom 2':} If $F\in\mathcal{P}\left(  E\right)  $ and $u\in
E\setminus\operatorname*{Shade}F$, then $\operatorname*{Shade}\left(
F\setminus\left\{  u\right\}  \right)  \subseteq\operatorname*{Shade}F$.
\end{statement}

Axiom 1' is weaker than Axiom 1, and likewise Axiom 2' is weaker than Axiom 2.
However, Axioms 1' and 2' combined are equivalent to Axioms 1 and 2 combined:

\begin{proposition}
\label{prop.Shade.shade-2}Let $E$ be a set. Let $\operatorname*{Shade}%
:\mathcal{P}\left(  E\right)  \rightarrow\mathcal{P}\left(  E\right)  $ be any
map. Then, $\operatorname*{Shade}$ is a shade map on $E$ if and only if
$\operatorname*{Shade}$ satisfies the two Axioms 1' and 2' stated above.
\end{proposition}

\begin{vershort}
\begin{proof}
Easy exercise.
\end{proof}
\end{vershort}

\begin{verlong}
\begin{proof}
[Proof of Proposition \ref{prop.Shade.shade-2}.]$\Longrightarrow:$ Assume that
$\operatorname*{Shade}$ is a shade map on $E$. We must show that
$\operatorname*{Shade}$ satisfies the two Axioms 1' and 2'.

Indeed, $\operatorname*{Shade}$ is a shade map (by assumption), and thus
satisfies Axioms 1 and 2 in Definition \ref{def.Shade.shademap}. Now, Axiom 1'
clearly follows from Axiom 1, whereas Axiom 2' clearly follows from Axiom 2.
Hence, $\operatorname*{Shade}$ satisfies the two Axioms 1' and 2'. This proves
the \textquotedblleft$\Longrightarrow$\textquotedblright\ direction of
Proposition \ref{prop.Shade.shade-2}.

$\Longleftarrow:$ Assume that $\operatorname*{Shade}$ satisfies the two Axioms
1' and 2'. We must prove that $\operatorname*{Shade}$ is a shade map on $E$.
To this purpose, we shall show that $\operatorname*{Shade}$ satisfies Axioms 1
and 2 in Definition \ref{def.Shade.shademap}.

[\textit{Proof of Axiom 1:} Let $F\in\mathcal{P}\left(  E\right)  $ and $u\in
E\setminus\operatorname*{Shade}F$. We must show that $\operatorname*{Shade}%
\left(  F\cup\left\{  u\right\}  \right)  =\operatorname*{Shade}F$.

Indeed, if $u\in F$, then this is obvious (since we have $F\cup\left\{
u\right\}  =F$ in this case). Thus, for the rest of this proof, we WLOG assume
that $u\notin F$. Hence, $\left(  F\cup\left\{  u\right\}  \right)
\setminus\left\{  u\right\}  =F$.

We have assumed that Axioms 1' and 2' hold. Thus, Axiom 1' shows that
$\operatorname*{Shade}\left(  F\cup\left\{  u\right\}  \right)  \subseteq
\operatorname*{Shade}F$. Hence, $E\setminus\underbrace{\operatorname*{Shade}%
\left(  F\cup\left\{  u\right\}  \right)  }_{\subseteq\operatorname*{Shade}%
F}\supseteq E\setminus\operatorname*{Shade}F$, so that $E\setminus
\operatorname*{Shade}F\subseteq E\setminus\operatorname*{Shade}\left(
F\cup\left\{  u\right\}  \right)  $. Now, $u\in E\setminus
\operatorname*{Shade}F\subseteq E\setminus\operatorname*{Shade}\left(
F\cup\left\{  u\right\}  \right)  $. Hence, Axiom 2' (applied to
$F\cup\left\{  u\right\}  $ instead of $F$) yields $\operatorname*{Shade}%
\left(  \left(  F\cup\left\{  u\right\}  \right)  \setminus\left\{  u\right\}
\right)  \subseteq\operatorname*{Shade}\left(  F\cup\left\{  u\right\}
\right)  $. In view of $\left(  F\cup\left\{  u\right\}  \right)
\setminus\left\{  u\right\}  =F$, this rewrites as $\operatorname*{Shade}%
F\subseteq\operatorname*{Shade}\left(  F\cup\left\{  u\right\}  \right)  $.
Combining this with $\operatorname*{Shade}\left(  F\cup\left\{  u\right\}
\right)  \subseteq\operatorname*{Shade}F$, we obtain $\operatorname*{Shade}%
\left(  F\cup\left\{  u\right\}  \right)  =\operatorname*{Shade}F$. This
proves Axiom 1.]

[\textit{Proof of Axiom 2:} Let $F\in\mathcal{P}\left(  E\right)  $ and $u\in
E\setminus\operatorname*{Shade}F$. We must show that $\operatorname*{Shade}%
\left(  F\setminus\left\{  u\right\}  \right)  =\operatorname*{Shade}F$.

Indeed, if $u\notin F$, then this is obvious (since we have $F\setminus
\left\{  u\right\}  =F$ in this case). Thus, for the rest of this proof, we
WLOG assume that $u\in F$. Hence, $\left(  F\setminus\left\{  u\right\}
\right)  \cup\left\{  u\right\}  =F$.

We have assumed that Axioms 1' and 2' hold. Thus, Axiom 2' shows that
$\operatorname*{Shade}\left(  F\setminus\left\{  u\right\}  \right)
\subseteq\operatorname*{Shade}F$. Hence, $E\setminus
\underbrace{\operatorname*{Shade}\left(  F\setminus\left\{  u\right\}
\right)  }_{\subseteq\operatorname*{Shade}F}\supseteq E\setminus
\operatorname*{Shade}F$, so that $E\setminus\operatorname*{Shade}F\subseteq
E\setminus\operatorname*{Shade}\left(  F\setminus\left\{  u\right\}  \right)
$. Now, $u\in E\setminus\operatorname*{Shade}F\subseteq E\setminus
\operatorname*{Shade}\left(  F\setminus\left\{  u\right\}  \right)  $. Hence,
Axiom 1' (applied to $F\setminus\left\{  u\right\}  $ instead of $F$) yields
$\operatorname*{Shade}\left(  \left(  F\setminus\left\{  u\right\}  \right)
\cup\left\{  u\right\}  \right)  \subseteq\operatorname*{Shade}\left(
F\setminus\left\{  u\right\}  \right)  $. In view of $\left(  F\setminus
\left\{  u\right\}  \right)  \cup\left\{  u\right\}  =F$, this rewrites as
$\operatorname*{Shade}F\subseteq\operatorname*{Shade}\left(  F\setminus
\left\{  u\right\}  \right)  $. Combining this with $\operatorname*{Shade}%
\left(  F\setminus\left\{  u\right\}  \right)  \subseteq\operatorname*{Shade}%
F$, we obtain $\operatorname*{Shade}\left(  F\setminus\left\{  u\right\}
\right)  =\operatorname*{Shade}F$. This proves Axiom 2.]

We have now shown that the map $\operatorname*{Shade}$ satisfies both Axioms 1
and 2 in Definition \ref{def.Shade.shademap}. Hence, $\operatorname*{Shade}$
is a shade map. This proves the \textquotedblleft$\Longleftarrow
$\textquotedblright\ direction of Proposition \ref{prop.Shade.shade-2}.

Thus, the proof of Proposition \ref{prop.Shade.shade-2} is complete.
\end{proof}
\end{verlong}

Here is one further restatement of both Axioms 1' and 2': For any $F\subseteq
E$ and any two elements $u,v\in E\setminus\operatorname*{Shade}F$, we have
$v\notin\operatorname*{Shade}\left(  F\cup\left\{  u\right\}  \right)  $
(Axiom 1') and $v\notin\operatorname*{Shade}\left(  F\setminus\left\{
u\right\}  \right)  $ (Axiom 2').

Axioms 1 and 2 can also be combined into one common axiom:

\begin{statement}
\textit{Axiom 3:} If $F\in\mathcal{P}\left(  E\right)  $ and $u\in E\setminus
F$, then we have $\operatorname*{Shade}F=\operatorname*{Shade}\left(
F\cup\left\{  u\right\}  \right)  $ or $u\in\left(  \operatorname*{Shade}%
F\right)  \cap\operatorname*{Shade}\left(  F\cup\left\{  u\right\}  \right)  $.
\end{statement}

\begin{proposition}
\label{prop.Shade.shade-3}Let $E$ be a set. Let $\operatorname*{Shade}%
:\mathcal{P}\left(  E\right)  \rightarrow\mathcal{P}\left(  E\right)  $ be any
map. Then, $\operatorname*{Shade}$ is a shade map on $E$ if and only if
$\operatorname*{Shade}$ satisfies Axiom 3.
\end{proposition}

\begin{vershort}
\begin{proof}
Easy exercise.
\end{proof}
\end{vershort}

\begin{verlong}
\begin{proof}
[Proof of Proposition \ref{prop.Shade.shade-3}.]$\Longrightarrow:$ Assume that
$\operatorname*{Shade}$ is a shade map on $E$. We must show that
$\operatorname*{Shade}$ satisfies Axiom 3.

[\textit{Proof of Axiom 3:} Let $F\in\mathcal{P}\left(  E\right)  $ and $u\in
E\setminus F$. We must prove that we have $\operatorname*{Shade}%
F=\operatorname*{Shade}\left(  F\cup\left\{  u\right\}  \right)  $ or
$u\in\left(  \operatorname*{Shade}F\right)  \cap\operatorname*{Shade}\left(
F\cup\left\{  u\right\}  \right)  $.

Indeed, assume the contrary. Thus, we have neither $\operatorname*{Shade}%
F=\operatorname*{Shade}\left(  F\cup\left\{  u\right\}  \right)  $ nor
$u\in\left(  \operatorname*{Shade}F\right)  \cap\operatorname*{Shade}\left(
F\cup\left\{  u\right\}  \right)  $. In other words, we have
$\operatorname*{Shade}F\neq\operatorname*{Shade}\left(  F\cup\left\{
u\right\}  \right)  $ and $u\notin\left(  \operatorname*{Shade}F\right)
\cap\operatorname*{Shade}\left(  F\cup\left\{  u\right\}  \right)  $.

We know that $\operatorname*{Shade}$ is a shade map (by assumption), and thus
satisfies Axioms 1 and 2 in Definition \ref{def.Shade.shademap}. Hence, if we
had $u\in E\setminus\operatorname*{Shade}F$, then we would have
$\operatorname*{Shade}\left(  F\cup\left\{  u\right\}  \right)
=\operatorname*{Shade}F$ (by Axiom 1), which would contradict
$\operatorname*{Shade}F\neq\operatorname*{Shade}\left(  F\cup\left\{
u\right\}  \right)  $. Thus, we cannot have $u\in E\setminus
\operatorname*{Shade}F$. Hence, we have $u\in E\setminus\left(  E\setminus
\operatorname*{Shade}F\right)  =\operatorname*{Shade}F$ (since
$\operatorname*{Shade}F\subseteq E$).

Therefore, if we had $u\in\operatorname*{Shade}\left(  F\cup\left\{
u\right\}  \right)  $, then we would have $u\in\left(  \operatorname*{Shade}%
F\right)  \cap\operatorname*{Shade}\left(  F\cup\left\{  u\right\}  \right)
$, which would contradict $u\notin\left(  \operatorname*{Shade}F\right)
\cap\operatorname*{Shade}\left(  F\cup\left\{  u\right\}  \right)  $. Thus, we
cannot have $u\in\operatorname*{Shade}\left(  F\cup\left\{  u\right\}
\right)  $. Hence, we have $u\notin\operatorname*{Shade}\left(  F\cup\left\{
u\right\}  \right)  $. Therefore, $u\in E\setminus\operatorname*{Shade}\left(
F\cup\left\{  u\right\}  \right)  $. Hence, Axiom 2 (applied to $F\cup\left\{
u\right\}  $ instead of $F$) yields $\operatorname*{Shade}\left(  \left(
F\cup\left\{  u\right\}  \right)  \setminus\left\{  u\right\}  \right)
=\operatorname*{Shade}\left(  F\cup\left\{  u\right\}  \right)  $. However,
$u\notin F$ (since $u\in E\setminus F$) and thus $\left(  F\cup\left\{
u\right\}  \right)  \setminus\left\{  u\right\}  =F$. Hence, the equality
$\operatorname*{Shade}\left(  \left(  F\cup\left\{  u\right\}  \right)
\setminus\left\{  u\right\}  \right)  =\operatorname*{Shade}\left(
F\cup\left\{  u\right\}  \right)  $ rewrites as $\operatorname*{Shade}%
F=\operatorname*{Shade}\left(  F\cup\left\{  u\right\}  \right)  $. This
contradicts $\operatorname*{Shade}F\neq\operatorname*{Shade}\left(
F\cup\left\{  u\right\}  \right)  $.

This contradiction shows that our assumption was wrong. Hence, we have proved
that we have $\operatorname*{Shade}F=\operatorname*{Shade}\left(
F\cup\left\{  u\right\}  \right)  $ or $u\in\left(  \operatorname*{Shade}%
F\right)  \cap\operatorname*{Shade}\left(  F\cup\left\{  u\right\}  \right)
$. This proves Axiom 3.]

We thus have shown that $\operatorname*{Shade}$ satisfies Axiom 3. This proves
the \textquotedblleft$\Longrightarrow$\textquotedblright\ direction of
Proposition \ref{prop.Shade.shade-3}.

$\Longleftarrow:$ Assume that $\operatorname*{Shade}$ satisfies Axiom 3. We
must prove that $\operatorname*{Shade}$ is a shade map on $E$. To this
purpose, we shall show that $\operatorname*{Shade}$ satisfies Axioms 1 and 2
in Definition \ref{def.Shade.shademap}.

[\textit{Proof of Axiom 1:} Let $F\in\mathcal{P}\left(  E\right)  $ and $u\in
E\setminus\operatorname*{Shade}F$. We must show that $\operatorname*{Shade}%
\left(  F\cup\left\{  u\right\}  \right)  =\operatorname*{Shade}F$.

Indeed, assume the contrary. Thus, $\operatorname*{Shade}\left(  F\cup\left\{
u\right\}  \right)  \neq\operatorname*{Shade}F$, so that $F\cup\left\{
u\right\}  \neq F$. Consequently, $u\notin F$ (because otherwise, we would
have $u\in F$ and thus $F\cup\left\{  u\right\}  =F$, contradicting
$F\cup\left\{  u\right\}  \neq F$). Hence, $u\in E\setminus F$.

We assumed that Axiom 3 holds. Thus, Axiom 3 shows that we have
$\operatorname*{Shade}F=\operatorname*{Shade}\left(  F\cup\left\{  u\right\}
\right)  $ or $u\in\left(  \operatorname*{Shade}F\right)  \cap
\operatorname*{Shade}\left(  F\cup\left\{  u\right\}  \right)  $. Hence,
$u\in\left(  \operatorname*{Shade}F\right)  \cap\operatorname*{Shade}\left(
F\cup\left\{  u\right\}  \right)  $ (since $\operatorname*{Shade}%
F=\operatorname*{Shade}\left(  F\cup\left\{  u\right\}  \right)  $ would
contradict $\operatorname*{Shade}\left(  F\cup\left\{  u\right\}  \right)
\neq\operatorname*{Shade}F$). Thus, $u\in\left(  \operatorname*{Shade}%
F\right)  \cap\operatorname*{Shade}\left(  F\cup\left\{  u\right\}  \right)
\subseteq\operatorname*{Shade}F$. However, $u\in E\setminus
\operatorname*{Shade}F$ entails $u\notin\operatorname*{Shade}F$. This
contradicts $u\in\operatorname*{Shade}F$. This contradiction shows that our
assumption was false. Hence, we have proved $\operatorname*{Shade}\left(
F\cup\left\{  u\right\}  \right)  =\operatorname*{Shade}F$; thus Axiom 1 is proven.]

[\textit{Proof of Axiom 2:} Let $F\in\mathcal{P}\left(  E\right)  $ and $u\in
E\setminus\operatorname*{Shade}F$. We must show that $\operatorname*{Shade}%
\left(  F\setminus\left\{  u\right\}  \right)  =\operatorname*{Shade}F$.

Indeed, assume the contrary. Thus, $\operatorname*{Shade}\left(
F\setminus\left\{  u\right\}  \right)  \neq\operatorname*{Shade}F$, so that
$F\setminus\left\{  u\right\}  \neq F$. Consequently, $u\in F$ (because
otherwise, we would have $u\notin F$ and thus $F\setminus\left\{  u\right\}
=F$, contradicting $F\setminus\left\{  u\right\}  \neq F$). Hence, $\left(
F\setminus\left\{  u\right\}  \right)  \cup\left\{  u\right\}  =F$.

We have $u\in E\setminus\left(  F\setminus\left\{  u\right\}  \right)  $
(since $u\in E$ and $u\notin F\setminus\left\{  u\right\}  $). We assumed that
Axiom 3 holds. Thus, Axiom 3 (applied to $F\setminus\left\{  u\right\}  $
instead of $F$) shows that we have $\operatorname*{Shade}\left(
F\setminus\left\{  u\right\}  \right)  =\operatorname*{Shade}\left(  \left(
F\setminus\left\{  u\right\}  \right)  \cup\left\{  u\right\}  \right)  $ or
$u\in\left(  \operatorname*{Shade}\left(  F\setminus\left\{  u\right\}
\right)  \right)  \cap\operatorname*{Shade}\left(  \left(  F\setminus\left\{
u\right\}  \right)  \cup\left\{  u\right\}  \right)  $. In view of $\left(
F\setminus\left\{  u\right\}  \right)  \cup\left\{  u\right\}  =F$, this
rewrites as follows: We have $\operatorname*{Shade}\left(  F\setminus\left\{
u\right\}  \right)  =\operatorname*{Shade}F$ or $u\in\left(
\operatorname*{Shade}\left(  F\setminus\left\{  u\right\}  \right)  \right)
\cap\operatorname*{Shade}F$. Hence, $u\in\left(  \operatorname*{Shade}\left(
F\setminus\left\{  u\right\}  \right)  \right)  \cap\operatorname*{Shade}F$
(since we assumed that we don't have $\operatorname*{Shade}\left(
F\setminus\left\{  u\right\}  \right)  =\operatorname*{Shade}F$). Thus,
$u\in\left(  \operatorname*{Shade}\left(  F\setminus\left\{  u\right\}
\right)  \right)  \cap\operatorname*{Shade}F\subseteq\operatorname*{Shade}F$.
However, $u\in E\setminus\operatorname*{Shade}F$ entails $u\notin%
\operatorname*{Shade}F$. This contradicts $u\in\operatorname*{Shade}F$. This
contradiction shows that our assumption was false. Hence, we have proved
$\operatorname*{Shade}\left(  F\setminus\left\{  u\right\}  \right)
=\operatorname*{Shade}F$; thus Axiom 2 is proven.]

We have now shown that the map $\operatorname*{Shade}$ satisfies both Axioms 1
and 2 in Definition \ref{def.Shade.shademap}. Hence, $\operatorname*{Shade}$
is a shade map. This proves the \textquotedblleft$\Longleftarrow
$\textquotedblright\ direction of Proposition \ref{prop.Shade.shade-3}.

Thus, the proof of Proposition \ref{prop.Shade.shade-3} is complete.
\end{proof}
\end{verlong}

We will soon see some examples. First, let us introduce two more basic
concepts that will help clarify these examples:

\begin{definition}
\label{def.inclusion-pres}Let $E$ be a set. Let $\operatorname*{Shade}%
:\mathcal{P}\left(  E\right)  \rightarrow\mathcal{P}\left(  E\right)  $ be any
map (not necessarily a shade map).

\textbf{(a)} We say that $\operatorname*{Shade}$ is
\emph{inclusion-preserving} if it satisfies the following property: If $A$ and
$B$ are two subsets of $E$ such that $A\subseteq B$, then
$\operatorname*{Shade}A\subseteq\operatorname*{Shade}B$.

\textbf{(b)} We say that $\operatorname*{Shade}$ is \emph{inclusion-reversing}
if it satisfies the following property: If $A$ and $B$ are two subsets of $E$
such that $A\subseteq B$, then $\operatorname*{Shade}B\subseteq
\operatorname*{Shade}A$.
\end{definition}

For instance, the map $\operatorname*{Shade}$ from Definition \ref{def.Shade}
is inclusion-preserving (because of Lemma \ref{lem.Shade-monoton}) and is a
shade map (by Lemma \ref{lem.Shade-tog}). The same holds for the analogue of
the map $\operatorname*{Shade}$ that uses vertex-infection instead of
infection. This does not mean that any shade map is inclusion-preserving.
Indeed, we shall soon see some inclusion-reversing shade maps, and it is not
hard to construct shade maps that are neither inclusion-preserving nor inclusion-reversing.

Let us observe that there is a simple bijection between inclusion-preserving
and inclusion-reversing maps, and this bijection preserves shadeness:

\begin{proposition}
\label{prop.Shade.dual}Let $E$ be a set. Let $\operatorname*{Shade}%
:\mathcal{P}\left(  E\right)  \rightarrow\mathcal{P}\left(  E\right)  $ be any
map (not necessarily a shade map). Let $\operatorname*{Shade}\nolimits^{\prime
}:\mathcal{P}\left(  E\right)  \rightarrow\mathcal{P}\left(  E\right)  $ be
the map that sends each $F\in\mathcal{P}\left(  E\right)  $ to
$\operatorname*{Shade}\left(  E\setminus F\right)  \in\mathcal{P}\left(
E\right)  $. Then:

\textbf{(a)} The map $\operatorname*{Shade}$ is inclusion-preserving if and
only if the map $\operatorname*{Shade}\nolimits^{\prime}$ is inclusion-reversing.

\textbf{(b)} The map $\operatorname*{Shade}$ is a shade map if and only if the
map $\operatorname*{Shade}\nolimits^{\prime}$ is a shade map.
\end{proposition}

\begin{proof}
This is easy and left to the reader. (For part \textbf{(b)}, observe that
Axioms 1 and 2 for the map $\operatorname*{Shade}$ translate into Axioms 2 and
1 for the map $\operatorname*{Shade}\nolimits^{\prime}$.)
\end{proof}

\begin{definition}
Let $E$, $\operatorname*{Shade}$ and $\operatorname*{Shade}\nolimits^{\prime}$
be as in Proposition \ref{prop.Shade.dual}. We then say that the map
$\operatorname*{Shade}\nolimits^{\prime}$ is \emph{dual} to
$\operatorname*{Shade}$.
\end{definition}

\subsection{Some examples of shade maps}

As we already mentioned, Lemma \ref{lem.Shade-tog} and its analogue for
vertex-infection provide two examples of inclusion-preserving shade maps
$\operatorname*{Shade}$. An example of an inclusion-reversing shade map comes
from the theory of posets:

\begin{example}
\label{exa.shade.low}Let $E$ be a poset. For any $F\subseteq E$, we define%
\[
F_{\downarrow}=\left\{  e\in E\ \mid\ \text{there exists an }f\in F\text{ with
}e<f\right\}
\]
and%
\[
\operatorname*{Shade}F=E\setminus F_{\downarrow}.
\]
Then, this map $\operatorname*{Shade}:\mathcal{P}\left(  E\right)
\rightarrow\mathcal{P}\left(  E\right)  $ is an inclusion-reversing shade map.
\end{example}

\begin{vershort}
For the straightforward proof of Example \ref{exa.shade.low}, see the detailed
version \cite{verlong} of this paper.
\end{vershort}

\begin{verlong}
\begin{proof}
[Proof of Example \ref{exa.shade.low} (sketched).]We shall prove the following
three statements:

\begin{statement}
\textit{Statement 0:} If $A$ and $B$ are two subsets of $E$ such that
$A\subseteq B$, then $\operatorname*{Shade}B\subseteq\operatorname*{Shade}A$.
\end{statement}

\begin{statement}
\textit{Statement 1:} If $F\in\mathcal{P}\left(  E\right)  $ and $u\in
E\setminus\operatorname*{Shade}F$, then $\operatorname*{Shade}\left(
F\cup\left\{  u\right\}  \right)  =\operatorname*{Shade}F$.
\end{statement}

\begin{statement}
\textit{Statement 2:} If $F\in\mathcal{P}\left(  E\right)  $ and $u\in
E\setminus\operatorname*{Shade}F$, then $\operatorname*{Shade}\left(
F\setminus\left\{  u\right\}  \right)  =\operatorname*{Shade}F$.
\end{statement}

[\textit{Proof of Statement 0:} Let $A$ and $B$ be two subsets of $E$ such
that $A\subseteq B$. We must show that $\operatorname*{Shade}B\subseteq
\operatorname*{Shade}A$.

Indeed, $A\subseteq B$ entails $A_{\downarrow}\subseteq B_{\downarrow}$ (by
the definitions of $A_{\downarrow}$ and $B_{\downarrow}$). Hence, $E\setminus
A_{\downarrow}\supseteq E\setminus B_{\downarrow}$, so that $E\setminus
B_{\downarrow}\subseteq E\setminus A_{\downarrow}$. But the definition of
$\operatorname*{Shade}$ yields $\operatorname*{Shade}A=E\setminus
A_{\downarrow}$ and $\operatorname*{Shade}B=E\setminus B_{\downarrow}$. Hence,
$\operatorname*{Shade}B=E\setminus B_{\downarrow}\subseteq E\setminus
A_{\downarrow}=\operatorname*{Shade}A$. This proves Statement 0.]

[\textit{Proof of Statement 1:} Let $F\in\mathcal{P}\left(  E\right)  $ and
$u\in E\setminus\operatorname*{Shade}F$. We must prove that
$\operatorname*{Shade}\left(  F\cup\left\{  u\right\}  \right)
=\operatorname*{Shade}F$.

We have $u\in E\setminus\operatorname*{Shade}F=F_{\downarrow}$ (since
$\operatorname*{Shade}F=E\setminus F_{\downarrow}$). In other words, there
exists a $g\in F$ with $u<g$ (by the definition of $F_{\downarrow}$). Consider
this $g$.

We must prove that $\operatorname*{Shade}\left(  F\cup\left\{  u\right\}
\right)  =\operatorname*{Shade}F$. In other words, we must prove that $\left(
F\cup\left\{  u\right\}  \right)  _{\downarrow}=F_{\downarrow}$ (since
$\operatorname*{Shade}\left(  F\cup\left\{  u\right\}  \right)  =E\setminus
\left(  F\cup\left\{  u\right\}  \right)  _{\downarrow}$ and
$\operatorname*{Shade}F=E\setminus F_{\downarrow}$). Since $F_{\downarrow
}\subseteq\left(  F\cup\left\{  u\right\}  \right)  _{\downarrow}$ is rather
obvious (by the definitions of $F_{\downarrow}$ and $\left(  F\cup\left\{
u\right\}  \right)  _{\downarrow}$), we only need to prove that $\left(
F\cup\left\{  u\right\}  \right)  _{\downarrow}\subseteq F_{\downarrow}$.

So let $v\in\left(  F\cup\left\{  u\right\}  \right)  _{\downarrow}$. We must
show that $v\in F_{\downarrow}$.

We have $v\in\left(  F\cup\left\{  u\right\}  \right)  _{\downarrow}$. In
other words, there exists an $f\in F\cup\left\{  u\right\}  $ with $v<f$ (by
the definition of $\left(  F\cup\left\{  u\right\}  \right)  _{\downarrow}$).
Consider this $f$. If $f\in F$, then we immediately obtain $v\in
F_{\downarrow}$ (since $v<f$ and $f\in F$), and thus our proof is complete.
Hence, we WLOG assume that $f\notin F$. Hence, $f=u$ (since $f\in
F\cup\left\{  u\right\}  $ but $f\notin F$) and thus $v<f=u<g$. This entails
$v\in F_{\downarrow}$ (by the definition of $F_{\downarrow}$, since $g\in F$).

Forget that we fixed $v$. We thus have shown that $v\in F_{\downarrow}$ for
each $v\in\left(  F\cup\left\{  u\right\}  \right)  _{\downarrow}$. In other
words, $\left(  F\cup\left\{  u\right\}  \right)  _{\downarrow}\subseteq
F_{\downarrow}$. As we explained above, this completes the proof of Statement 1.]

[\textit{Proof of Statement 2:} Let $F\in\mathcal{P}\left(  E\right)  $ and
$u\in E\setminus\operatorname*{Shade}F$. We must prove that
$\operatorname*{Shade}\left(  F\setminus\left\{  u\right\}  \right)
=\operatorname*{Shade}F$.

We have $u\in E\setminus\operatorname*{Shade}F=F_{\downarrow}$ (since
$\operatorname*{Shade}F=E\setminus F_{\downarrow}$). In other words, there
exists a $g\in F$ with $u<g$ (by the definition of $F_{\downarrow}$). Consider
this $g$. Note that $u<g$ entails $g\neq u$, so that $g\in F\setminus\left\{
u\right\}  $ (since $g\in F$).

We must prove that $\operatorname*{Shade}\left(  F\setminus\left\{  u\right\}
\right)  =\operatorname*{Shade}F$. In other words, we must prove that $\left(
F\setminus\left\{  u\right\}  \right)  _{\downarrow}=F_{\downarrow}$ (since
$\operatorname*{Shade}\left(  F\setminus\left\{  u\right\}  \right)
=E\setminus\left(  F\setminus\left\{  u\right\}  \right)  _{\downarrow}$ and
$\operatorname*{Shade}F=E\setminus F_{\downarrow}$). Since $\left(
F\setminus\left\{  u\right\}  \right)  _{\downarrow}\subseteq F_{\downarrow}$
is rather obvious (by the definitions of $F_{\downarrow}$ and $\left(
F\setminus\left\{  u\right\}  \right)  _{\downarrow}$), we only need to prove
that $F_{\downarrow}\subseteq\left(  F\setminus\left\{  u\right\}  \right)
_{\downarrow}$.

So let $v\in F_{\downarrow}$. We must show that $v\in\left(  F\setminus
\left\{  u\right\}  \right)  _{\downarrow}$.

We have $v\in F_{\downarrow}$. In other words, there exists an $f\in F$ with
$v<f$ (by the definition of $F_{\downarrow}$). Consider this $f$. If $f\in
F\setminus\left\{  u\right\}  $, then we immediately obtain $v\in\left(
F\setminus\left\{  u\right\}  \right)  _{\downarrow}$ (since $v<f$ and $f\in
F\setminus\left\{  u\right\}  $), and thus our proof is complete. Hence, we
WLOG assume that $f\notin F\setminus\left\{  u\right\}  $. Hence, $f=u$ (since
$f\in F$ but $f\notin F\setminus\left\{  u\right\}  $) and thus $v<f=u<g$.
This entails $v\in\left(  F\setminus\left\{  u\right\}  \right)  _{\downarrow
}$ (by the definition of $\left(  F\setminus\left\{  u\right\}  \right)
_{\downarrow}$, since $g\in F\setminus\left\{  u\right\}  $).

Forget that we fixed $v$. We thus have shown that $v\in\left(  F\setminus
\left\{  u\right\}  \right)  _{\downarrow}$ for each $v\in F_{\downarrow}$. In
other words, $F_{\downarrow}\subseteq\left(  F\setminus\left\{  u\right\}
\right)  _{\downarrow}$. As we explained above, this completes the proof of
Statement 2.]

Now, we have proved Statements 1 and 2. Thus, the map $\operatorname*{Shade}%
:\mathcal{P}\left(  E\right)  \rightarrow\mathcal{P}\left(  E\right)  $
satisfies the two axioms in Definition \ref{def.Shade.shademap}. In other
words, this map is a shade map. Moreover, this map is inclusion-reversing (by
Statement 0). This proves Example \ref{exa.shade.low}.
\end{proof}
\end{verlong}

Another example of a shade map comes from discrete geometry:

\begin{example}
\label{exa.shade.convex}Let $A$ be an affine space over $\mathbb{R}$. If $S$
is a finite subset of $A$, then a \emph{nontrivial convex combination} of $S$
will mean a point of the form $\sum_{s\in S}\lambda_{s}s\in A$, where the
coefficients $\lambda_{s}$ are nonnegative reals smaller than $1$ and
satisfying $\sum_{s\in S}\lambda_{s}=1$.

Fix a finite subset $E$ of $A$. For any $F\subseteq E$, we define%
\[
\operatorname*{Shade}F=\left\{  e\in E\ \mid\ e\text{ is \textbf{not} a
nontrivial convex combination of }F\right\}  .
\]
Then, this map $\operatorname*{Shade}:\mathcal{P}\left(  E\right)
\rightarrow\mathcal{P}\left(  E\right)  $ is an inclusion-reversing shade map.
\end{example}

\begin{vershort}
For a (not very difficult) proof of Example \ref{exa.shade.convex}, see the
detailed version \cite{verlong} of this paper.
\end{vershort}

\begin{verlong}
\begin{proof}
[Proof of Example \ref{exa.shade.convex} (sketched).]We shall prove the
following three statements:

\begin{statement}
\textit{Statement 0:} If $X$ and $Y$ are two subsets of $E$ such that
$X\subseteq Y$, then $\operatorname*{Shade}Y\subseteq\operatorname*{Shade}X$.
\end{statement}

\begin{statement}
\textit{Statement 1:} If $F\in\mathcal{P}\left(  E\right)  $ and $u\in
E\setminus\operatorname*{Shade}F$, then $\operatorname*{Shade}\left(
F\cup\left\{  u\right\}  \right)  =\operatorname*{Shade}F$.
\end{statement}

\begin{statement}
\textit{Statement 2:} If $F\in\mathcal{P}\left(  E\right)  $ and $u\in
E\setminus\operatorname*{Shade}F$, then $\operatorname*{Shade}\left(
F\setminus\left\{  u\right\}  \right)  =\operatorname*{Shade}F$.
\end{statement}

[\textit{Proof of Statement 0:} Let $X$ and $Y$ be two subsets of $E$ such
that $X\subseteq Y$. We must prove that $\operatorname*{Shade}Y\subseteq
\operatorname*{Shade}X$.

Indeed, we have $X\subseteq Y$. Thus, any nontrivial convex combination of $X$
is a nontrivial convex combination of $Y$. Therefore, any $e\in E$ that is
\textbf{not} a nontrivial convex combination of $Y$ also cannot be a
nontrivial convex combination of $X$. In other words, $\operatorname*{Shade}%
Y\subseteq\operatorname*{Shade}X$ (by the definition of $\operatorname*{Shade}%
$). This proves Statement 0.]

[\textit{Proof of Statement 1:} Let $F\in\mathcal{P}\left(  E\right)  $ and
$u\in E\setminus\operatorname*{Shade}F$. We must prove that
$\operatorname*{Shade}\left(  F\cup\left\{  u\right\}  \right)
=\operatorname*{Shade}F$. If $u\in F$, then this is obvious (since
$F\cup\left\{  u\right\}  =F$ in this case); thus, we WLOG assume that
$u\notin F$.

We must prove that $\operatorname*{Shade}\left(  F\cup\left\{  u\right\}
\right)  =\operatorname*{Shade}F$. We have $\operatorname*{Shade}\left(
F\cup\left\{  u\right\}  \right)  \subseteq\operatorname*{Shade}F$ (by
Statement 0, since $F\subseteq F\cup\left\{  u\right\}  $). Thus, it suffices
to show that $\operatorname*{Shade}F\subseteq\operatorname*{Shade}\left(
F\cup\left\{  u\right\}  \right)  $. So let $p\in\operatorname*{Shade}F$. We
must then show that $p\in\operatorname*{Shade}\left(  F\cup\left\{  u\right\}
\right)  $. In other words, we must prove that $p$ is \textbf{not} a
nontrivial convex combination of $F\cup\left\{  u\right\}  $.

Assume the contrary. Thus, $p$ is a nontrivial convex combination of
$F\cup\left\{  u\right\}  $. In other words, $p=\sum_{s\in F\cup\left\{
u\right\}  }\mu_{s}s$ for some nonnegative reals $\mu_{s}$ smaller than $1$
and satisfying $\sum_{s\in F\cup\left\{  u\right\}  }\mu_{s}=1$. Consider
these reals $\mu_{s}$.

We have $u\in E\setminus\operatorname*{Shade}F$. Hence, $u\in E$ and
$u\notin\operatorname*{Shade}F$. From $u\notin\operatorname*{Shade}F$, we see
that $u$ is a nontrivial convex combination of $F$. In other words,
$u=\sum_{f\in F}\lambda_{f}f$ for some nonnegative reals $\lambda_{f}$ smaller
than $1$ and satisfying $\sum_{f\in F}\lambda_{f}=1$. Consider these reals
$\lambda_{f}$.

Now,%
\begin{align}
p  &  =\sum_{s\in F\cup\left\{  u\right\}  }\mu_{s}s=\underbrace{\sum_{s\in
F}\mu_{s}s}_{=\sum_{f\in F}\mu_{f}f}+\mu_{u}\underbrace{u}_{=\sum_{f\in
F}\lambda_{f}f}\nonumber\\
&  \ \ \ \ \ \ \ \ \ \ \left(
\begin{array}
[c]{c}%
\text{here, we have split off the addend for }s=u\\
\text{from the sum, since }u\notin F
\end{array}
\right) \nonumber\\
&  =\sum_{f\in F}\mu_{f}f+\mu_{u}\sum_{f\in F}\lambda_{f}f=\sum_{f\in
F}\left(  \mu_{f}+\mu_{u}\lambda_{f}\right)  f.
\label{pf.exa.shade.convex.s1.pf.1}%
\end{align}
Clearly, the coefficients $\mu_{f}+\mu_{u}\lambda_{f}$ on the right hand side
of this equality are nonnegative reals. Moreover, they are easily seen to be
smaller than $1$\ \ \ \ \footnote{\textit{Proof.} Assume the contrary. Thus,
$\mu_{f}+\mu_{u}\lambda_{f}\geq1$ for some $f\in F$. Consider this $f$. Note
that $\mu_{f}<1$ (since all coefficients $\mu_{s}$ are smaller than $1$) and
$\lambda_{f}<1$ (for similar reasons). Now, $\mu_{f}+\mu_{u}\lambda_{f}%
\geq1>\mu_{f}$ (since $\mu_{f}<1$), so that $\mu_{u}\lambda_{f}>0$. This
entails $\mu_{u}\lambda_{f}\neq0$ and thus $\mu_{u}\neq0$, so that $\mu_{u}>0$
(since $\mu_{u}$ is nonnegative). Now, $\mu_{u}\underbrace{\lambda_{f}}%
_{<1}<\mu_{u}$ (since $\mu_{u}>0$).
\par
However, $f\neq u$ (since $f\in F$ and $u\notin F$), so that $\mu_{f}+\mu_{u}$
is a subsum of the sum $\sum_{s\in F\cup\left\{  u\right\}  }\mu_{s}$. Since
all the $\mu_{s}$ are nonnegative, this subsum thus must be $\leq\sum_{s\in
F\cup\left\{  u\right\}  }\mu_{s}$. In other words, $\mu_{f}+\mu_{u}\leq
\sum_{s\in F\cup\left\{  u\right\}  }\mu_{s}$. Hence, $\mu_{f}+\underbrace{\mu
_{u}\lambda_{f}}_{<\mu_{u}}<\mu_{f}+\mu_{u}\leq\sum_{s\in F\cup\left\{
u\right\}  }\mu_{s}=1$. This contradicts $\mu_{f}+\mu_{u}\lambda_{f}\geq1$.
This contradiction shows that our assumption was wrong, qed.}. Finally, they
satisfy
\begin{align*}
\sum_{f\in F}\left(  \mu_{f}+\mu_{u}\lambda_{f}\right)   &  =\sum_{f\in F}%
\mu_{f}+\mu_{u}\underbrace{\sum_{f\in F}\lambda_{f}}_{=1}=\sum_{f\in F}\mu
_{f}+\mu_{u}=\sum_{f\in F\cup\left\{  u\right\}  }\mu_{f}%
\ \ \ \ \ \ \ \ \ \ \left(  \text{since }u\notin F\right) \\
&  =\sum_{s\in F\cup\left\{  u\right\}  }\mu_{s}=1.
\end{align*}
Hence, (\ref{pf.exa.shade.convex.s1.pf.1}) shows that $p$ is a nontrivial
convex combination of $F$. In other words, $p\notin\operatorname*{Shade}F$.
This contradicts $p\in\operatorname*{Shade}F$. This contradiction shows that
our assumption was false. This completes the proof of Statement 1.]

[\textit{Proof of Statement 2:} Let $F\in\mathcal{P}\left(  E\right)  $ and
$u\in E\setminus\operatorname*{Shade}F$. We must prove that
$\operatorname*{Shade}\left(  F\setminus\left\{  u\right\}  \right)
=\operatorname*{Shade}F$. If $u\notin F$, then this is obvious (since
$F\setminus\left\{  u\right\}  =F$ in this case); thus, we WLOG assume that
$u\in F$.

We have $u\in E\setminus\operatorname*{Shade}F$. Hence, $u\in E$ and
$u\notin\operatorname*{Shade}F$. From $u\notin\operatorname*{Shade}F$, we see
that $u$ is a nontrivial convex combination of $F$. In other words,
$u=\sum_{f\in F}\lambda_{f}f$ for some nonnegative reals $\lambda_{f}$ smaller
than $1$ and satisfying $\sum_{f\in F}\lambda_{f}=1$. Consider these reals
$\lambda_{f}$.

We have $\sum_{g\in F}\lambda_{g}=\sum_{f\in F}\lambda_{f}=1$.

We have $\lambda_{u}<1$ (since all reals $\lambda_{f}$ are smaller than $1$).
Hence, $1-\lambda_{u}>0$, so that $\dfrac{1}{1-\lambda_{u}}$ is a well-defined
positive real.

We have%
\[
u=\sum_{f\in F}\lambda_{f}f=\sum_{f\in F\setminus\left\{  u\right\}  }%
\lambda_{f}f+\lambda_{u}u
\]
(here, we have split off the addend for $f=u$ from the sum, since $u\in F$).
In other words, $u-\lambda_{u}u=\sum_{f\in F\setminus\left\{  u\right\}
}\lambda_{f}f$. In other words, $\left(  1-\lambda_{u}\right)  u=\sum_{f\in
F\setminus\left\{  u\right\}  }\lambda_{f}f$. Since $1-\lambda_{u}>0$, we can
divide this equality by $1-\lambda_{u}$, and obtain%
\begin{equation}
u=\dfrac{1}{1-\lambda_{u}}\sum_{f\in F\setminus\left\{  u\right\}  }%
\lambda_{f}f=\sum_{f\in F\setminus\left\{  u\right\}  }\dfrac{\lambda_{f}%
}{1-\lambda_{u}}f. \label{pf.exa.shade.convex.s2.pf.1}%
\end{equation}
Clearly, the coefficients $\dfrac{\lambda_{f}}{1-\lambda_{u}}$ on the right
hand side of this equality are nonnegative reals (since $1-\lambda_{u}>0$).
Moreover, they are easily seen to be smaller than $1$%
\ \ \ \ \footnote{\textit{Proof.} Assume the contrary. Thus, $\dfrac
{\lambda_{f}}{1-\lambda_{u}}\geq1$ for some $f\in F\setminus\left\{
u\right\}  $. Consider this $f$. From $\dfrac{\lambda_{f}}{1-\lambda_{u}}%
\geq1$, we obtain $\lambda_{f}\geq1-\lambda_{u}$ (since $1-\lambda_{u}>0$),
thus $\lambda_{f}+\lambda_{u}\geq1$.
\par
However, $f\neq u$ (since $f\in F\setminus\left\{  u\right\}  $), so that
$\lambda_{f}+\lambda_{u}$ is a subsum of the sum $\sum_{g\in F}\lambda_{g}$.
Since all the $\lambda_{g}$ are nonnegative, this subsum thus must be
$\leq\sum_{g\in F}\lambda_{g}$. In other words, $\lambda_{f}+\lambda_{u}%
\leq\sum_{g\in F}\lambda_{g}$. Hence, $\lambda_{f}+\lambda_{u}\leq\sum_{g\in
F}\lambda_{g}=1$. Combining this with $\lambda_{f}+\lambda_{u}\geq1$, we
obtain $\lambda_{f}+\lambda_{u}=1$. Thus, $\lambda_{f}=1-\lambda_{u}$.
\par
Now, we observe that $\lambda_{f}+\lambda_{u}\geq1>\lambda_{u}$ (since
$\lambda_{u}<1$), so that $\lambda_{f}>0$. Moreover, if $h\in F$ satisfies
$h\neq f$ and $h\neq u$, then $\lambda_{f}+\lambda_{u}+\lambda_{h}$ is a
subsum of the sum $\sum_{g\in F}\lambda_{g}$ and thus satisfies $\lambda
_{f}+\lambda_{u}+\lambda_{h}\leq\sum_{g\in F}\lambda_{g}$ (since all the
$\lambda_{g}$ are nonnegative) and therefore $\lambda_{f}+\lambda_{u}%
+\lambda_{h}\leq\sum_{g\in F}\lambda_{g}=1=\lambda_{f}+\lambda_{u}$, so that
$\lambda_{h}\leq0$ and therefore $\lambda_{h}=0$ (since $\lambda_{h}$ is a
nonnegative real). Thus, we have shown that all elements $h\in F$ except for
(possibly) $f$ and $u$ satisfy $\lambda_{h}=0$. Thus,
\[
\sum_{h\in F}\lambda_{h}h=\lambda_{f}f+\lambda_{u}u\ \ \ \ \ \ \ \ \ \ \left(
\text{since }f\text{ and }u\text{ are distinct elements of }F\right)  .
\]
\par
However, recall that $u=\sum_{h\in F}\lambda_{h}h$. (This is just the old
equation $u=\sum_{f\in F}\lambda_{f}f$, except that we had to rename the
summation index $f$ as $h$ since $f$ means something else now.) Thus,%
\[
u=\sum_{h\in F}\lambda_{h}h=\lambda_{f}f+\lambda_{u}u,
\]
so that $u-\lambda_{u}u=\lambda_{f}f$. In view of $u-\lambda_{u}u=\left(
1-\lambda_{u}\right)  u$ and $\lambda_{f}=1-\lambda_{u}$, this rewrites as
$\left(  1-\lambda_{u}\right)  u=\left(  1-\lambda_{u}\right)  f$. We can
cancel $1-\lambda_{u}$ from this equality (since $1-\lambda_{u}>0$), and thus
obtain $u=f$. This contradicts $f\neq u$. This contradiction shows that our
assumption was wrong, qed.}. Finally, they satisfy
\begin{align*}
\sum_{f\in F\setminus\left\{  u\right\}  }\dfrac{\lambda_{f}}{1-\lambda_{u}}
&  =\dfrac{1}{1-\lambda_{u}}\underbrace{\sum_{f\in F\setminus\left\{
u\right\}  }\lambda_{f}}_{\substack{=\sum_{f\in F}\lambda_{f}-\lambda
_{u}\\\text{(since }u\in F\text{)}}}=\dfrac{1}{1-\lambda_{u}}\left(
\underbrace{\sum_{f\in F}\lambda_{f}}_{=1}-\lambda_{u}\right) \\
&  =\dfrac{1}{1-\lambda_{u}}\left(  1-\lambda_{u}\right)  =1.
\end{align*}
Hence, (\ref{pf.exa.shade.convex.s2.pf.1}) shows that $u$ is a nontrivial
convex combination of $F\setminus\left\{  u\right\}  $. In other words,
$u\notin\operatorname*{Shade}\left(  F\setminus\left\{  u\right\}  \right)  $.
Thus, $u\in E\setminus\operatorname*{Shade}\left(  F\setminus\left\{
u\right\}  \right)  $. Hence, Statement 1 (applied to $F\setminus\left\{
u\right\}  $ instead of $F$) yields $\operatorname*{Shade}\left(  \left(
F\setminus\left\{  u\right\}  \right)  \cup\left\{  u\right\}  \right)
=\operatorname*{Shade}\left(  F\setminus\left\{  u\right\}  \right)  $.
Therefore,%
\[
\operatorname*{Shade}\left(  F\setminus\left\{  u\right\}  \right)
=\operatorname*{Shade}\left(  \underbrace{\left(  F\setminus\left\{
u\right\}  \right)  \cup\left\{  u\right\}  }_{\substack{=F\\\text{(since
}u\in F\text{)}}}\right)  =\operatorname*{Shade}F.
\]
This proves Statement 2.]

Now, we have proved Statements 1 and 2. Thus, the map $\operatorname*{Shade}%
:\mathcal{P}\left(  E\right)  \rightarrow\mathcal{P}\left(  E\right)  $
satisfies the two axioms in Definition \ref{def.Shade.shademap}. In other
words, this map is a shade map. Moreover, this map is inclusion-reversing (by
Statement 0). This proves Example \ref{exa.shade.convex}.
\end{proof}
\end{verlong}

As a contrast to Example \ref{exa.shade.convex}, let us mention a
not-quite-example (satisfying only one of the two axioms in
Definition \ref{def.Shade.shademap}):

\begin{example}
\label{exa.shade.cone}Let $V$ be a vector space over $\mathbb{R}$. If $S$ is a
finite subset of $V$, then a \emph{nontrivial conic combination} of $S$ will
mean a vector of the form $\sum_{s\in S}\lambda_{s}s\in V$, where the
coefficients $\lambda_{s}$ are nonnegative reals with the property that at
least two elements $s\in S$ satisfy $\lambda_{s}>0$.

Fix a finite subset $E$ of $V$. For any $F\subseteq E$, we define%
\[
\operatorname*{Shade}F=\left\{  e\in E\ \mid\ e\text{ is \textbf{not} a
nontrivial conic combination of }F\right\}  .
\]
It can be shown that this map $\operatorname*{Shade}:\mathcal{P}\left(
E\right)  \rightarrow\mathcal{P}\left(  E\right)  $ satisfies Axiom 1 in
Definition \ref{def.Shade.shademap}. In general, it does not satisfy Axiom 2.
Thus, it is not a shade map in general.
\end{example}

\subsection{\label{subsect.abstract.amatr}Antimatroids and inclusion-reversing
shade maps}

Examples \ref{exa.shade.low} and \ref{exa.shade.convex} are instances of a
general class of examples: shade maps coming from \emph{antimatroids}. Not
unlike matroids, antimatroids are a combinatorial concept with many equivalent
avatars (see, e.g., \cite[Chapter III]{KoLoSc91}). Here we shall view them
through one of these avatars: that of \emph{antimatroidal quasi-closure
operators} (roughly equivalent to \emph{convex geometries}). We begin by
defining the notions we need:

\begin{definition}
\label{def.antimat}Let $E$ be any set.

\textbf{(a)} A \emph{quasi-closure operator} on $E$ means a map $\tau
:\mathcal{P}\left(  E\right)  \rightarrow\mathcal{P}\left(  E\right)  $ with
the following properties:

\begin{enumerate}
\item We have $A\subseteq\tau\left(  A\right)  $ for any $A\subseteq E$.

\item If $A$ and $B$ are two subsets of $E$ satisfying $A\subseteq B$, then
$\tau\left(  A\right)  \subseteq\tau\left(  B\right)  $.

\item We have $\tau\left(  \tau\left(  A\right)  \right)  =\tau\left(
A\right)  $ for any $A\subseteq E$.
\end{enumerate}

\textbf{(b)} A quasi-closure operator $\tau$ on $E$ is said to be
\emph{antimatroidal} if it has the following additional property:

\begin{enumerate}
\item[4.] If $X$ is a subset of $E$, and if $y$ and $z$ are two distinct
elements of $E\setminus\tau\left(  X\right)  $ satisfying $z\in\tau\left(
X\cup\left\{  y\right\}  \right)  $, then $y\notin\tau\left(  X\cup\left\{
z\right\}  \right)  $.
\end{enumerate}

\textbf{(c)} A \emph{closure operator} on $E$ means a quasi-closure operator
$\tau$ on $E$ that satisfies $\tau\left(  \varnothing\right)  =\varnothing$.

\textbf{(d)} If $\tau$ is an antimatroidal closure operator on $E$, then the
pair $\left(  E,\tau\right)  $ is called a \emph{convex geometry}.
\end{definition}

Here are some examples of antimatroidal quasi-closure operators:

\begin{example}
\label{exa.antimat.low}Let $E$ be a poset. For any $F\subseteq E$, we define%
\[
\tau\left(  F\right)  =\left\{  e\in E\ \mid\ \text{there exists an }f\in
F\text{ with }e\leq f\right\}  .
\]
Then, $\tau$ is an antimatroidal closure operator on $E$. (This example is the
\textquotedblleft downset alignment\textquotedblright\ from \cite[\S 3,
Example II]{EdeJam85}, and is equivalent to the \textquotedblleft poset
antimatroid\textquotedblright\ from \cite[\S III.2.3]{KoLoSc91}.)
\end{example}

\begin{example}
\label{exa.antimat.interval}Let $E$ be a poset. For any $F\subseteq E$, we
define%
\[
\tau\left(  F\right)  =\left\{  e\in E\ \mid\ \text{there exist }f\in F\text{
and }g\in F\text{ with }g\leq e\leq f\right\}  .
\]
Then, $\tau$ is an antimatroidal closure operator on $E$. (This example is the
\textquotedblleft order convex alignment\textquotedblright\ from \cite[\S 3,
Example II]{EdeJam85}, and is the \textquotedblleft double shelling of a
poset\textquotedblright\ example from \cite[\S III.2.4]{KoLoSc91}.)
\end{example}

\begin{example}
\label{exa.antimat.convex}Let $A$ be an affine space over $\mathbb{R}$. If $S$
is a finite subset of $A$, then a \emph{convex combination} of $S$ will mean a
point of the form $\sum_{s\in S}\lambda_{s}s\in A$, where the coefficients
$\lambda_{s}$ are nonnegative reals satisfying $\sum_{s\in S}\lambda_{s}=1$.

Fix a finite subset $E$ of $A$. For any $F\subseteq E$, we define%
\[
\tau\left(  F\right)  =\left\{  e\in E\ \mid\ e\text{ is a convex combination
of }F\right\}  .
\]
Then, $\tau$ is an antimatroidal closure operator on $E$. (This example is
\cite[\S 3, Example I]{EdeJam85}; it gave the name \textquotedblleft convex
geometry\textquotedblright\ to the notion defined in Definition
\ref{def.antimat} \textbf{(d)}.)
\end{example}

\begin{example}
\label{exa.antimat.linesearch}Let $\Gamma$ be any graph with edge set $E$. Fix
a vertex $v$ of $\Gamma$. We say that a subset $F\subseteq E$ \emph{blocks} an
edge $e\in E$ if each path of $\Gamma$ that contains $v$ and $e$ must contain
at least one edge of $F$. (In particular, this is automatically the case when
$e\in F$.) For each $F\subseteq E$, we define%
\[
\tau\left(  F\right)  =\left\{  e\in E\ \mid\ F\text{ blocks }e\right\}  .
\]
Then, $\tau$ is an antimatroidal quasi-closure operator on $E$. (This example
is the \textquotedblleft line-search antimatroid\textquotedblright\ from
\cite[\S III.2.11]{KoLoSc91}.)

If $\Gamma$ is connected, then $\tau$ is actually a closure operator.
\end{example}

Further examples of antimatroidal closure operators can be found in
\cite[\S III.2]{KoLoSc91} and \cite[\S 3]{EdeJam85}.

We shall be dealing with quasi-closure operators rather than closure operators
most of the time. However, since the latter concept is somewhat more
widespread, let us comment on the connection between the two. Roughly
speaking, the relation between quasi-closure and closure operators is
comparable to the relation between semigroups and monoids, or between
nonunital rings and unital rings, or (perhaps the best analogue) between
simplicial complexes in general and simplicial complexes without ghost
vertices (i.e., simplicial complexes for which every element of the ground set
is a dimension-$0$ face). More concretely, specifying a quasi-closure operator
on a set $E$ is tantamount to specifying a subset of $E$ and a closure
operator on this subset. To wit:

\begin{proposition}
\label{prop.antimat.cl-vs-qcl}Let $E$ be a set. Let $L$ be a subset of $E$.

\textbf{(a)} Then, there is a bijection from%
\[
\left\{  \text{quasi-closure operators }\tau\text{ on }E\text{ satisfying
}\tau\left(  \varnothing\right)  =L\right\}
\]
to%
\[
\left\{  \text{closure operators }\sigma\text{ on }E\setminus L\right\}
\]
that is defined as follows: It sends each quasi-closure operator $\tau$ to the
closure operator $\sigma$ that sends each $F\subseteq E\setminus L$ to
$\tau\left(  F\right)  \setminus L$.

\textbf{(b)} This bijection restricts to a bijection from%
\[
\left\{  \text{antimatroidal quasi-closure operators }\tau\text{ on }E\text{
satisfying }\tau\left(  \varnothing\right)  =L\right\}
\]
to%
\[
\left\{  \text{antimatroidal closure operators }\sigma\text{ on }E\setminus
L\right\}  .
\]

\end{proposition}

Proposition \ref{prop.antimat.cl-vs-qcl} will not be important to what
follows, so we omit the (rather straightforward) proof.

Now, we claim the following:

\begin{theorem}
\label{thm.antimat.Shade}Let $E$ be a set. Let $\tau:\mathcal{P}\left(
E\right)  \rightarrow\mathcal{P}\left(  E\right)  $ be an antimatroidal
quasi-closure operator on $E$. For any $F\subseteq E$, we define%
\begin{equation}
\operatorname*{Shade}F=\left\{  e\in E\ \mid\ e\notin\tau\left(
F\setminus\left\{  e\right\}  \right)  \right\}  .
\label{eq.thm.antimat.Shade.ShadeF=}%
\end{equation}
Then, this map $\operatorname*{Shade}:\mathcal{P}\left(  E\right)
\rightarrow\mathcal{P}\left(  E\right)  $ is an inclusion-reversing shade map.
\end{theorem}

Theorem \ref{thm.antimat.Shade} generalizes Examples \ref{exa.shade.low} and
\ref{exa.shade.convex}. Indeed, applying Theorem \ref{thm.antimat.Shade} to
the setting of Example \ref{exa.antimat.low}, we can easily recover the claim
of Example \ref{exa.shade.low}. Likewise, applying Theorem
\ref{thm.antimat.Shade} to the setting of Example \ref{exa.antimat.convex}, we
can recover the claim of Example \ref{exa.shade.convex}. Less directly, Lemma
\ref{lem.Shade-tog} and its vertex-infection analogue are particular cases of
Theorem \ref{thm.antimat.Shade} as well (even though they involve shade maps
that are inclusion-preserving rather than inclusion-reversing). Indeed, if we
apply Theorem \ref{thm.antimat.Shade} to the setting of Example
\ref{exa.antimat.linesearch}, then we obtain the claim of Lemma
\ref{lem.Shade-tog} with $\operatorname*{Shade}F$ replaced by
$\operatorname*{Shade}\left(  E\setminus F\right)  $; this is easily seen to
be equivalent to Lemma \ref{lem.Shade-tog} (by the duality stated in
Proposition \ref{prop.Shade.dual}).

We shall soon prove Theorem \ref{thm.antimat.Shade}; first we set up two lemmas:

\begin{lemma}
\label{lem.antimat.clos-add}Let $E$ be a set. Let $\tau$ be a quasi-closure
operator on $E$. Let $X$ be a subset of $E$, and let $z\in\tau\left(
X\right)  $. Then, $\tau\left(  X\cup\left\{  z\right\}  \right)  =\tau\left(
X\right)  $.
\end{lemma}

\begin{vershort}


\begin{proof}
[Proof of Lemma \ref{lem.antimat.clos-add}.]Well-known and thus omitted (see
\cite{verlong}).
\end{proof}
\end{vershort}

\begin{verlong}


\begin{proof}
[Proof of Lemma \ref{lem.antimat.clos-add}.]We have $z\in\tau\left(  X\right)
$, so that $\tau\left(  X\right)  \cup\left\{  z\right\}  =\tau\left(
X\right)  $. Hence, $\tau\left(  \tau\left(  X\right)  \cup\left\{  z\right\}
\right)  =\tau\left(  \tau\left(  X\right)  \right)  =\tau\left(  X\right)  $
(by Property 3 in Definition \ref{def.antimat} \textbf{(a)}). However,
$X\subseteq\tau\left(  X\right)  $ (by Property 1 in Definition
\ref{def.antimat} \textbf{(a)}) and thus $\underbrace{X}_{\subseteq\tau\left(
X\right)  }\cup\left\{  z\right\}  \subseteq\tau\left(  X\right)  \cup\left\{
z\right\}  $. Hence, Property 2 in Definition \ref{def.antimat} \textbf{(a)}
(applied to $A=X\cup\left\{  z\right\}  $ and $B=\tau\left(  X\right)
\cup\left\{  z\right\}  $) yields $\tau\left(  X\cup\left\{  z\right\}
\right)  \subseteq\tau\left(  \tau\left(  X\right)  \cup\left\{  z\right\}
\right)  =\tau\left(  X\right)  $. On the other hand, $X\subseteq
X\cup\left\{  z\right\}  $; thus, Property 2 in Definition \ref{def.antimat}
\textbf{(a)} (applied to $A=X$ and $B=X\cup\left\{  z\right\}  $) yields
$\tau\left(  X\right)  \subseteq\tau\left(  X\cup\left\{  z\right\}  \right)
$. Combining this with $\tau\left(  X\cup\left\{  z\right\}  \right)
\subseteq\tau\left(  X\right)  $, we obtain $\tau\left(  X\cup\left\{
z\right\}  \right)  =\tau\left(  X\right)  $. This proves Lemma
\ref{lem.antimat.clos-add}.
\end{proof}
\end{verlong}

\begin{lemma}
\label{lem.antimat.nicer}Let $E$ be a set. Let $\tau$ be an antimatroidal
quasi-closure operator on $E$. Let $X$ be a subset of $E$, and let $y$ and $z$
be two distinct elements of $E$ satisfying $z\in\tau\left(  X\cup\left\{
y\right\}  \right)  $ and $y\in\tau\left(  X\cup\left\{  z\right\}  \right)
$. Then, $y\in\tau\left(  X\right)  $.
\end{lemma}

\begin{proof}
[Proof of Lemma \ref{lem.antimat.nicer}.]Assume the contrary. Thus,
$y\notin\tau\left(  X\right)  $, so that $y\in E\setminus\tau\left(  X\right)
$.

If we had $z\in\tau\left(  X\right)  $, then Lemma \ref{lem.antimat.clos-add}
would yield $\tau\left(  X\cup\left\{  z\right\}  \right)  =\tau\left(
X\right)  $ and therefore $y\in\tau\left(  X\cup\left\{  z\right\}  \right)
=\tau\left(  X\right)  $, which would contradict $y\notin\tau\left(  X\right)
$. Hence, we cannot have $z\in\tau\left(  X\right)  $. Thus, we have
$z\notin\tau\left(  X\right)  $. Therefore, $z\in E\setminus\tau\left(
X\right)  $.

Now, we know that $y$ and $z$ are two distinct elements of $E\setminus
\tau\left(  X\right)  $ satisfying $z\in\tau\left(  X\cup\left\{  y\right\}
\right)  $. Hence, Property 4 in Definition \ref{def.antimat} \textbf{(b)}
shows that $y\notin\tau\left(  X\cup\left\{  z\right\}  \right)  $. This
contradicts $y\in\tau\left(  X\cup\left\{  z\right\}  \right)  $. This
contradiction shows that our assumption was false. Thus, Lemma
\ref{lem.antimat.nicer} is proven.
\end{proof}

Note that Lemma \ref{lem.antimat.nicer} has a converse: If $\tau$ is a
quasi-closure operator on $E$ satisfying the claim of Lemma
\ref{lem.antimat.nicer}, then $\tau$ is antimatroidal. This is easy to see but
will not be used in what follows.

\begin{proof}
[Proof of Theorem \ref{thm.antimat.Shade}.]We shall prove the following three statements:

\begin{statement}
\textit{Statement 0:} If $A$ and $B$ are two subsets of $E$ such that
$A\subseteq B$, then $\operatorname*{Shade}B\subseteq\operatorname*{Shade}A$.
\end{statement}

\begin{statement}
\textit{Statement 1:} If $F\in\mathcal{P}\left(  E\right)  $ and $u\in
E\setminus\operatorname*{Shade}F$, then $\operatorname*{Shade}\left(
F\cup\left\{  u\right\}  \right)  =\operatorname*{Shade}F$.
\end{statement}

\begin{statement}
\textit{Statement 2:} If $F\in\mathcal{P}\left(  E\right)  $ and $u\in
E\setminus\operatorname*{Shade}F$, then $\operatorname*{Shade}\left(
F\setminus\left\{  u\right\}  \right)  =\operatorname*{Shade}F$.
\end{statement}

[\textit{Proof of Statement 0:} Let $A$ and $B$ be two subsets of $E$ such
that $A\subseteq B$. We must prove that $\operatorname*{Shade}B\subseteq
\operatorname*{Shade}A$.

Let $u\in\operatorname*{Shade}B$. Thus,%
\[
u\in\operatorname*{Shade}B=\left\{  e\in E\ \mid\ e\notin\tau\left(
B\setminus\left\{  e\right\}  \right)  \right\}
\]
(by the definition of $\operatorname*{Shade}B$). In other words, $u\in E$ and
$u\notin\tau\left(  B\setminus\left\{  u\right\}  \right)  $.

However, $A\setminus\left\{  u\right\}  \subseteq B\setminus\left\{
u\right\}  $ (since $A\subseteq B$) and thus $\tau\left(  A\setminus\left\{
u\right\}  \right)  \subseteq\tau\left(  B\setminus\left\{  u\right\}
\right)  $ (by Property 2 in Definition \ref{def.antimat} \textbf{(a)},
applied to $A\setminus\left\{  u\right\}  $ and $B\setminus\left\{  u\right\}
$ instead of $A$ and $B$). Hence, from $u\notin\tau\left(  B\setminus\left\{
u\right\}  \right)  $, we obtain $u\notin\tau\left(  A\setminus\left\{
u\right\}  \right)  $. Therefore, $u\in\left\{  e\in E\ \mid\ e\notin%
\tau\left(  A\setminus\left\{  e\right\}  \right)  \right\}
=\operatorname*{Shade}A$ (by the definition of $\operatorname*{Shade}A$).

Forget that we fixed $u$. We thus have shown that $u\in\operatorname*{Shade}A$
for each $u\in\operatorname*{Shade}B$. In other words, $\operatorname*{Shade}%
B\subseteq\operatorname*{Shade}A$. This proves Statement 0.]

[\textit{Proof of Statement 2:} Let $F\in\mathcal{P}\left(  E\right)  $ and
$u\in E\setminus\operatorname*{Shade}F$. We must prove that
$\operatorname*{Shade}\left(  F\setminus\left\{  u\right\}  \right)
=\operatorname*{Shade}F$.

We have $u\in E\setminus\operatorname*{Shade}F=\left\{  e\in E\ \mid\ e\in
\tau\left(  F\setminus\left\{  e\right\}  \right)  \right\}  $ (since
\newline$\operatorname*{Shade}F=\left\{  e\in E\ \mid\ e\notin\tau\left(
F\setminus\left\{  e\right\}  \right)  \right\}  $). In other words, $u\in E$
and $u\in\tau\left(  F\setminus\left\{  u\right\}  \right)  $.

We have $F\setminus\left\{  u\right\}  \subseteq F$ and thus
$\operatorname*{Shade}F\subseteq\operatorname*{Shade}\left(  F\setminus
\left\{  u\right\}  \right)  $ (by Statement 0, applied to $A=F\setminus
\left\{  u\right\}  $ and $B=F$).

Now, let $v\in\operatorname*{Shade}\left(  F\setminus\left\{  u\right\}
\right)  \setminus\operatorname*{Shade}F$. We shall derive a contradiction.

We have
\[
v\in\underbrace{\operatorname*{Shade}\left(  F\setminus\left\{  u\right\}
\right)  }_{\subseteq E}\setminus\operatorname*{Shade}F\subseteq
E\setminus\operatorname*{Shade}F=\left\{  e\in E\ \mid\ e\in\tau\left(
F\setminus\left\{  e\right\}  \right)  \right\}  .
\]
In other words, $v\in E$ and $v\in\tau\left(  F\setminus\left\{  v\right\}
\right)  $.

On the other hand,%
\begin{align*}
v  &  \in\operatorname*{Shade}\left(  F\setminus\left\{  u\right\}  \right)
\setminus\operatorname*{Shade}F\\
&  \subseteq\operatorname*{Shade}\left(  F\setminus\left\{  u\right\}
\right)  =\left\{  e\in E\ \mid\ e\notin\tau\left(  \left(  F\setminus\left\{
u\right\}  \right)  \setminus\left\{  e\right\}  \right)  \right\}
\end{align*}
(by the definition of $\operatorname*{Shade}\left(  F\setminus\left\{
u\right\}  \right)  $). In other words,
\[
v\in E\ \ \ \ \ \ \ \ \ \ \text{and}\ \ \ \ \ \ \ \ \ \ v\notin\tau\left(
\left(  F\setminus\left\{  u\right\}  \right)  \setminus\left\{  v\right\}
\right)  .
\]


Let $X=\left(  F\setminus\left\{  u\right\}  \right)  \setminus\left\{
v\right\}  $. Then, $X\cup\left\{  v\right\}  \supseteq F\setminus\left\{
u\right\}  $, so that $F\setminus\left\{  u\right\}  \subseteq X\cup\left\{
v\right\}  $ and therefore $\tau\left(  F\setminus\left\{  u\right\}  \right)
\subseteq\tau\left(  X\cup\left\{  v\right\}  \right)  $ (by Property 2 in
Definition \ref{def.antimat} \textbf{(a)}, applied to $A=F\setminus\left\{
u\right\}  $ and $B=X\cup\left\{  v\right\}  $). Hence, $u\in\tau\left(
F\setminus\left\{  u\right\}  \right)  \subseteq\tau\left(  X\cup\left\{
v\right\}  \right)  $.

Also, from $X=\left(  F\setminus\left\{  u\right\}  \right)  \setminus\left\{
v\right\}  =\left(  F\setminus\left\{  v\right\}  \right)  \setminus\left\{
u\right\}  $, we obtain $X\cup\left\{  u\right\}  \supseteq F\setminus\left\{
v\right\}  $, so that $F\setminus\left\{  v\right\}  \subseteq X\cup\left\{
u\right\}  $ and therefore $\tau\left(  F\setminus\left\{  v\right\}  \right)
\subseteq\tau\left(  X\cup\left\{  u\right\}  \right)  $ (by Property 2 in
Definition \ref{def.antimat} \textbf{(a)}, applied to $A=F\setminus\left\{
v\right\}  $ and $B=X\cup\left\{  u\right\}  $). Hence, $v\in\tau\left(
F\setminus\left\{  v\right\}  \right)  \subseteq\tau\left(  X\cup\left\{
u\right\}  \right)  $.

If we had $v=u$, then we would have $\left(  F\setminus\left\{  u\right\}
\right)  \setminus\left\{  v\right\}  =\left(  F\setminus\left\{  u\right\}
\right)  \setminus\left\{  u\right\}  =F\setminus\left\{  u\right\}
=F\setminus\left\{  v\right\}  $ (since $u=v$) and therefore $v\notin%
\tau\left(  \underbrace{\left(  F\setminus\left\{  u\right\}  \right)
\setminus\left\{  v\right\}  }_{=F\setminus\left\{  v\right\}  }\right)
=\tau\left(  F\setminus\left\{  v\right\}  \right)  $, which would contradict
$v\in\tau\left(  F\setminus\left\{  v\right\}  \right)  $. Thus, we cannot
have $v=u$. Hence, $v$ and $u$ are distinct.

Thus, Lemma \ref{lem.antimat.nicer} (applied to $y=v$ and $z=u$) yields
$v\in\tau\left(  X\right)  $ (since $v\in\tau\left(  X\cup\left\{  u\right\}
\right)  $ and $u\in\tau\left(  X\cup\left\{  v\right\}  \right)  $). In other
words, $v\in\tau\left(  \left(  F\setminus\left\{  u\right\}  \right)
\setminus\left\{  v\right\}  \right)  $ (since $X=\left(  F\setminus\left\{
u\right\}  \right)  \setminus\left\{  v\right\}  $). But this contradicts
$v\notin\tau\left(  \left(  F\setminus\left\{  u\right\}  \right)
\setminus\left\{  v\right\}  \right)  $.

Forget that we fixed $v$. We thus have found a contradiction for each
$v\in\operatorname*{Shade}\left(  F\setminus\left\{  u\right\}  \right)
\setminus\operatorname*{Shade}F$. Hence, there exists no such $v$. In other
words, the set $\operatorname*{Shade}\left(  F\setminus\left\{  u\right\}
\right)  \setminus\operatorname*{Shade}F$ is empty. Hence,
$\operatorname*{Shade}\left(  F\setminus\left\{  u\right\}  \right)
\subseteq\operatorname*{Shade}F$. Combining this with $\operatorname*{Shade}%
F\subseteq\operatorname*{Shade}\left(  F\setminus\left\{  u\right\}  \right)
$, we obtain $\operatorname*{Shade}\left(  F\setminus\left\{  u\right\}
\right)  =\operatorname*{Shade}F$. This proves Statement 2.]

[\textit{Proof of Statement 1:} Let $F\in\mathcal{P}\left(  E\right)  $ and
$u\in E\setminus\operatorname*{Shade}F$. We must prove that
$\operatorname*{Shade}\left(  F\cup\left\{  u\right\}  \right)
=\operatorname*{Shade}F$. If $u\in F$, then this is obvious (since
$F\cup\left\{  u\right\}  =F$ in this case). Thus, we WLOG assume that
$u\notin F$. Hence, $\left(  F\cup\left\{  u\right\}  \right)  \setminus
\left\{  u\right\}  =F$.

We have $F\subseteq F\cup\left\{  u\right\}  $ and thus $\operatorname*{Shade}%
\left(  F\cup\left\{  u\right\}  \right)  \subseteq\operatorname*{Shade}F$ (by
Statement 0, applied to $A=F$ and $B=F\cup\left\{  u\right\}  $). Hence,
$E\setminus\underbrace{\operatorname*{Shade}\left(  F\cup\left\{  u\right\}
\right)  }_{\subseteq\operatorname*{Shade}F}\supseteq E\setminus
\operatorname*{Shade}F$, so that $E\setminus\operatorname*{Shade}F\subseteq
E\setminus\operatorname*{Shade}\left(  F\cup\left\{  u\right\}  \right)  $.

Now, $u\in E\setminus\operatorname*{Shade}F\subseteq E\setminus
\operatorname*{Shade}\left(  F\cup\left\{  u\right\}  \right)  $. Hence,
Statement 2 (applied to $F\cup\left\{  u\right\}  $ instead of $F$) yields
$\operatorname*{Shade}\left(  \left(  F\cup\left\{  u\right\}  \right)
\setminus\left\{  u\right\}  \right)  =\operatorname*{Shade}\left(
F\cup\left\{  u\right\}  \right)  $. In view of $\left(  F\cup\left\{
u\right\}  \right)  \setminus\left\{  u\right\}  =F$, this rewrites as
$\operatorname*{Shade}F=\operatorname*{Shade}\left(  F\cup\left\{  u\right\}
\right)  $. Hence, $\operatorname*{Shade}\left(  F\cup\left\{  u\right\}
\right)  =\operatorname*{Shade}F$. This proves Statement 1.]

Now, we have proved Statements 1 and 2. Thus, the map $\operatorname*{Shade}%
:\mathcal{P}\left(  E\right)  \rightarrow\mathcal{P}\left(  E\right)  $
satisfies the two axioms in Definition \ref{def.Shade.shademap}. In other
words, this map is a shade map. Moreover, this map is inclusion-reversing (by
Statement 0). Thus, Theorem \ref{thm.antimat.Shade} is proved.
\end{proof}

We note that the quasi-closure operator $\tau$ in Theorem
\ref{thm.antimat.Shade} can be reconstructed from the map
$\operatorname*{Shade}$. This does not even require $\tau$ to be
antimatroidal; the following holds for any quasi-closure operator:

\begin{proposition}
\label{prop.antimat.recover-tau}Let $E$ be a set. Let $\tau:\mathcal{P}\left(
E\right)  \rightarrow\mathcal{P}\left(  E\right)  $ be a quasi-closure
operator on $E$. For any $F\subseteq E$, we define%
\[
\operatorname*{Shade}F=\left\{  e\in E\ \mid\ e\notin\tau\left(
F\setminus\left\{  e\right\}  \right)  \right\}  .
\]
Then, each $F\subseteq E$ satisfies%
\begin{equation}
\tau\left(  F\right)  =F\cup\left(  E\setminus\operatorname*{Shade}F\right)  .
\label{eq.prop.antimat.recover-tau.tauF=}%
\end{equation}

\end{proposition}

\begin{vershort}


\begin{proof}
Exercise (see \cite{verlong} for details).
\end{proof}
\end{vershort}

\begin{verlong}


\begin{proof}
[Proof of Proposition \ref{prop.antimat.recover-tau}.]Let $F$ be a subset of
$E$. We must prove (\ref{eq.prop.antimat.recover-tau.tauF=}).

We shall prove the two inclusions $\tau\left(  F\right)  \subseteq
F\cup\left(  E\setminus\operatorname*{Shade}F\right)  $ and $F\cup\left(
E\setminus\operatorname*{Shade}F\right)  \subseteq\tau\left(  F\right)  $ separately:

[\textit{Proof of }$\tau\left(  F\right)  \subseteq F\cup\left(
E\setminus\operatorname*{Shade}F\right)  $\textit{:} Let $u\in\tau\left(
F\right)  $. We shall show that $u\in F\cup\left(  E\setminus
\operatorname*{Shade}F\right)  $. Indeed, assume the contrary. Thus, $u\notin
F\cup\left(  E\setminus\operatorname*{Shade}F\right)  $. In other words,
$u\notin F$ and $u\notin E\setminus\operatorname*{Shade}F$. From $u\in
\tau\left(  F\right)  \subseteq E$ and $u\notin E\setminus
\operatorname*{Shade}F$, we obtain%
\[
u\in E\setminus\left(  E\setminus\operatorname*{Shade}F\right)  \subseteq
\operatorname*{Shade}F=\left\{  e\in E\ \mid\ e\notin\tau\left(
F\setminus\left\{  e\right\}  \right)  \right\}  .
\]
In other words, $u\in E$ and $u\notin\tau\left(  F\setminus\left\{  u\right\}
\right)  $. However, $u\notin F$ shows that $F\setminus\left\{  u\right\}
=F$. Now, $u\notin\tau\left(  F\setminus\left\{  u\right\}  \right)
=\tau\left(  F\right)  $ (since $F\setminus\left\{  u\right\}  =F$)
contradicts $u\in\tau\left(  F\right)  $. This contradiction shows that our
assumption was false. Hence, $u\in F\cup\left(  E\setminus
\operatorname*{Shade}F\right)  $ is proved.

Now, forget that we fixed $u$. We thus have proved that $u\in F\cup\left(
E\setminus\operatorname*{Shade}F\right)  $ for each $u\in\tau\left(  F\right)
$. In other words, $\tau\left(  F\right)  \subseteq F\cup\left(
E\setminus\operatorname*{Shade}F\right)  $.]

[\textit{Proof of }$F\cup\left(  E\setminus\operatorname*{Shade}F\right)
\subseteq\tau\left(  F\right)  $\textit{:} We have $F\subseteq\tau\left(
F\right)  $ (by Property 1 in Definition \ref{def.antimat} \textbf{(a)}).

Now, let $v\in E\setminus\operatorname*{Shade}F$. Thus,
\[
v\in E\setminus\operatorname*{Shade}F=\left\{  e\in E\ \mid\ e\in\tau\left(
F\setminus\left\{  e\right\}  \right)  \right\}
\]
(since $\operatorname*{Shade}F=\left\{  e\in E\ \mid\ e\notin\tau\left(
F\setminus\left\{  e\right\}  \right)  \right\}  $). In other words, $v\in E$
and $v\in\tau\left(  F\setminus\left\{  v\right\}  \right)  $. However,
$F\setminus\left\{  v\right\}  \subseteq F$ and thus $\tau\left(
F\setminus\left\{  v\right\}  \right)  \subseteq\tau\left(  F\right)  $ (by
Property 2 in Definition \ref{def.antimat} \textbf{(a)}, applied to
$A=F\setminus\left\{  v\right\}  $ and $B=F$). Thus, $v\in\tau\left(
F\setminus\left\{  v\right\}  \right)  \subseteq\tau\left(  F\right)  $.

Forget that we fixed $v$. We thus have shown that $v\in\tau\left(  F\right)  $
for each $v\in E\setminus\operatorname*{Shade}F$. In other words,
$E\setminus\operatorname*{Shade}F\subseteq\tau\left(  F\right)  $. Combining
this with $F\subseteq\tau\left(  F\right)  $, we obtain $F\cup\left(
E\setminus\operatorname*{Shade}F\right)  \subseteq\tau\left(  F\right)  $.]

Combining the two inclusions $\tau\left(  F\right)  \subseteq F\cup\left(
E\setminus\operatorname*{Shade}F\right)  $ and $F\cup\left(  E\setminus
\operatorname*{Shade}F\right)  \subseteq\tau\left(  F\right)  $, we obtain
$\tau\left(  F\right)  =F\cup\left(  E\setminus\operatorname*{Shade}F\right)
$. This proves (\ref{eq.prop.antimat.recover-tau.tauF=}). Proposition
\ref{prop.antimat.recover-tau} is thus proven.
\end{proof}
\end{verlong}

It turns out that if one applies the formula
(\ref{eq.prop.antimat.recover-tau.tauF=}) to an inclusion-reversing shade map
$\operatorname*{Shade}$, then the resulting map $\tau$ is an antimatroidal
quasi-closure operator, at least when $E$ is finite. In fact, we have the following:

\begin{proposition}
\label{prop.antimat.tau-from-Shade}Let $E$ be a finite set. Let
$\operatorname*{Shade}:\mathcal{P}\left(  E\right)  \rightarrow\mathcal{P}%
\left(  E\right)  $ be an inclusion-reversing shade map. Define a map
$\tau:\mathcal{P}\left(  E\right)  \rightarrow\mathcal{P}\left(  E\right)  $
by setting%
\begin{equation}
\tau\left(  F\right)  =F\cup\left(  E\setminus\operatorname*{Shade}F\right)
\label{eq.prop.antimat.tau-from-Shade.tauF=}%
\end{equation}
for each $F\subseteq E$. Then, $\tau$ is an antimatroidal quasi-closure
operator on $E$.
\end{proposition}

The proof of this proposition rests on the following lemma:

\begin{lemma}
\label{lem.Shade.FuG}Let $E$ be a set. Let $\operatorname*{Shade}%
:\mathcal{P}\left(  E\right)  \rightarrow\mathcal{P}\left(  E\right)  $ be a
shade map. Let $A$ and $B$ be two subsets of $E$ such that $B$ is finite and
$B\cap\operatorname*{Shade}A=\varnothing$. Then, $\operatorname*{Shade}\left(
A\cup B\right)  =\operatorname*{Shade}A$.
\end{lemma}

\begin{vershort}


\begin{proof}
[Proof of Lemma \ref{lem.Shade.FuG}.]Induction on $\left\vert B\right\vert $,
using Axiom 1 from Definition \ref{def.Shade.shademap}.
\end{proof}
\end{vershort}

\begin{verlong}


\begin{proof}
[Proof of Lemma \ref{lem.Shade.FuG}.]We proceed by induction on $\left\vert
B\right\vert $:

\textit{Induction base:} Lemma \ref{lem.Shade.FuG} holds for $\left\vert
B\right\vert =0$\ \ \ \ \footnote{\textit{Proof.} Assume that $\left\vert
B\right\vert =0$. Thus, $B=\varnothing$, so that $A\cup B=A\cup\varnothing=A$.
Hence, $\operatorname*{Shade}\left(  A\cup B\right)  =\operatorname*{Shade}A$.
Thus, Lemma \ref{lem.Shade.FuG} is proven for $\left\vert B\right\vert =0$.}.

\textit{Induction step:} Let $k\in\mathbb{N}$. Assume that Lemma
\ref{lem.Shade.FuG} holds for $\left\vert B\right\vert =k$. We must prove that
Lemma \ref{lem.Shade.FuG} holds for $\left\vert B\right\vert =k+1$.

We have assumed that Lemma \ref{lem.Shade.FuG} holds for $\left\vert
B\right\vert =k$. In other words, if $A$ and $B$ are two subsets of $E$ such
that $B$ is finite and $B\cap\operatorname*{Shade}A=\varnothing$ and
$\left\vert B\right\vert =k$, then%
\begin{equation}
\operatorname*{Shade}\left(  A\cup B\right)  =\operatorname*{Shade}A.
\label{pf.lem.Shade.FuG.IH}%
\end{equation}


Now, let $A$ and $B$ be two subsets of $E$ such that $B$ is finite and
$B\cap\operatorname*{Shade}A=\varnothing$ and $\left\vert B\right\vert =k+1$.
We shall prove that $\operatorname*{Shade}\left(  A\cup B\right)
=\operatorname*{Shade}A$.

We assumed that $\operatorname*{Shade}$ is a shade map. Thus, in particular,
$\operatorname*{Shade}$ satisfies Axiom 1 in Definition
\ref{def.Shade.shademap}.

The set $B$ is nonempty (since $\left\vert B\right\vert =k+1>k\geq0$). Thus,
there exists some $u\in B$. Consider this $u$. If we had $u\in
\operatorname*{Shade}A$, then we would have $u\in B\cap\operatorname*{Shade}A$
(since $u\in B$ and $u\in\operatorname*{Shade}A$), which would entail
$B\cap\operatorname*{Shade}A\neq\varnothing$; but this would contradict
$B\cap\operatorname*{Shade}A=\varnothing$. Hence, we cannot have
$u\in\operatorname*{Shade}A$. In other words, we have $u\notin%
\operatorname*{Shade}A$. Hence, $u\in E\setminus\operatorname*{Shade}A$.

From $u\in B$, we obtain $\left\vert B\setminus\left\{  u\right\}  \right\vert
=\left\vert B\right\vert -1=k$ (since $\left\vert B\right\vert =k+1$).
Moreover, $\left(  B\setminus\left\{  u\right\}  \right)  \cap
\operatorname*{Shade}A=\varnothing$ (since $\underbrace{\left(  B\setminus
\left\{  u\right\}  \right)  }_{\subseteq B}\cap\operatorname*{Shade}%
A\subseteq B\cap\operatorname*{Shade}A=\varnothing$). Thus, we can apply
(\ref{pf.lem.Shade.FuG.IH}) to $B\setminus\left\{  u\right\}  $ instead of
$B$. As a consequence, we obtain $\operatorname*{Shade}\left(  A\cup\left(
B\setminus\left\{  u\right\}  \right)  \right)  =\operatorname*{Shade}A$.
Thus, $u\in E\setminus\operatorname*{Shade}\left(  A\cup\left(  B\setminus
\left\{  u\right\}  \right)  \right)  $ (since $u\in E\setminus
\operatorname*{Shade}A$). Thus, Axiom 1 in Definition \ref{def.Shade.shademap}
(applied to $F=A\cup\left(  B\setminus\left\{  u\right\}  \right)  $) yields
\[
\operatorname*{Shade}\left(  \left(  A\cup\left(  B\setminus\left\{
u\right\}  \right)  \right)  \cup\left\{  u\right\}  \right)
=\operatorname*{Shade}\left(  A\cup\left(  B\setminus\left\{  u\right\}
\right)  \right)  =\operatorname*{Shade}A.
\]
In view of
\[
\left(  A\cup\left(  B\setminus\left\{  u\right\}  \right)  \right)
\cup\left\{  u\right\}  =A\cup\underbrace{\left(  B\setminus\left\{
u\right\}  \right)  \cup\left\{  u\right\}  }_{\substack{=B\\\text{(since
}u\in B\text{)}}}=A\cup B,
\]
this rewrites as $\operatorname*{Shade}\left(  A\cup B\right)
=\operatorname*{Shade}A$.

Forget that we fixed $A$ and $B$. We thus have shown that if $A$ and $B$ are
two subsets of $E$ such that $B$ is finite and $B\cap\operatorname*{Shade}%
A=\varnothing$ and $\left\vert B\right\vert =k+1$, then $\operatorname*{Shade}%
\left(  A\cup B\right)  =\operatorname*{Shade}A$. In other words, Lemma
\ref{lem.Shade.FuG} holds for $\left\vert B\right\vert =k+1$. This completes
the induction step. Thus, Lemma \ref{lem.Shade.FuG} is proven.
\end{proof}
\end{verlong}

\begin{vershort}


\begin{proof}
[Proof of Proposition \ref{prop.antimat.tau-from-Shade}.]Again, see
\cite{verlong}. (The proof of Property 3 relies on Lemma \ref{lem.Shade.FuG}.)
\end{proof}
\end{vershort}

\begin{verlong}


\begin{proof}
[Proof of Proposition \ref{prop.antimat.tau-from-Shade}.]We shall show that
the four Properties 1, 2, 3 and 4 from Definition \ref{def.antimat} are
satisfied for our map $\tau$.

[\textit{Proof of Property 1:} Let $A$ be any subset of $E$. We must show that
$A\subseteq\tau\left(  A\right)  $.

The definition of $\tau$ yields $\tau\left(  A\right)  =A\cup\left(
E\setminus\operatorname*{Shade}A\right)  \supseteq A$, so that $A\subseteq
\tau\left(  A\right)  $. Thus, Property 1 is proved.]

[\textit{Proof of Property 2:} Let $A$ and $B$ be two subsets of $E$
satisfying $A\subseteq B$. We must prove that $\tau\left(  A\right)
\subseteq\tau\left(  B\right)  $.

We have $A\subseteq B$ and therefore $\operatorname*{Shade}B\subseteq
\operatorname*{Shade}A$ (since $\operatorname*{Shade}$ is
inclusion-reversing). Thus, $E\setminus\underbrace{\operatorname*{Shade}%
B}_{\subseteq\operatorname*{Shade}A}\supseteq E\setminus\operatorname*{Shade}%
A$, so that $E\setminus\operatorname*{Shade}A\subseteq E\setminus
\operatorname*{Shade}B$. However, the definition of $\tau$ yields $\tau\left(
A\right)  =A\cup\left(  E\setminus\operatorname*{Shade}A\right)  $ and
$\tau\left(  B\right)  =B\cup\left(  E\setminus\operatorname*{Shade}B\right)
$. Therefore,%
\[
\tau\left(  A\right)  =\underbrace{A}_{\subseteq B}\cup\underbrace{\left(
E\setminus\operatorname*{Shade}A\right)  }_{\subseteq E\setminus
\operatorname*{Shade}B}\subseteq B\cup\left(  E\setminus\operatorname*{Shade}%
B\right)  =\tau\left(  B\right)  .
\]
This proves Property 2.]

[\textit{Proof of Property 3:} Let $A$ be any subset of $E$. We must show that
$\tau\left(  \tau\left(  A\right)  \right)  =\tau\left(  A\right)  $.

The set $E\setminus\operatorname*{Shade}A$ is finite (since $E$ is finite) and
satisfies $\left(  E\setminus\operatorname*{Shade}A\right)  \cap
\operatorname*{Shade}A=\varnothing$. Thus, Lemma \ref{lem.Shade.FuG} (applied
to $B=E\setminus\operatorname*{Shade}A$) yields \newline$\operatorname*{Shade}%
\left(  A\cup\left(  E\setminus\operatorname*{Shade}A\right)  \right)
=\operatorname*{Shade}A$. However, the definition of $\tau$ yields
$\tau\left(  A\right)  =A\cup\left(  E\setminus\operatorname*{Shade}A\right)
$. Thus,%
\[
\operatorname*{Shade}\left(  \tau\left(  A\right)  \right)
=\operatorname*{Shade}\left(  A\cup\left(  E\setminus\operatorname*{Shade}%
A\right)  \right)  =\operatorname*{Shade}A.
\]
Now, the definition of $\tau$ yields%
\begin{align*}
\tau\left(  \tau\left(  A\right)  \right)   &  =\underbrace{\tau\left(
A\right)  }_{=A\cup\left(  E\setminus\operatorname*{Shade}A\right)  }%
\cup\left(  E\setminus\underbrace{\operatorname*{Shade}\left(  \tau\left(
A\right)  \right)  }_{=\operatorname*{Shade}A}\right) \\
&  =A\cup\underbrace{\left(  E\setminus\operatorname*{Shade}A\right)
\cup\left(  E\setminus\operatorname*{Shade}A\right)  }_{=\left(
E\setminus\operatorname*{Shade}A\right)  }=A\cup\left(  E\setminus
\operatorname*{Shade}A\right)  =\tau\left(  A\right)  .
\end{align*}
This proves Property 3.]

[\textit{Proof of Property 4:} Let $X$ be a subset of $E$, and let $y$ and $z$
be two distinct elements of $E\setminus\tau\left(  X\right)  $ satisfying
$z\in\tau\left(  X\cup\left\{  y\right\}  \right)  $. We must prove that
$y\notin\tau\left(  X\cup\left\{  z\right\}  \right)  $.

Indeed, assume the contrary. Thus, $y\in\tau\left(  X\cup\left\{  z\right\}
\right)  $.

From $y\in E\setminus\tau\left(  X\right)  $, we obtain $y\notin\tau\left(
X\right)  =X\cup\left(  E\setminus\operatorname*{Shade}X\right)  $ (by the
definition of $\tau$). In other words, $y\notin X$ and $y\notin E\setminus
\operatorname*{Shade}X$. Combining $y\in E$ and $y\notin E\setminus
\operatorname*{Shade}X$, we obtain $y\in E\setminus\left(  E\setminus
\operatorname*{Shade}X\right)  \subseteq\operatorname*{Shade}X$.

We have thus proved that $y\notin X$ and $y\in\operatorname*{Shade}X$. The
same argument (applied to $z$ instead of $y$) shows that $z\notin X$ and
$z\in\operatorname*{Shade}X$.

We assumed that $\operatorname*{Shade}$ is a shade map. Thus, in particular,
$\operatorname*{Shade}$ satisfies both Axioms 1 and 2 in Definition
\ref{def.Shade.shademap}.

We have $z\neq y$ (since $y$ and $z$ are distinct) and thus $z\notin\left\{
y\right\}  $. Combining this with $z\notin X$, we obtain $z\notin
X\cup\left\{  y\right\}  $. However, by assumption, we have%
\[
z\in\tau\left(  X\cup\left\{  y\right\}  \right)  =\left(  X\cup\left\{
y\right\}  \right)  \cup\left(  E\setminus\operatorname*{Shade}\left(
X\cup\left\{  y\right\}  \right)  \right)
\]
(by the definition of $\tau$). In other words, we have $z\in X\cup\left\{
y\right\}  $ or $z\in E\setminus\operatorname*{Shade}\left(  X\cup\left\{
y\right\}  \right)  $. Since $z\in X\cup\left\{  y\right\}  $ is impossible
(because $z\notin X\cup\left\{  y\right\}  $), we thus obtain $z\in
E\setminus\operatorname*{Shade}\left(  X\cup\left\{  y\right\}  \right)  $.
Hence, Axiom 1 in Definition \ref{def.Shade.shademap} (applied to
$F=X\cup\left\{  y\right\}  $ and $u=z$) yields%
\[
\operatorname*{Shade}\left(  \left(  X\cup\left\{  y\right\}  \right)
\cup\left\{  z\right\}  \right)  =\operatorname*{Shade}\left(  X\cup\left\{
y\right\}  \right)  .
\]
The same argument (with the roles of $y$ and $z$ swapped) shows that%
\[
\operatorname*{Shade}\left(  \left(  X\cup\left\{  z\right\}  \right)
\cup\left\{  y\right\}  \right)  =\operatorname*{Shade}\left(  X\cup\left\{
z\right\}  \right)
\]
(since $y\in\tau\left(  X\cup\left\{  z\right\}  \right)  $). Hence,%
\begin{align*}
\operatorname*{Shade}\left(  X\cup\left\{  z\right\}  \right)   &
=\operatorname*{Shade}\left(  \underbrace{\left(  X\cup\left\{  z\right\}
\right)  \cup\left\{  y\right\}  }_{=\left(  X\cup\left\{  y\right\}  \right)
\cup\left\{  z\right\}  }\right) \\
&  =\operatorname*{Shade}\left(  \left(  X\cup\left\{  y\right\}  \right)
\cup\left\{  z\right\}  \right)  =\operatorname*{Shade}\left(  X\cup\left\{
y\right\}  \right)  .
\end{align*}


However, recall that $z\in E\setminus\underbrace{\operatorname*{Shade}\left(
X\cup\left\{  y\right\}  \right)  }_{=\operatorname*{Shade}\left(
X\cup\left\{  z\right\}  \right)  }=E\setminus\operatorname*{Shade}\left(
X\cup\left\{  z\right\}  \right)  $. Therefore, Axiom 2 in Definition
\ref{def.Shade.shademap} (applied to $F=X\cup\left\{  z\right\}  $ and $u=z$)
yields%
\[
\operatorname*{Shade}\left(  \left(  X\cup\left\{  z\right\}  \right)
\setminus\left\{  z\right\}  \right)  =\operatorname*{Shade}\left(
X\cup\left\{  z\right\}  \right)  =\operatorname*{Shade}\left(  X\cup\left\{
y\right\}  \right)  .
\]
Thus,%
\[
\operatorname*{Shade}\left(  X\cup\left\{  y\right\}  \right)
=\operatorname*{Shade}\left(  \underbrace{\left(  X\cup\left\{  z\right\}
\right)  \setminus\left\{  z\right\}  }_{\substack{=X\\\text{(since }z\notin
X\text{)}}}\right)  =\operatorname*{Shade}X.
\]
Recall that $z\in E\setminus\underbrace{\operatorname*{Shade}\left(
X\cup\left\{  y\right\}  \right)  }_{=\operatorname*{Shade}X}=E\setminus
\operatorname*{Shade}X$. Hence, $z\notin\operatorname*{Shade}X$. This clearly
contradicts $z\in\operatorname*{Shade}X$. This contradiction shows that our
assumption was false. Hence, we must have $y\notin\tau\left(  X\cup\left\{
z\right\}  \right)  $. This proves Property 4.]

Now, the map $\tau$ is a quasi-closure operator (since Properties 1, 2 and 3
in Definition \ref{def.antimat} \textbf{(a)} are satisfied for it), and is
antimatroidal (since Property 4 in Definition \ref{def.antimat} \textbf{(b)}
is satisfied for it). Thus, Proposition \ref{prop.antimat.tau-from-Shade} is proved.
\end{proof}
\end{verlong}

Proposition \ref{prop.antimat.recover-tau} has a (sort of) converse:

\begin{proposition}
\label{prop.antimat.recover-Shade}Let $E$ be a set. Let $\operatorname*{Shade}%
:\mathcal{P}\left(  E\right)  \rightarrow\mathcal{P}\left(  E\right)  $ be an
inclusion-reversing shade map. For any $F\subseteq E$, we define%
\[
\tau\left(  F\right)  =F\cup\left(  E\setminus\operatorname*{Shade}F\right)
.
\]
Then, each $F\subseteq E$ satisfies%
\begin{equation}
\operatorname*{Shade}F=\left\{  e\in E\ \mid\ e\notin\tau\left(
F\setminus\left\{  e\right\}  \right)  \right\}  .
\label{eq.prop.antimat.recover-Shade.ShadeF=}%
\end{equation}

\end{proposition}

\begin{vershort}


\begin{proof}
Again, we omit the proof (see \cite{verlong}).
\end{proof}
\end{vershort}

\begin{verlong}


\begin{proof}
[Proof of Proposition \ref{prop.antimat.recover-Shade}.]We assumed that
$\operatorname*{Shade}$ is a shade map. Thus, in particular,
$\operatorname*{Shade}$ satisfies both Axioms 1 and 2 in Definition
\ref{def.Shade.shademap}.

Let $F$ be a subset of $E$. We set $X=\left\{  e\in E\ \mid\ e\notin%
\tau\left(  F\setminus\left\{  e\right\}  \right)  \right\}  $. We shall prove
the inclusions $\operatorname*{Shade}F\subseteq X$ and $X\subseteq
\operatorname*{Shade}F$ separately.

[\textit{Proof of }$\operatorname*{Shade}F\subseteq X$\textit{:} Let
$u\in\operatorname*{Shade}F$. We shall show that $u\in X$.

Indeed, assume the contrary. Thus, $u\notin X$. Combining this with
$u\in\operatorname*{Shade}F\subseteq E$, we find $u\in E\setminus X=\left\{
e\in E\ \mid\ e\in\tau\left(  F\setminus\left\{  e\right\}  \right)  \right\}
$ (since $X=\left\{  e\in E\ \mid\ e\notin\tau\left(  F\setminus\left\{
e\right\}  \right)  \right\}  $). In other words, $u\in E$ and $u\in
\tau\left(  F\setminus\left\{  u\right\}  \right)  $. However, the definition
of $\tau$ yields $\tau\left(  F\setminus\left\{  u\right\}  \right)  =\left(
F\setminus\left\{  u\right\}  \right)  \cup\left(  E\setminus
\operatorname*{Shade}\left(  F\setminus\left\{  u\right\}  \right)  \right)
$. Thus,%
\[
u\in\tau\left(  F\setminus\left\{  u\right\}  \right)  =\left(  F\setminus
\left\{  u\right\}  \right)  \cup\left(  E\setminus\operatorname*{Shade}%
\left(  F\setminus\left\{  u\right\}  \right)  \right)  .
\]
In other words, we have $u\in F\setminus\left\{  u\right\}  $ or $u\in
E\setminus\operatorname*{Shade}\left(  F\setminus\left\{  u\right\}  \right)
$. Since $u\in F\setminus\left\{  u\right\}  $ is impossible, we thus have
$u\in E\setminus\operatorname*{Shade}\left(  F\setminus\left\{  u\right\}
\right)  $. Hence, $u\in E$ and $u\notin\operatorname*{Shade}\left(
F\setminus\left\{  u\right\}  \right)  $. However, $\operatorname*{Shade}$ is
inclusion-reversing; thus, from $F\setminus\left\{  u\right\}  \subseteq F$,
we obtain $\operatorname*{Shade}F\subseteq\operatorname*{Shade}\left(
F\setminus\left\{  u\right\}  \right)  $. Thus, $u\in\operatorname*{Shade}%
F\subseteq\operatorname*{Shade}\left(  F\setminus\left\{  u\right\}  \right)
$. But this contradicts $u\notin\operatorname*{Shade}\left(  F\setminus
\left\{  u\right\}  \right)  $. This contradiction shows that our assumption
was false. Hence, $u\in X$ is proved.

Now, forget that we fixed $u$. We thus have shown that $u\in X$ for each
$u\in\operatorname*{Shade}F$. In other words, $\operatorname*{Shade}F\subseteq
X$.]

[\textit{Proof of }$X\subseteq\operatorname*{Shade}F$\textit{:} Let $u\in X$.
We shall show that $u\in\operatorname*{Shade}F$.

Indeed, assume the contrary. Thus, $u\notin\operatorname*{Shade}F$. This
entails $u\in E\setminus\operatorname*{Shade}F$ (since $u\in X\subseteq E$).
Hence, Axiom 2 in Definition \ref{def.Shade.shademap} yields
$\operatorname*{Shade}\left(  F\setminus\left\{  u\right\}  \right)
=\operatorname*{Shade}F$.

However, from $u\in X=\left\{  e\in E\ \mid\ e\notin\tau\left(  F\setminus
\left\{  e\right\}  \right)  \right\}  $, we see that $u\in E$ and
$u\notin\tau\left(  F\setminus\left\{  u\right\}  \right)  $. Thus,%
\[
u\notin\tau\left(  F\setminus\left\{  u\right\}  \right)  =\left(
F\setminus\left\{  u\right\}  \right)  \cup\left(  E\setminus
\operatorname*{Shade}\left(  F\setminus\left\{  u\right\}  \right)  \right)
\]
(by the definition of $\tau$). Hence, in particular, $u\notin E\setminus
\operatorname*{Shade}\left(  F\setminus\left\{  u\right\}  \right)  $. Since
$u\in E$, we thus have $u\in\operatorname*{Shade}\left(  F\setminus\left\{
u\right\}  \right)  =\operatorname*{Shade}F$. But this contradicts
$u\notin\operatorname*{Shade}F$. This contradiction shows that our assumption
was false. Hence, $u\in\operatorname*{Shade}F$ is proved.

Now, forget that we fixed $u$. We thus have shown that $u\in
\operatorname*{Shade}F$ for each $u\in X$. In other words, $X\subseteq
\operatorname*{Shade}F$.]

Combining $\operatorname*{Shade}F\subseteq X$ with $X\subseteq
\operatorname*{Shade}F$, we obtain $\operatorname*{Shade}F=X=\left\{  e\in
E\ \mid\ e\notin\tau\left(  F\setminus\left\{  e\right\}  \right)  \right\}
$. This proves Proposition \ref{prop.antimat.recover-Shade}.
\end{proof}
\end{verlong}

Combining many of the results in this section, we obtain the following
description of inclusion-reversing shade maps:

\begin{theorem}
\label{thm.antimat.Shade.inc-rev-class}Let $E$ be a finite set. Then, there is
a bijection from the set%
\[
\left\{  \text{inclusion-reversing shade maps }\operatorname*{Shade}%
:\mathcal{P}\left(  E\right)  \rightarrow\mathcal{P}\left(  E\right)
\right\}
\]
to the set%
\[
\left\{  \text{antimatroidal quasi-closure operators }\tau:\mathcal{P}\left(
E\right)  \rightarrow\mathcal{P}\left(  E\right)  \right\}  .
\]
It sends each map $\operatorname*{Shade}$ to the map $\tau$ defined by
(\ref{eq.prop.antimat.tau-from-Shade.tauF=}). Its inverse map sends each map
$\tau$ to the map $\operatorname*{Shade}$ defined by
(\ref{eq.thm.antimat.Shade.ShadeF=}).
\end{theorem}

\begin{vershort}


\begin{proof}
This follows from Theorem \ref{thm.antimat.Shade}, Proposition
\ref{prop.antimat.tau-from-Shade}, Proposition \ref{prop.antimat.recover-tau}
and Proposition \ref{prop.antimat.recover-Shade}.
\end{proof}
\end{vershort}

\begin{verlong}


\begin{proof}
[Proof of Theorem \ref{thm.antimat.Shade.inc-rev-class}.]Let%
\[
\mathcal{S}=\left\{  \text{inclusion-reversing shade maps }%
\operatorname*{Shade}:\mathcal{P}\left(  E\right)  \rightarrow\mathcal{P}%
\left(  E\right)  \right\}
\]
and%
\[
\mathcal{A}=\left\{  \text{antimatroidal quasi-closure operators }%
\tau:\mathcal{P}\left(  E\right)  \rightarrow\mathcal{P}\left(  E\right)
\right\}  .
\]
Then, the map $\mathcal{S}\rightarrow\mathcal{A}$ that sends each map
$\operatorname*{Shade}\in\mathcal{S}$ to the map $\tau$ defined by
(\ref{eq.prop.antimat.tau-from-Shade.tauF=}) is well-defined (by Proposition
\ref{prop.antimat.tau-from-Shade}). Furthermore, the map $\mathcal{A}%
\rightarrow\mathcal{S}$ that sends each map $\tau\in\mathcal{A}$ to the map
$\operatorname*{Shade}$ defined by (\ref{eq.thm.antimat.Shade.ShadeF=}) is
well-defined (by Theorem \ref{thm.antimat.Shade}). These two maps
$\mathcal{S}\rightarrow\mathcal{A}$ and $\mathcal{A}\rightarrow\mathcal{S}$
are mutually inverse (by Proposition \ref{prop.antimat.recover-tau} and
Proposition \ref{prop.antimat.recover-Shade}), and thus are bijections. This
proves Theorem \ref{thm.antimat.Shade.inc-rev-class}.
\end{proof}
\end{verlong}

Theorem \ref{thm.antimat.Shade.inc-rev-class} classifies
\textbf{inclusion-reversing} shade maps in terms of antimatroidal
quasi-closure operators\footnote{In the parlance of matroid theorists, it
shows that inclusion-reversing shade maps are \emph{cryptomorphic} to
antimatroidal quasi-closure operators.}. The latter can in turn be described
in terms of antimatroidal closure operators (by Proposition
\ref{prop.antimat.cl-vs-qcl}), i.e., in terms of antimatroids. Thus,
inclusion-reversing shade maps \textquotedblleft boil down\textquotedblright%
\ to antimatroids. The same can be said of \textbf{inclusion-preserving} shade
maps (because Proposition \ref{prop.Shade.dual} establishes a bijection
between them and the inclusion-reversing ones). In the next subsection, we
shall classify \textbf{arbitrary} shade maps in terms of what we will call
\emph{Boolean interval partitions}.

\subsection{\label{subsect.abstract.bip}Boolean interval partitions and
arbitrary shade maps}

Let us first define Boolean interval partitions:

\begin{definition}
\label{def.bip.bip}Let $E$ be a set.

\textbf{(a)} If $U$ and $V$ are two subsets of $E$ satisfying $U\subseteq V$,
then $\left[  U,V\right]  $ shall denote the subset $\left\{  I\in
\mathcal{P}\left(  E\right)  \ \mid\ U\subseteq I\subseteq V\right\}  $ of
$\mathcal{P}\left(  E\right)  $. This is the set of all subsets of $E$ that
lie between $U$ and $V$ (meaning that they contain $U$ as a subset, but in
turn are contained in $V$ as subsets).

\textbf{(b)} A \emph{Boolean interval} of $\mathcal{P}\left(  E\right)  $
shall mean a subset of $\mathcal{P}\left(  E\right)  $ that has the form
$\left[  U,V\right]  $ for two subsets $U$ and $V$ of $E$ satisfying
$U\subseteq V$. Note that each Boolean interval $\left[  U,V\right]  $ of
$\mathcal{P}\left(  E\right)  $ is nonempty (as it contains $U$ and $V$), and
the two subsets $U$ and $V$ can easily be reconstructed from it (namely, $U$
is the intersection of all $I\in\left[  U,V\right]  $, whereas $V$ is the
union of all $I\in\left[  U,V\right]  $).

\textbf{(c)} A \emph{Boolean interval partition} of $\mathcal{P}\left(
E\right)  $ means a set of pairwise disjoint Boolean intervals of
$\mathcal{P}\left(  E\right)  $ whose union is $\mathcal{P}\left(  E\right)  $.

\textbf{(d)} If $\mathbf{P}$ is a Boolean interval partition of $\mathcal{P}%
\left(  E\right)  $, then the elements of $\mathbf{P}$ (that is, the Boolean
intervals that belong to $\mathbf{P}$) are called the \emph{blocks} of
$\mathbf{P}$.
\end{definition}

\begin{example}
\label{exa.bip.bip.123}For this example, let $E=\left\{  1,2,3\right\}  $. We
shall use the shorthand $i_{1}i_{2}\cdots i_{k}$ for a subset $\left\{
i_{1},i_{2},\ldots,i_{k}\right\}  $ of $E$. (For example, $13$ means the
subset $\left\{  1,3\right\}  $.)

\textbf{(a)} We have
\[
\left[  1,\ 123\right]  =\left\{  I\in\mathcal{P}\left(  123\right)
\ \mid\ 1\subseteq I\subseteq123\right\}  =\left\{  1,\ 12,\ 13,\ 123\right\}
\]
and%
\[
\left[  1,\ 13\right]  =\left\{  I\in\mathcal{P}\left(  123\right)
\ \mid\ 1\subseteq I\subseteq13\right\}  =\left\{  1,\ 13\right\}  .
\]


\textbf{(b)} There are $3^{3}=27$ Boolean intervals of $\mathcal{P}\left(
E\right)  $. (More generally, if $E$ is an $n$-element set, then there are
$3^{n}$ Boolean intervals of $\mathcal{P}\left(  E\right)  $.)

\textbf{(c)} Here is one of many Boolean interval partitions of $\mathcal{P}%
\left(  E\right)  $ (where $E$ is still $\left\{  1,2,3\right\}  $):%
\[
\left\{  \underbrace{\left\{  \varnothing\right\}  }_{=\left[  \varnothing
,\ \varnothing\right]  },\ \ \underbrace{\left\{  1,\ 13\right\}  }_{=\left[
1,\ 13\right]  },\ \ \underbrace{\left\{  3\right\}  }_{=\left[  3,\ 3\right]
},\ \ \underbrace{\left\{  2,\ 12,\ 23,\ 123\right\}  }_{=\left[
2,\ 123\right]  }\right\}  .
\]
Here is another:%
\[
\left\{  \underbrace{\left\{  \varnothing,\ 1\right\}  }_{=\left[
\varnothing,\ 1\right]  },\ \ \underbrace{\left\{  3,\ 13\right\}  }_{=\left[
3,\ 13\right]  },\ \ \underbrace{\left\{  2,\ 23\right\}  }_{=\left[
2,\ 23\right]  },\ \ \underbrace{\left\{  12\right\}  }_{=\left[
12,\ 12\right]  },\ \ \underbrace{\left\{  123\right\}  }_{=\left[
123,\ 123\right]  }\right\}  .
\]
The former has four blocks; the latter has five.
\end{example}

Here are two ways to think of Boolean interval partitions of $\mathcal{P}%
\left(  E\right)  $:

\begin{itemize}
\item The following is just a slick restatement of Definition
\ref{def.bip.bip} \textbf{(c)} using standard combinatorial lingo: A Boolean
interval partition of $\mathcal{P}\left(  E\right)  $ is a set partition of
the Boolean lattice $\mathcal{P}\left(  E\right)  $ into intervals.

\item It is well-known that the set partitions of a given set are in a
canonical bijection with the equivalence relations on this set. In light of
this, the Boolean interval partitions of $\mathcal{P}\left(  E\right)  $ can
be viewed as the equivalence relations on $\mathcal{P}\left(  E\right)  $
whose equivalence classes are Boolean intervals. Equivalently, if $E$ is
finite, then they can be
viewed as the equivalence relations $\sim$ on $\mathcal{P}\left(  E\right)  $
satisfying the axiom \textquotedblleft if $U,V,I\in\mathcal{P}\left(
E\right)  $ satisfy $U\sim V$ and $U\cap V\subseteq I\subseteq U\cup V$, then
$U\sim I\sim V$\textquotedblright. The reader can prove this alternative
characterization as an easy exercise in Boolean algebra.
\end{itemize}

\begin{verlong}
We note a crucial property of Boolean interval partitions (which we shall use tacitly):

\begin{statement}
\textit{Unique block property of Boolean interval partitions:} Let $E$ be a
set. Let $\mathbf{P}$ be a Boolean interval partition of $\mathcal{P}\left(
E\right)  $. Let $F\in\mathcal{P}\left(  E\right)  $. Then, there is a
\textbf{unique} block of $\mathbf{P}$ that contains $F$. Moreover, this block
is a Boolean interval, i.e., it has the form $\left[  \alpha\left(  F\right)
,\ \tau\left(  F\right)  \right]  $ for some $\alpha\left(  F\right)
\in\mathcal{P}\left(  E\right)  $ and some $\tau\left(  F\right)
\in\mathcal{P}\left(  E\right)  $. Finally, these $\alpha\left(  F\right)  $
and $\tau\left(  F\right)  $ are uniquely determined by $F$ and $\mathbf{P}$.
\end{statement}
\end{verlong}

Boolean interval partitions have come up in combinatorics before (e.g.,
\cite{BrOlNo09}, \cite{Dawson80}, \cite{DedTit20}, \cite{GorMah97}).

We shall now construct a shade map from any Boolean interval partition:

\begin{theorem}
\label{thm.bip.Bip-to-Shade}Let $E$ be a set. Let $\mathbf{P}$ be a Boolean
interval partition of $\mathcal{P}\left(  E\right)  $.

For any $F\in\mathcal{P}\left(  E\right)  $, let $\left[  \alpha\left(
F\right)  ,\ \tau\left(  F\right)  \right]  $ denote the (unique) block of
$\mathbf{P}$ that contains $F$.

We define a map $\operatorname*{Shade}:\mathcal{P}\left(  E\right)
\rightarrow\mathcal{P}\left(  E\right)  $ by setting%
\[
\operatorname*{Shade}F=E\setminus\left(  \tau\left(  F\right)  \setminus
\alpha\left(  F\right)  \right)  \ \ \ \ \ \ \ \ \ \ \text{for any }%
F\in\mathcal{P}\left(  E\right)  .
\]


Then:

\textbf{(a)} The map $\operatorname*{Shade}$ is a shade map on $E$.

\textbf{(b)} We have $\alpha\left(  F\right)  =F\cap\operatorname*{Shade}F$
and $\tau\left(  F\right)  =F\cup\left(  E\setminus\operatorname*{Shade}%
F\right)  $ for any $F\in\mathcal{P}\left(  E\right)  $.

\textbf{(c)} We have $\mathbf{P}=\left\{  \left[  \alpha\left(  F\right)
,\ \tau\left(  F\right)  \right]  \ \mid\ F\in\mathcal{P}\left(  E\right)
\right\}  $.
\end{theorem}

The proof of Theorem \ref{thm.bip.Bip-to-Shade} is rather easy. We lighten our
burden somewhat with a simple lemma:

\begin{lemma}
\label{lem.bip.Bip-to-Shade.lemb}Let $X$, $Y$, $Z$ and $E$ be four sets such
that $X\subseteq Y\subseteq Z\subseteq E$. Then:

\textbf{(a)} We have $Y\cup\left(  Z\setminus X\right)  =Z$.

\textbf{(b)} We have $Y\cap\left(  E\setminus\left(  Z\setminus X\right)
\right)  =X$.
\end{lemma}

\begin{proof}
[Proof of Lemma \ref{lem.bip.Bip-to-Shade.lemb}.]This is straightforward set
theory; it is easily checked (e.g.) using Venn diagrams.
\end{proof}

\begin{proof}
[Proof of Theorem \ref{thm.bip.Bip-to-Shade}.]\textbf{(a)} We shall prove the
following two statements:

\begin{statement}
\textit{Statement 1:} If $F\in\mathcal{P}\left(  E\right)  $ and $u\in
E\setminus\operatorname*{Shade}F$, then $\operatorname*{Shade}\left(
F\cup\left\{  u\right\}  \right)  =\operatorname*{Shade}F$.
\end{statement}

\begin{statement}
\textit{Statement 2:} If $F\in\mathcal{P}\left(  E\right)  $ and $u\in
E\setminus\operatorname*{Shade}F$, then $\operatorname*{Shade}\left(
F\setminus\left\{  u\right\}  \right)  =\operatorname*{Shade}F$.
\end{statement}

[\textit{Proof of Statement 1:} Let $F\in\mathcal{P}\left(  E\right)  $ and
$u\in E\setminus\operatorname*{Shade}F$. We must prove that
$\operatorname*{Shade}\left(  F\cup\left\{  u\right\}  \right)
=\operatorname*{Shade}F$.

\begin{vershort}
The definition of $\alpha\left(  F\right)  $ and $\tau\left(  F\right)  $
reveals that $\left[  \alpha\left(  F\right)  ,\ \tau\left(  F\right)
\right]  $ is the (unique) block of $\mathbf{P}$ that contains $F$. Thus,
$\left[  \alpha\left(  F\right)  ,\ \tau\left(  F\right)  \right]  $ is a
block of $\mathbf{P}$, so that $\alpha\left(  F\right)  $ and $\tau\left(
F\right)  $ are subsets of $E$. Hence, $\tau\left(  F\right)  \setminus
\alpha\left(  F\right)  $ is a subset of $E$ as well.
\end{vershort}

\begin{verlong}
The definition of $\alpha\left(  F\right)  $ and $\tau\left(  F\right)  $
reveals that $\left[  \alpha\left(  F\right)  ,\ \tau\left(  F\right)
\right]  $ is the (unique) block of $\mathbf{P}$ that contains $F$. Thus,
$\left[  \alpha\left(  F\right)  ,\ \tau\left(  F\right)  \right]  $ is a
block of $\mathbf{P}$, therefore a Boolean interval of $\mathcal{P}\left(
E\right)  $ (since $\mathbf{P}$ is a Boolean interval partition of
$\mathcal{P}\left(  E\right)  $). Hence, $\alpha\left(  F\right)  $ and
$\tau\left(  F\right)  $ belong to $\mathcal{P}\left(  E\right)  $, hence are
subsets of $E$. Hence, $\tau\left(  F\right)  \setminus\alpha\left(  F\right)
$ is a subset of $E$ as well.
\end{verlong}

The definition of $\operatorname*{Shade}$ yields $\operatorname*{Shade}%
F=E\setminus\left(  \tau\left(  F\right)  \setminus\alpha\left(  F\right)
\right)  $. Now,%
\[
u\in E\setminus\underbrace{\operatorname*{Shade}F}_{=E\setminus\left(
\tau\left(  F\right)  \setminus\alpha\left(  F\right)  \right)  }%
=E\setminus\left(  E\setminus\left(  \tau\left(  F\right)  \setminus
\alpha\left(  F\right)  \right)  \right)  =\tau\left(  F\right)
\setminus\alpha\left(  F\right)
\]
(since $\tau\left(  F\right)  \setminus\alpha\left(  F\right)  $ is a subset
of $E$). In other words, $u\in\tau\left(  F\right)  $ and $u\notin%
\alpha\left(  F\right)  $.

On the other hand, the Boolean interval $\left[  \alpha\left(  F\right)
,\ \tau\left(  F\right)  \right]  $ contains $F$ (since $\left[  \alpha\left(
F\right)  ,\ \tau\left(  F\right)  \right]  $ is the (unique) block of
$\mathbf{P}$ that contains $F$). In other words, $\alpha\left(  F\right)
\subseteq F\subseteq\tau\left(  F\right)  $.

\begin{vershort}
Now, set $F^{\prime}=F\cup\left\{  u\right\}  $. From $F\subseteq\tau\left(
F\right)  $ and $u\in\tau\left(  F\right)  $, we thus obtain $F^{\prime
}\subseteq\tau\left(  F\right)  $. Combined with $\alpha\left(  F\right)
\subseteq F\subseteq F^{\prime}$, this entails $F^{\prime}\in\left[
\alpha\left(  F\right)  ,\ \tau\left(  F\right)  \right]  $. Hence, $\left[
\alpha\left(  F\right)  ,\ \tau\left(  F\right)  \right]  $ is a block of
$\mathbf{P}$ that contains $F^{\prime}$ (since we already know that $\left[
\alpha\left(  F\right)  ,\ \tau\left(  F\right)  \right]  $ is a block of
$\mathbf{P}$).
\end{vershort}

\begin{verlong}
Now, set $F^{\prime}=F\cup\left\{  u\right\}  $. From $u\in\tau\left(
F\right)  $, we obtain $\left\{  u\right\}  \subseteq\tau\left(  F\right)  $.
Combining $F\subseteq\tau\left(  F\right)  $ with $\left\{  u\right\}
\subseteq\tau\left(  F\right)  $, we obtain $F\cup\left\{  u\right\}
\subseteq\tau\left(  F\right)  $. In other words, $F^{\prime}\subseteq
\tau\left(  F\right)  $ (since $F^{\prime}=F\cup\left\{  u\right\}  $).
Combining this with $\alpha\left(  F\right)  \subseteq F\subseteq
F\cup\left\{  u\right\}  =F^{\prime}$, we obtain the chain of inclusions
$\alpha\left(  F\right)  \subseteq F^{\prime}\subseteq\tau\left(  F\right)  $.
In other words, $F^{\prime}\in\left[  \alpha\left(  F\right)  ,\ \tau\left(
F\right)  \right]  $. Hence, $\left[  \alpha\left(  F\right)  ,\ \tau\left(
F\right)  \right]  $ is a block of $\mathbf{P}$ that contains $F^{\prime}$
(since we already know that $\left[  \alpha\left(  F\right)  ,\ \tau\left(
F\right)  \right]  $ is a block of $\mathbf{P}$).
\end{verlong}

However, the definition of $\alpha\left(  F^{\prime}\right)  $ and
$\tau\left(  F^{\prime}\right)  $ shows that $\left[  \alpha\left(  F^{\prime
}\right)  ,\ \tau\left(  F^{\prime}\right)  \right]  $ is the (unique) block
of $\mathbf{P}$ that contains $F^{\prime}$. Since we know that $\left[
\alpha\left(  F\right)  ,\ \tau\left(  F\right)  \right]  $ is a block of
$\mathbf{P}$ that contains $F^{\prime}$, we therefore conclude that $\left[
\alpha\left(  F^{\prime}\right)  ,\ \tau\left(  F^{\prime}\right)  \right]
=\left[  \alpha\left(  F\right)  ,\ \tau\left(  F\right)  \right]  $. Hence,
we have%
\[
\alpha\left(  F^{\prime}\right)  =\alpha\left(  F\right)
\ \ \ \ \ \ \ \ \ \ \text{and}\ \ \ \ \ \ \ \ \ \ \tau\left(  F^{\prime
}\right)  =\tau\left(  F\right)
\]
(because a Boolean interval $\left[  U,\ V\right]  $ uniquely determines both
$U$ and $V$). Now, the definition of $\operatorname*{Shade}$ yields%
\[
\operatorname*{Shade}\left(  F^{\prime}\right)  =E\setminus\left(
\underbrace{\tau\left(  F^{\prime}\right)  }_{=\tau\left(  F\right)
}\setminus\underbrace{\alpha\left(  F^{\prime}\right)  }_{=\alpha\left(
F\right)  }\right)  =E\setminus\left(  \tau\left(  F\right)  \setminus
\alpha\left(  F\right)  \right)  =\operatorname*{Shade}F.
\]
In view of $F^{\prime}=F\cup\left\{  u\right\}  $, this rewrites as
$\operatorname*{Shade}\left(  F\cup\left\{  u\right\}  \right)
=\operatorname*{Shade}F$. This proves Statement 1.]

[\textit{Proof of Statement 2:} Let $F\in\mathcal{P}\left(  E\right)  $ and
$u\in E\setminus\operatorname*{Shade}F$. We must prove that
$\operatorname*{Shade}\left(  F\setminus\left\{  u\right\}  \right)
=\operatorname*{Shade}F$.

\begin{vershort}
We proceed exactly as in our above proof of Statement 1 up until the point
where we define $F^{\prime}$. Instead of setting $F^{\prime}=F\cup\left\{
u\right\}  $, we now set $F^{\prime}=F\setminus\left\{  u\right\}  $.
Combining $\alpha\left(  F\right)  \subseteq F$ with $u\notin\alpha\left(
F\right)  $, we obtain $\alpha\left(  F\right)  \subseteq F\setminus\left\{
u\right\}  =F^{\prime}$. Combining this with $F^{\prime}=F\setminus\left\{
u\right\}  \subseteq F\subseteq\tau\left(  F\right)  $, we see that
$F^{\prime}\in\left[  \alpha\left(  F\right)  ,\ \tau\left(  F\right)
\right]  $. From this, we can obtain $\operatorname*{Shade}\left(  F^{\prime
}\right)  =\operatorname*{Shade}F$ by the same argument that we used back in
the proof of Statement 1. In view of $F^{\prime}=F\setminus\left\{  u\right\}
$, this rewrites as $\operatorname*{Shade}\left(  F\setminus\left\{
u\right\}  \right)  =\operatorname*{Shade}F$. This proves Statement 2.]
\end{vershort}

\begin{verlong}
We proceed exactly as in our above proof of Statement 1 up until the point
where we define $F^{\prime}$. Instead of setting $F^{\prime}=F\cup\left\{
u\right\}  $, we now set $F^{\prime}=F\setminus\left\{  u\right\}  $. Now, let
us call to mind the obvious fact that if two sets $U$ and $V$ and some object
$w$ satisfy $U\subseteq V$ and $w\notin U$, then $U\subseteq V\setminus
\left\{  w\right\}  $. We can apply this to $U=\alpha\left(  F\right)  $ and
$V=F$ and $w=u$ (since $\alpha\left(  F\right)  \subseteq F$ and
$u\notin\alpha\left(  F\right)  $); thus we obtain $\alpha\left(  F\right)
\subseteq F\setminus\left\{  u\right\}  $. In other words, $\alpha\left(
F\right)  \subseteq F^{\prime}$ (since $F^{\prime}=F\setminus\left\{
u\right\}  $). Combining this with $F^{\prime}=F\setminus\left\{  u\right\}
\subseteq F\subseteq\tau\left(  F\right)  $, we obtain the chain of inclusions
$\alpha\left(  F\right)  \subseteq F^{\prime}\subseteq\tau\left(  F\right)  $.
From this chain, we can obtain $\operatorname*{Shade}\left(  F^{\prime
}\right)  =\operatorname*{Shade}F$ by the same argument that we used back in
the proof of Statement 1. In view of $F^{\prime}=F\setminus\left\{  u\right\}
$, this rewrites as $\operatorname*{Shade}\left(  F\setminus\left\{
u\right\}  \right)  =\operatorname*{Shade}F$. This proves Statement 2.]
\end{verlong}

Now, we have proved Statements 1 and 2. Thus, the map $\operatorname*{Shade}%
:\mathcal{P}\left(  E\right)  \rightarrow\mathcal{P}\left(  E\right)  $
satisfies the two axioms in Definition \ref{def.Shade.shademap}. In other
words, this map is a shade map. This proves Theorem \ref{thm.bip.Bip-to-Shade}
\textbf{(a)}.

\textbf{(b)} Let $F\in\mathcal{P}\left(  E\right)  $. We must prove that
$\alpha\left(  F\right)  =F\cap\operatorname*{Shade}F$ and $\tau\left(
F\right)  =F\cup\left(  E\setminus\operatorname*{Shade}F\right)  $.

The definition of $\alpha\left(  F\right)  $ and $\tau\left(  F\right)  $
reveals that $\left[  \alpha\left(  F\right)  ,\ \tau\left(  F\right)
\right]  $ is the (unique) block of $\mathbf{P}$ that contains $F$. Hence,
$\left[  \alpha\left(  F\right)  ,\ \tau\left(  F\right)  \right]  $ contains
$F$. In other words, $\alpha\left(  F\right)  \subseteq F\subseteq\tau\left(
F\right)  $. Moreover, $\left[  \alpha\left(  F\right)  ,\ \tau\left(
F\right)  \right]  $ is a block of $\mathbf{P}$ (since $\left[  \alpha\left(
F\right)  ,\ \tau\left(  F\right)  \right]  $ is the (unique) block of
$\mathbf{P}$ that contains $F$), therefore a Boolean interval of
$\mathcal{P}\left(  E\right)  $ (since $\mathbf{P}$ is a Boolean interval
partition of $\mathcal{P}\left(  E\right)  $). Hence, $\alpha\left(  F\right)
$ and $\tau\left(  F\right)  $ belong to $\mathcal{P}\left(  E\right)  $,
hence are subsets of $E$. Hence, $\tau\left(  F\right)  \setminus\alpha\left(
F\right)  $ is a subset of $E$ as well.

The definition of $\operatorname*{Shade}$ yields $\operatorname*{Shade}%
F=E\setminus\left(  \tau\left(  F\right)  \setminus\alpha\left(  F\right)
\right)  $. Hence,
\[
E\setminus\underbrace{\operatorname*{Shade}F}_{=E\setminus\left(  \tau\left(
F\right)  \setminus\alpha\left(  F\right)  \right)  }=E\setminus\left(
E\setminus\left(  \tau\left(  F\right)  \setminus\alpha\left(  F\right)
\right)  \right)  =\tau\left(  F\right)  \setminus\alpha\left(  F\right)
\]
(since $\tau\left(  F\right)  \setminus\alpha\left(  F\right)  $ is a subset
of $E$).

Now, $\alpha\left(  F\right)  \subseteq F\subseteq\tau\left(  F\right)
\subseteq E$ (since $\tau\left(  F\right)  $ is a subset of $E$). Hence, Lemma
\ref{lem.bip.Bip-to-Shade.lemb} \textbf{(b)} (applied to $X=\alpha\left(
F\right)  $ and $Y=F$ and $Z=\tau\left(  F\right)  $) yields that
$F\cap\left(  E\setminus\left(  \tau\left(  F\right)  \setminus\alpha\left(
F\right)  \right)  \right)  =\alpha\left(  F\right)  $. In view of
$\operatorname*{Shade}F=E\setminus\left(  \tau\left(  F\right)  \setminus
\alpha\left(  F\right)  \right)  $, this rewrites as $F\cap
\operatorname*{Shade}F=\alpha\left(  F\right)  $. In other words,
$\alpha\left(  F\right)  =F\cap\operatorname*{Shade}F$.

Furthermore, Lemma \ref{lem.bip.Bip-to-Shade.lemb} \textbf{(a)} (applied to
$X=\alpha\left(  F\right)  $ and $Y=F$ and $Z=\tau\left(  F\right)  $) yields
that $F\cup\left(  \tau\left(  F\right)  \setminus\alpha\left(  F\right)
\right)  =\tau\left(  F\right)  $. In view of $E\setminus\operatorname*{Shade}%
F=\tau\left(  F\right)  \setminus\alpha\left(  F\right)  $, this rewrites as
$F\cup\left(  E\setminus\operatorname*{Shade}F\right)  =\tau\left(  F\right)
$. In other words, $\tau\left(  F\right)  =F\cup\left(  E\setminus
\operatorname*{Shade}F\right)  $. Thus, Theorem \ref{thm.bip.Bip-to-Shade}
\textbf{(b)} is proven.

\begin{vershort}
\textbf{(c)} Each block of $\mathbf{P}$ has the form $\left[  \alpha\left(
F\right)  ,\ \tau\left(  F\right)  \right]  $ for some $F\in\mathcal{P}\left(
E\right)  $ (since it is a Boolean interval, thus nonempty, therefore contains
some $F\in\mathcal{P}\left(  E\right)  $; but then it must be $\left[
\alpha\left(  F\right)  ,\ \tau\left(  F\right)  \right]  $ for this $F$).
Conversely, any set of the form $\left[  \alpha\left(  F\right)
,\ \tau\left(  F\right)  \right]  $ is a block of $\mathbf{P}$ (by the
definition of $\left[  \alpha\left(  F\right)  ,\ \tau\left(  F\right)
\right]  $). Combining these two facts, we conclude that the blocks of
$\mathbf{P}$ are precisely the sets of the form $\left[  \alpha\left(
F\right)  ,\ \tau\left(  F\right)  \right]  $ with $F\in\mathcal{P}\left(
E\right)  $. But this is precisely the claim of Theorem
\ref{thm.bip.Bip-to-Shade} \textbf{(c)}. \qedhere

\end{vershort}

\begin{verlong}
\textbf{(c)} It is clear that $\mathbf{P}\subseteq\left\{  \left[
\alpha\left(  F\right)  ,\ \tau\left(  F\right)  \right]  \ \mid
\ F\in\mathcal{P}\left(  E\right)  \right\}  $%
\ \ \ \ \footnote{\textit{Proof.} Let $Q\in\mathbf{P}$. Hence, $Q$ is a
Boolean interval of $\mathcal{P}\left(  E\right)  $ (since $\mathbf{P}$ is a
Boolean interval partition of $\mathcal{P}\left(  E\right)  $). Therefore, $Q$
is nonempty (since every Boolean interval is nonempty). In other words, there
exists some $F\in Q$. Consider this $F$. Thus, $Q$ contains $F$ (since $F\in
Q$). Moreover, $Q$ is a block of $\mathbf{P}$ (since $Q\in\mathbf{P}$), and we
have $F\in Q\subseteq\mathcal{P}\left(  E\right)  $ (since $Q$ is a Boolean
interval of $\mathcal{P}\left(  E\right)  $).
\par
The definition of $\alpha\left(  F\right)  $ and $\tau\left(  F\right)  $
reveals that $\left[  \alpha\left(  F\right)  ,\ \tau\left(  F\right)
\right]  $ is the (unique) block of $\mathbf{P}$ that contains $F$. Thus,
$\left[  \alpha\left(  F\right)  ,\ \tau\left(  F\right)  \right]  =Q$ (since
$Q$ is a block of $\mathbf{P}$ that contains $F$). In other words, $Q=\left[
\alpha\left(  F\right)  ,\ \tau\left(  F\right)  \right]  $.
\par
Thus, $Q$ has the form $\left[  \alpha\left(  G\right)  ,\ \tau\left(
G\right)  \right]  $ for some $G\in\mathcal{P}\left(  E\right)  $ (namely, for
$G=F$). In other words, $Q\in\left\{  \left[  \alpha\left(  G\right)
,\ \tau\left(  G\right)  \right]  \ \mid\ G\in\mathcal{P}\left(  E\right)
\right\}  $.
\par
Forget that we fixed $Q$. We thus have shown that $Q\in\left\{  \left[
\alpha\left(  G\right)  ,\ \tau\left(  G\right)  \right]  \ \mid
\ G\in\mathcal{P}\left(  E\right)  \right\}  $ for each $Q\in\mathbf{P}$. In
other words, $\mathbf{P}\subseteq\left\{  \left[  \alpha\left(  G\right)
,\ \tau\left(  G\right)  \right]  \ \mid\ G\in\mathcal{P}\left(  E\right)
\right\}  $. Hence,%
\[
\mathbf{P}\subseteq\left\{  \left[  \alpha\left(  G\right)  ,\ \tau\left(
G\right)  \right]  \ \mid\ G\in\mathcal{P}\left(  E\right)  \right\}
=\left\{  \left[  \alpha\left(  F\right)  ,\ \tau\left(  F\right)  \right]
\ \mid\ F\in\mathcal{P}\left(  E\right)  \right\}
\]
(here, we have renamed the index $G$ as $F$).} and \newline$\left\{  \left[
\alpha\left(  F\right)  ,\ \tau\left(  F\right)  \right]  \ \mid
\ F\in\mathcal{P}\left(  E\right)  \right\}  \subseteq\mathbf{P}%
$\ \ \ \ \footnote{\textit{Proof.} Let $F\in\mathcal{P}\left(  E\right)  $.
The definition of $\alpha\left(  F\right)  $ and $\tau\left(  F\right)  $
reveals that $\left[  \alpha\left(  F\right)  ,\ \tau\left(  F\right)
\right]  $ is the (unique) block of $\mathbf{P}$ that contains $F$. Thus,
$\left[  \alpha\left(  F\right)  ,\ \tau\left(  F\right)  \right]  $ is a
block of $\mathbf{P}$. In other words, $\left[  \alpha\left(  F\right)
,\ \tau\left(  F\right)  \right]  \in\mathbf{P}$.
\par
Forget that we fixed $F$. We thus have shown that $\left[  \alpha\left(
F\right)  ,\ \tau\left(  F\right)  \right]  \in\mathbf{P}$ for each
$F\in\mathcal{P}\left(  E\right)  $. In other words, $\left\{  \left[
\alpha\left(  F\right)  ,\ \tau\left(  F\right)  \right]  \ \mid
\ F\in\mathcal{P}\left(  E\right)  \right\}  \subseteq\mathbf{P}$.}. Combining
these two inclusions, we obtain $\mathbf{P}=\left\{  \left[  \alpha\left(
F\right)  ,\ \tau\left(  F\right)  \right]  \ \mid\ F\in\mathcal{P}\left(
E\right)  \right\}  $. This proves Theorem \ref{thm.bip.Bip-to-Shade}
\textbf{(c)}. \qedhere

\end{verlong}
\end{proof}

A converse to Theorem \ref{thm.bip.Bip-to-Shade} is provided by the following theorem:

\begin{theorem}
\label{thm.bip.Shade-to-Bip}Let $E$ be a finite set. Let
$\operatorname*{Shade}:\mathcal{P}\left(  E\right)  \rightarrow\mathcal{P}%
\left(  E\right)  $ be a shade map on $E$. Define a map $\alpha:\mathcal{P}%
\left(  E\right)  \rightarrow\mathcal{P}\left(  E\right)  $ by setting%
\[
\alpha\left(  F\right)  =F\cap\operatorname*{Shade}%
F\ \ \ \ \ \ \ \ \ \ \text{for any }F\in\mathcal{P}\left(  E\right)  .
\]
Define a map $\tau:\mathcal{P}\left(  E\right)  \rightarrow\mathcal{P}\left(
E\right)  $ by setting%
\[
\tau\left(  F\right)  =F\cup\left(  E\setminus\operatorname*{Shade}F\right)
\ \ \ \ \ \ \ \ \ \ \text{for any }F\in\mathcal{P}\left(  E\right)  .
\]
Let
\begin{equation}
\mathbf{P}=\left\{  \left[  \alpha\left(  F\right)  ,\ \tau\left(  F\right)
\right]  \ \mid\ F\in\mathcal{P}\left(  E\right)  \right\}  .
\label{eq.thm.bip.Shade-to-Bip.P=}%
\end{equation}


Then:

\textbf{(a)} We have $F\in\left[  \alpha\left(  F\right)  ,\ \tau\left(
F\right)  \right]  $ for any $F\in\mathcal{P}\left(  E\right)  $.

\textbf{(b)} If $F\in\mathcal{P}\left(  E\right)  $ and $G\in\left[
\alpha\left(  F\right)  ,\ \tau\left(  F\right)  \right]  $, then $\left[
\alpha\left(  F\right)  ,\ \tau\left(  F\right)  \right]  =\left[
\alpha\left(  G\right)  ,\ \tau\left(  G\right)  \right]  $.

\textbf{(c)} The set $\mathbf{P}$ is a Boolean interval partition of
$\mathcal{P}\left(  E\right)  $.

\textbf{(d)} We have $\operatorname*{Shade}F=E\setminus\left(  \tau\left(
F\right)  \setminus\alpha\left(  F\right)  \right)  $ for any $F\in
\mathcal{P}\left(  E\right)  $.
\end{theorem}

To prove this theorem, we will need the following variant of Lemma
\ref{lem.Shade.FuG}:

\begin{lemma}
\label{lem.Shade.FoG}Let $E$ be a set. Let $\operatorname*{Shade}%
:\mathcal{P}\left(  E\right)  \rightarrow\mathcal{P}\left(  E\right)  $ be a
shade map. Let $A$ and $B$ be two subsets of $E$ such that $B$ is finite and
$B\cap\operatorname*{Shade}A=\varnothing$. Then, $\operatorname*{Shade}\left(
A\setminus B\right)  =\operatorname*{Shade}A$.
\end{lemma}

\begin{vershort}
\begin{proof}
[Proof of Lemma \ref{lem.Shade.FoG}.]Induction on $\left\vert B\right\vert $,
using Axiom 2 from Definition \ref{def.Shade}.
\end{proof}
\end{vershort}

\begin{verlong}
\begin{proof}
[Proof of Lemma \ref{lem.Shade.FoG}.]We proceed by induction on $\left\vert
B\right\vert $:

\textit{Induction base:} Lemma \ref{lem.Shade.FoG} holds for $\left\vert
B\right\vert =0$\ \ \ \ \footnote{\textit{Proof.} Assume that $\left\vert
B\right\vert =0$. Thus, $B=\varnothing$, so that $A\setminus B=A\setminus
\varnothing=A$. Hence, $\operatorname*{Shade}\left(  A\setminus B\right)
=\operatorname*{Shade}A$. Thus, Lemma \ref{lem.Shade.FoG} is proven for
$\left\vert B\right\vert =0$.}.

\textit{Induction step:} Let $k\in\mathbb{N}$. Assume that Lemma
\ref{lem.Shade.FoG} holds for $\left\vert B\right\vert =k$. We must prove that
Lemma \ref{lem.Shade.FoG} holds for $\left\vert B\right\vert =k+1$.

We have assumed that Lemma \ref{lem.Shade.FoG} holds for $\left\vert
B\right\vert =k$. In other words, if $A$ and $B$ are two subsets of $E$ such
that $B$ is finite and $B\cap\operatorname*{Shade}A=\varnothing$ and
$\left\vert B\right\vert =k$, then%
\begin{equation}
\operatorname*{Shade}\left(  A\setminus B\right)  =\operatorname*{Shade}A.
\label{pf.lem.Shade.FoG.IH}%
\end{equation}


Now, let $A$ and $B$ be two subsets of $E$ such that $B$ is finite and
$B\cap\operatorname*{Shade}A=\varnothing$ and $\left\vert B\right\vert =k+1$.
We shall prove that $\operatorname*{Shade}\left(  A\setminus B\right)
=\operatorname*{Shade}A$.

We assumed that $\operatorname*{Shade}$ is a shade map. Thus, in particular,
$\operatorname*{Shade}$ satisfies Axiom 2 in Definition
\ref{def.Shade.shademap}.

The set $B$ is nonempty (since $\left\vert B\right\vert =k+1>k\geq0$). Thus,
there exists some $u\in B$. Consider this $u$. If we had $u\in
\operatorname*{Shade}A$, then we would have $u\in B\cap\operatorname*{Shade}A$
(since $u\in B$ and $u\in\operatorname*{Shade}A$), which would entail
$B\cap\operatorname*{Shade}A\neq\varnothing$; but this would contradict
$B\cap\operatorname*{Shade}A=\varnothing$. Hence, we cannot have
$u\in\operatorname*{Shade}A$. In other words, we have $u\notin%
\operatorname*{Shade}A$. Hence, $u\in E\setminus\operatorname*{Shade}A$.

From $u\in B$, we obtain $\left\vert B\setminus\left\{  u\right\}  \right\vert
=\left\vert B\right\vert -1=k$ (since $\left\vert B\right\vert =k+1$).
Moreover, $\left(  B\setminus\left\{  u\right\}  \right)  \cap
\operatorname*{Shade}A=\varnothing$ (since $\underbrace{\left(  B\setminus
\left\{  u\right\}  \right)  }_{\subseteq B}\cap\operatorname*{Shade}%
A\subseteq B\cap\operatorname*{Shade}A=\varnothing$). Thus, we can apply
(\ref{pf.lem.Shade.FoG.IH}) to $B\setminus\left\{  u\right\}  $ instead of
$B$. As a consequence, we obtain $\operatorname*{Shade}\left(  A\setminus
\left(  B\setminus\left\{  u\right\}  \right)  \right)  =\operatorname*{Shade}%
A$. Thus, $u\in E\setminus\operatorname*{Shade}\left(  A\setminus\left(
B\setminus\left\{  u\right\}  \right)  \right)  $ (since $u\in E\setminus
\operatorname*{Shade}A$). Thus, Axiom 2 in Definition \ref{def.Shade.shademap}
(applied to $F=A\setminus\left(  B\setminus\left\{  u\right\}  \right)  $)
yields
\[
\operatorname*{Shade}\left(  \left(  A\setminus\left(  B\setminus\left\{
u\right\}  \right)  \right)  \setminus\left\{  u\right\}  \right)
=\operatorname*{Shade}\left(  A\setminus\left(  B\setminus\left\{  u\right\}
\right)  \right)  =\operatorname*{Shade}A.
\]
In view of
\[
\left(  A\setminus\left(  B\setminus\left\{  u\right\}  \right)  \right)
\setminus\left\{  u\right\}  =A\setminus\underbrace{\left(  \left(
B\setminus\left\{  u\right\}  \right)  \cup\left\{  u\right\}  \right)
}_{\substack{=B\\\text{(since }u\in B\text{)}}}=A\setminus B,
\]
this rewrites as $\operatorname*{Shade}\left(  A\setminus B\right)
=\operatorname*{Shade}A$.

Forget that we fixed $A$ and $B$. We thus have shown that if $A$ and $B$ are
two subsets of $E$ such that $B$ is finite and $B\cap\operatorname*{Shade}%
A=\varnothing$ and $\left\vert B\right\vert =k+1$, then $\operatorname*{Shade}%
\left(  A\setminus B\right)  =\operatorname*{Shade}A$. In other words, Lemma
\ref{lem.Shade.FoG} holds for $\left\vert B\right\vert =k+1$. This completes
the induction step. Thus, Lemma \ref{lem.Shade.FoG} is proven.
\end{proof}
\end{verlong}

We will also use another simple set-theoretical lemma:

\begin{lemma}
\label{lem.bip.Shade-to-Bip.lem1}Let $E$ be a set. Let $X$ and $Y$ be two
subsets of $E$. Then,%
\[
E\setminus Y=\left(  X\cup\left(  E\setminus Y\right)  \right)  \setminus
\left(  X\cap Y\right)  .
\]

\end{lemma}

\begin{proof}
[Proof of Lemma \ref{lem.bip.Shade-to-Bip.lem1}.]This is straightforward set
theory; it is easily checked (e.g.) using Venn diagrams.
\end{proof}

\begin{proof}
[Proof of Theorem \ref{thm.bip.Shade-to-Bip}.]\textbf{(a)} Let $F\in
\mathcal{P}\left(  E\right)  $. The definition of $\alpha$ yields
$\alpha\left(  F\right)  =F\cap\operatorname*{Shade}F\subseteq F$. The
definition of $\tau$ yields $\tau\left(  F\right)  =F\cup\left(
E\setminus\operatorname*{Shade}F\right)  \supseteq F$, so that $F\subseteq
\tau\left(  F\right)  $. Thus, we have the chain of inclusions $\alpha\left(
F\right)  \subseteq F\subseteq\tau\left(  F\right)  $. In other words,
$F\in\left[  \alpha\left(  F\right)  ,\ \tau\left(  F\right)  \right]  $. This
proves Theorem \ref{thm.bip.Shade-to-Bip} \textbf{(a)}.

\textbf{(b)} Let $F\in\mathcal{P}\left(  E\right)  $ and $G\in\left[
\alpha\left(  F\right)  ,\ \tau\left(  F\right)  \right]  $. From $G\in\left[
\alpha\left(  F\right)  ,\ \tau\left(  F\right)  \right]  $, we obtain the
chain of inclusions $\alpha\left(  F\right)  \subseteq G\subseteq\tau\left(
F\right)  $. Thus, $G\subseteq\tau\left(  F\right)  =F\cup\left(
E\setminus\operatorname*{Shade}F\right)  $ (by the definition of $\tau$).
Hence,
\[
\underbrace{G}_{\subseteq F\cup\left(  E\setminus\operatorname*{Shade}F\right)
}\setminus F\subseteq \left(  F\cup\left(  E\setminus\operatorname*{Shade}F\right)
\right)  \setminus F\subseteq E\setminus\operatorname*{Shade}F
\]
(since $\left(  X\cup Y\right)  \setminus X\subseteq Y$ for any two sets $X$
and $Y$). Therefore,%
\[
\underbrace{\left(  G\setminus F\right)  }_{\subseteq E\setminus
\operatorname*{Shade}F}\cap\operatorname*{Shade}F\subseteq\left(
E\setminus\operatorname*{Shade}F\right)  \cap\operatorname*{Shade}%
F=\varnothing
\]
(since $\left(  X\setminus Y\right)  \cap Y=\varnothing$ for any two sets $X$
and $Y$). Thus,%
\begin{equation}
\left(  G\setminus F\right)  \cap\operatorname*{Shade}F=\varnothing.
\label{pf.thm.bip.Shade-to-Bip.b.1}%
\end{equation}


Note that $G\setminus F\subseteq G\subseteq\tau\left(  F\right)  \subseteq E$
(since $\tau\left(  F\right)  \in\mathcal{P}\left(  E\right)  $) and $F\cap
G\subseteq F\subseteq E$ (since $F\in\mathcal{P}\left(  E\right)  $) and
$F\setminus G\subseteq F\subseteq E$.

Furthermore, we have $\alpha\left(  F\right)  =F\cap\operatorname*{Shade}F$
(by the definition of $\alpha$). Hence, $F\cap\operatorname*{Shade}%
F=\alpha\left(  F\right)  \subseteq G$. Thus, $G\supseteq F\cap
\operatorname*{Shade}F$, so that%
\[
F\setminus\underbrace{G}_{\supseteq F\cap\operatorname*{Shade}F}\subseteq
F\setminus\left(  F\cap\operatorname*{Shade}F\right)  =F\setminus
\operatorname*{Shade}F
\]
(since $X\setminus\left(  X\cap Y\right)  =X\setminus Y$ for any two sets $X$
and $Y$). Hence,%
\[
\underbrace{\left(  F\setminus G\right)  }_{\subseteq F\setminus
\operatorname*{Shade}F}\cap\operatorname*{Shade}F\subseteq\left(
F\setminus\operatorname*{Shade}F\right)  \cap\operatorname*{Shade}%
F=\varnothing
\]
(since $\left(  X\setminus Y\right)  \cap Y=\varnothing$ for any two sets $X$
and $Y$). Thus,%
\[
\left(  F\setminus G\right)  \cap\operatorname*{Shade}F=\varnothing.
\]
Hence, Lemma \ref{lem.Shade.FoG} (applied to $A=F$ and $B=F\setminus G$)
yields $\operatorname*{Shade}\left(  F\setminus\left(  F\setminus G\right)
\right)  =\operatorname*{Shade}F$. In view of $F\setminus\left(  F\setminus
G\right)  =F\cap G$, this rewrites as $\operatorname*{Shade}\left(  F\cap
G\right)  =\operatorname*{Shade}F$. Hence, (\ref{pf.thm.bip.Shade-to-Bip.b.1})
rewrites as
\[
\left(  G\setminus F\right)  \cap\operatorname*{Shade}\left(  F\cap G\right)
=\varnothing.
\]
Thus, Lemma \ref{lem.Shade.FuG} (applied to $A=F\cap G$ and $B=G\setminus F$)
yields \newline$\operatorname*{Shade}\left(  \left(  F\cap G\right)
\cup\left(  G\setminus F\right)  \right)  =\operatorname*{Shade}\left(  F\cap
G\right)  =\operatorname*{Shade}F$. In view of $\left(  F\cap G\right)
\cup\left(  G\setminus F\right)  =G$, this rewrites as $\operatorname*{Shade}%
G=\operatorname*{Shade}F$. In other words,
\[
\operatorname*{Shade}F=\operatorname*{Shade}G.
\]


The definition of $\alpha$ yields $\alpha\left(  F\right)  =F\cap
\operatorname*{Shade}F$ and $\alpha\left(  G\right)  =G\cap
\operatorname*{Shade}G$. Now, $\alpha\left(  F\right)  =F\cap
\operatorname*{Shade}F\subseteq\operatorname*{Shade}F$, so that
\[
\alpha\left(  F\right)  =\underbrace{\alpha\left(  F\right)  }_{\subseteq
G}\cap\underbrace{\operatorname*{Shade}F}_{=\operatorname*{Shade}G}\subseteq
G\cap\operatorname*{Shade}G=\alpha\left(  G\right)  .
\]
Combining this with%
\begin{align*}
\alpha\left(  G\right)   &  =\underbrace{G}_{\substack{\subseteq\tau\left(
F\right)  \\=F\cup\left(  E\setminus\operatorname*{Shade}F\right)  }%
}\cap\underbrace{\operatorname*{Shade}G}_{=\operatorname*{Shade}F}%
\subseteq\left(  F\cup\left(  E\setminus\operatorname*{Shade}F\right)
\right)  \cap\operatorname*{Shade}F\\
&  =\underbrace{\left(  F\cap\operatorname*{Shade}F\right)  }_{=\alpha\left(
F\right)  }\cup\underbrace{\left(  \left(  E\setminus\operatorname*{Shade}%
F\right)  \cap\operatorname*{Shade}F\right)  }_{=\varnothing}\\
&  \ \ \ \ \ \ \ \ \ \ \ \ \ \ \ \ \ \ \ \ \left(  \text{since }\left(  X\cup
Y\right)  \cap Z=\left(  X\cap Z\right)  \cup\left(  Y\cap Z\right)  \text{
for any three sets }X,Y,Z\right) \\
&  =\alpha\left(  F\right)  \cup\varnothing=\alpha\left(  F\right)  ,
\end{align*}
we obtain $\alpha\left(  F\right)  =\alpha\left(  G\right)  $.

The definition of $\tau$ yields $\tau\left(  F\right)  =F\cup\left(
E\setminus\operatorname*{Shade}F\right)  $ and $\tau\left(  G\right)
=G\cup\left(  E\setminus\operatorname*{Shade}G\right)  $. Now, $\tau\left(
F\right)  =F\cup\left(  E\setminus\operatorname*{Shade}F\right)  \supseteq
E\setminus\operatorname*{Shade}F$, so that%
\[
\tau\left(  F\right)  =\underbrace{\tau\left(  F\right)  }%
_{\substack{\supseteq G\\\text{(since }G\subseteq\tau\left(  F\right)
\text{)}}}\cup\left(  E\setminus\underbrace{\operatorname*{Shade}%
F}_{=\operatorname*{Shade}G}\right)  \supseteq G\cup\left(  E\setminus
\operatorname*{Shade}G\right)  =\tau\left(  G\right)  .
\]
Combining this with%
\begin{align*}
\tau\left(  F\right)   &  =\underbrace{F}_{\substack{=\left(  F\cap
\operatorname*{Shade}F\right)  \cup\left(  F\setminus\operatorname*{Shade}%
F\right)  \\\text{(since }X=\left(  X\cap Y\right)  \cup\left(  X\setminus
Y\right)  \\\text{for any two sets }X\text{ and }Y\text{)}}}\cup\left(
E\setminus\operatorname*{Shade}F\right) \\
&  =\left(  F\cap\operatorname*{Shade}F\right)  \cup\left(  \underbrace{F}%
_{\subseteq E}\setminus\operatorname*{Shade}F\right)  \cup\left(
E\setminus\operatorname*{Shade}F\right) \\
&  \subseteq\underbrace{\left(  F\cap\operatorname*{Shade}F\right)  }%
_{=\alpha\left(  F\right)  }\cup\underbrace{\left(  E\setminus
\operatorname*{Shade}F\right)  \cup\left(  E\setminus\operatorname*{Shade}%
F\right)  }_{=E\setminus\operatorname*{Shade}F}\\
&  =\underbrace{\alpha\left(  F\right)  }_{\subseteq G}\cup\left(
E\setminus\underbrace{\operatorname*{Shade}F}_{=\operatorname*{Shade}%
G}\right)  \subseteq G\cup\left(  E\setminus\operatorname*{Shade}G\right)
=\tau\left(  G\right)  ,
\end{align*}
we obtain $\tau\left(  F\right)  =\tau\left(  G\right)  $.

Now, $\left[  \underbrace{\alpha\left(  F\right)  }_{=\alpha\left(  G\right)
},\ \underbrace{\tau\left(  F\right)  }_{=\tau\left(  G\right)  }\right]
=\left[  \alpha\left(  G\right)  ,\ \tau\left(  G\right)  \right]  $. Theorem
\ref{thm.bip.Shade-to-Bip} \textbf{(b)} is thus proven.

\textbf{(c)} We have $\mathbf{P}=\left\{  \left[  \alpha\left(  F\right)
,\ \tau\left(  F\right)  \right]  \ \mid\ F\in\mathcal{P}\left(  E\right)
\right\}  $. Thus, $\mathbf{P}$ is a set of Boolean intervals of
$\mathcal{P}\left(  E\right)  $ (since $\left[  \alpha\left(  F\right)
,\ \tau\left(  F\right)  \right]  $ is a Boolean interval of $\mathcal{P}%
\left(  E\right)  $ whenever $F\in\mathcal{P}\left(  E\right)  $). We shall
refer to the elements of $\mathbf{P}$ as the \emph{blocks} of $\mathbf{P}$
(even though we don't know yet that $\mathbf{P}$ is a Boolean interval partition).

Next, we shall show that the blocks of $\mathbf{P}$ are pairwise disjoint.

[\textit{Proof.} Let $Q$ and $R$ be two distinct blocks of $\mathbf{P}$.

Let $G\in Q\cap R$. Now, we have $Q=\left[  \alpha\left(  F\right)
,\ \tau\left(  F\right)  \right]  $ for some $F\in\mathcal{P}\left(  E\right)
$ (since $Q\in\mathbf{P}=\left\{  \left[  \alpha\left(  F\right)
,\ \tau\left(  F\right)  \right]  \ \mid\ F\in\mathcal{P}\left(  E\right)
\right\}  $). Consider this $F$. We have $G\in Q\cap R\subseteq Q=\left[
\alpha\left(  F\right)  ,\ \tau\left(  F\right)  \right]  $. Therefore,
Theorem \ref{thm.bip.Shade-to-Bip} \textbf{(b)} yields $\left[  \alpha\left(
F\right)  ,\ \tau\left(  F\right)  \right]  =\left[  \alpha\left(  G\right)
,\ \tau\left(  G\right)  \right]  $. Hence, $Q=\left[  \alpha\left(  F\right)
,\ \tau\left(  F\right)  \right]  =\left[  \alpha\left(  G\right)
,\ \tau\left(  G\right)  \right]  $. The same argument (applied to $R$ instead
of $Q$) yields $R=\left[  \alpha\left(  G\right)  ,\ \tau\left(  G\right)
\right]  $. Comparing these two equalities, we obtain $Q=R$; this contradicts
the assumption that $Q$ and $R$ are distinct.

Forget that we fixed $G$. We thus have obtained a contradiction for each $G\in
Q\cap R$. Hence, there exists no $G\in Q\cap R$. In other words, $Q\cap
R=\varnothing$. In other words, $Q$ and $R$ are disjoint.

Forget that we fixed $Q$ and $R$. We thus have shown that any two distinct
blocks $Q$ and $R$ of $\mathbf{P}$ are disjoint. In other words, the blocks of
$\mathbf{P}$ are pairwise disjoint.]

Recall that $\mathbf{P}$ is a set of Boolean intervals of $\mathcal{P}\left(
E\right)  $. Hence, $\mathbf{P}$ is a set of pairwise disjoint Boolean
intervals of $\mathcal{P}\left(  E\right)  $ (since the blocks of $\mathbf{P}$
are pairwise disjoint).

Now, let $G\in\mathcal{P}\left(  E\right)  $. Then, $\left[  \alpha\left(
G\right)  ,\ \tau\left(  G\right)  \right]  \in\left\{  \left[  \alpha\left(
F\right)  ,\ \tau\left(  F\right)  \right]  \ \mid\ F\in\mathcal{P}\left(
E\right)  \right\}  =\mathbf{P}$. Thus, the set $\left[  \alpha\left(
G\right)  ,\ \tau\left(  G\right)  \right]  $ is one of the sets that are
being united in the union $\bigcup_{R\in\mathbf{P}}R$. Therefore, $\left[
\alpha\left(  G\right)  ,\ \tau\left(  G\right)  \right]  \subseteq
\bigcup_{R\in\mathbf{P}}R$. However, Theorem \ref{thm.bip.Shade-to-Bip}
\textbf{(a)} (applied to $F=G$) yields $G\in\left[  \alpha\left(  G\right)
,\ \tau\left(  G\right)  \right]  \subseteq\bigcup_{R\in\mathbf{P}}R$.

Forget that we fixed $G$. We thus have shown that $G\in\bigcup_{R\in
\mathbf{P}}R$ for each $G\in\mathcal{P}\left(  E\right)  $. In other words,
$\mathcal{P}\left(  E\right)  \subseteq\bigcup_{R\in\mathbf{P}}R$. Combining
this inclusion with the inclusion $\bigcup_{R\in\mathbf{P}}R\subseteq
\mathcal{P}\left(  E\right)  $ (which is obvious, since each $R\in\mathbf{P}$
is a subset of $\mathcal{P}\left(  E\right)  $), we obtain $\bigcup
_{R\in\mathbf{P}}R=\mathcal{P}\left(  E\right)  $. In other words, the union
of all blocks of $\mathbf{P}$ is $\mathcal{P}\left(  E\right)  $.

Now, recall that $\mathbf{P}$ is a set of pairwise disjoint Boolean intervals
of $\mathcal{P}\left(  E\right)  $. Hence, $\mathbf{P}$ is a set of pairwise
disjoint Boolean intervals of $\mathcal{P}\left(  E\right)  $ whose union is
$\mathcal{P}\left(  E\right)  $ (since the union of all blocks of $\mathbf{P}$
is $\mathcal{P}\left(  E\right)  $). In other words, $\mathbf{P}$ is a Boolean
interval partition of $\mathcal{P}\left(  E\right)  $. This proves Theorem
\ref{thm.bip.Shade-to-Bip} \textbf{(c)}.

\textbf{(d)} Let $F\in\mathcal{P}\left(  E\right)  $. Then,
$\operatorname*{Shade}F$ belongs to $\mathcal{P}\left(  E\right)  $ as well
(since $\operatorname*{Shade}$ is a map from $\mathcal{P}\left(  E\right)  $
to $\mathcal{P}\left(  E\right)  $). In other words, $\operatorname*{Shade}F$
is a subset of $E$.

Also, $F$ is a subset of $E$ (since $F\in\mathcal{P}\left(  E\right)  $).
Hence, Lemma \ref{lem.bip.Shade-to-Bip.lem1} (applied to $X=F$ and
$Y=\operatorname*{Shade}F$) yields%
\[
E\setminus\operatorname*{Shade}F=\underbrace{\left(  F\cup\left(
E\setminus\operatorname*{Shade}F\right)  \right)  }_{\substack{=\tau\left(
F\right)  \\\text{(since }\tau\left(  F\right)  =F\cup\left(  E\setminus
\operatorname*{Shade}F\right)  \text{)}}}\setminus\underbrace{\left(
F\cap\operatorname*{Shade}F\right)  }_{\substack{=\alpha\left(  F\right)
\\\text{(since }\alpha\left(  F\right)  =F\cap\operatorname*{Shade}F\text{)}%
}}=\tau\left(  F\right)  \setminus\alpha\left(  F\right)  .
\]
However, $\operatorname*{Shade}F$ is a subset of $E$; thus, $E\setminus\left(
E\setminus\operatorname*{Shade}F\right)  =\operatorname*{Shade}F$. Therefore,%
\[
\operatorname*{Shade}F=E\setminus\underbrace{\left(  E\setminus
\operatorname*{Shade}F\right)  }_{=\tau\left(  F\right)  \setminus
\alpha\left(  F\right)  }=E\setminus\left(  \tau\left(  F\right)
\setminus\alpha\left(  F\right)  \right)  .
\]
This proves Theorem \ref{thm.bip.Shade-to-Bip} \textbf{(d)}.
\end{proof}

Combining Theorem \ref{thm.bip.Bip-to-Shade} with Theorem
\ref{thm.bip.Shade-to-Bip}, we obtain the following:

\begin{theorem}
\label{thm.bip.cryptomor}Let $E$ be a finite set. Then, there is a bijection
from the set%
\[
\left\{  \text{shade maps }\operatorname*{Shade}:\mathcal{P}\left(  E\right)
\rightarrow\mathcal{P}\left(  E\right)  \right\}
\]
to the set%
\[
\left\{  \text{Boolean interval partitions of }\mathcal{P}\left(  E\right)
\right\}  .
\]
It sends each map $\operatorname*{Shade}$ to the Boolean interval partition
$\mathbf{P}$ defined in Theorem \ref{thm.bip.Shade-to-Bip}\textbf{ (c)}. Its
inverse map sends each Boolean interval partition $\mathbf{P}$ to the map
$\operatorname*{Shade}$ defined in Theorem \ref{thm.bip.Bip-to-Shade}.
\end{theorem}

\begin{vershort}
\begin{proof}
[Proof of Theorem \ref{thm.bip.cryptomor}.]This follows easily from Theorem
\ref{thm.bip.Bip-to-Shade} and Theorem \ref{thm.bip.Shade-to-Bip}.
\end{proof}
\end{vershort}

\begin{verlong}
\begin{proof}
[Proof of Theorem \ref{thm.bip.cryptomor}.]If $\operatorname*{Shade}%
:\mathcal{P}\left(  E\right)  \rightarrow\mathcal{P}\left(  E\right)  $ is a
shade map, then the $\mathbf{P}$ defined in Theorem \ref{thm.bip.Shade-to-Bip}
\textbf{(c)} is a Boolean interval partition of $\mathcal{P}\left(  E\right)
$ (by Theorem \ref{thm.bip.Shade-to-Bip} \textbf{(c)}). Hence, we can define a
map%
\[
\mathbf{B}:\left\{  \text{shade maps }\operatorname*{Shade}:\mathcal{P}\left(
E\right)  \rightarrow\mathcal{P}\left(  E\right)  \right\}  \rightarrow
\left\{  \text{Boolean interval partitions of }\mathcal{P}\left(  E\right)
\right\}
\]
that sends each map $\operatorname*{Shade}$ to the Boolean interval partition
$\mathbf{P}$ defined in Theorem \ref{thm.bip.Shade-to-Bip}\textbf{ (c)}.
Consider this map $\mathbf{B}$.

If $\mathbf{P}$ is a Boolean interval partition of $\mathcal{P}\left(
E\right)  $, then the map $\operatorname*{Shade}:\mathcal{P}\left(  E\right)
\rightarrow\mathcal{P}\left(  E\right)  $ defined in Theorem
\ref{thm.bip.Bip-to-Shade} is a shade map (by Theorem
\ref{thm.bip.Bip-to-Shade} \textbf{(a)}). Hence, we can define a map%
\[
\mathbf{S}:\left\{  \text{Boolean interval partitions of }\mathcal{P}\left(
E\right)  \right\}  \rightarrow\left\{  \text{shade maps }%
\operatorname*{Shade}:\mathcal{P}\left(  E\right)  \rightarrow\mathcal{P}%
\left(  E\right)  \right\}
\]
that sends each Boolean interval partition $\mathbf{P}$ to the map
$\operatorname*{Shade}$ defined in Theorem \ref{thm.bip.Bip-to-Shade}.
Consider this map $\mathbf{S}$.

We shall now show that the maps $\mathbf{B}$ and $\mathbf{S}$ are mutually
inverse. Indeed, we claim the following:

\begin{statement}
\textit{Claim 1:} We have $\mathbf{B}\circ\mathbf{S}=\operatorname*{id}$.
\end{statement}

[\textit{Proof of Claim 1:} Let $\mathbf{P}$ be a Boolean interval partition
of $\mathcal{P}\left(  E\right)  $.

For any $F\in\mathcal{P}\left(  E\right)  $, let $\left[  \alpha\left(
F\right)  ,\ \tau\left(  F\right)  \right]  $ denote the (unique) block of
$\mathbf{P}$ that contains $F$. Thus, we have defined two sets $\alpha\left(
F\right)  \in\mathcal{P}\left(  E\right)  $ and $\tau\left(  F\right)
\in\mathcal{P}\left(  E\right)  $ for each $F\in\mathcal{P}\left(  E\right)
$. In other words, we have defined two maps $\alpha:\mathcal{P}\left(
E\right)  \rightarrow\mathcal{P}\left(  E\right)  $ and $\tau:\mathcal{P}%
\left(  E\right)  \rightarrow\mathcal{P}\left(  E\right)  $.

We define a map $\operatorname*{Shade}:\mathcal{P}\left(  E\right)
\rightarrow\mathcal{P}\left(  E\right)  $ by setting%
\[
\operatorname*{Shade}F=E\setminus\left(  \tau\left(  F\right)  \setminus
\alpha\left(  F\right)  \right)  \ \ \ \ \ \ \ \ \ \ \text{for any }%
F\in\mathcal{P}\left(  E\right)  .
\]
This map $\operatorname*{Shade}$ is precisely the map $\operatorname*{Shade}$
defined in Theorem \ref{thm.bip.Bip-to-Shade}. Hence, $\mathbf{S}\left(
\mathbf{P}\right)  =\operatorname*{Shade}$ (by the definition of $\mathbf{S}$).

Now, any $F\in\mathcal{P}\left(  E\right)  $ satisfies $\alpha\left(
F\right)  =F\cap\operatorname*{Shade}F$ (by Theorem \ref{thm.bip.Bip-to-Shade}
\textbf{(b)}). Thus, our map $\alpha:\mathcal{P}\left(  E\right)
\rightarrow\mathcal{P}\left(  E\right)  $ is precisely the map $\alpha$
defined in Theorem \ref{thm.bip.Shade-to-Bip}.

Furthermore, any $F\in\mathcal{P}\left(  E\right)  $ satisfies $\tau\left(
F\right)  =F\cup\left(  E\setminus\operatorname*{Shade}F\right)  $ (by Theorem
\ref{thm.bip.Bip-to-Shade} \textbf{(b)}). Thus, our map $\tau:\mathcal{P}%
\left(  E\right)  \rightarrow\mathcal{P}\left(  E\right)  $ is precisely the
map $\tau$ defined in Theorem \ref{thm.bip.Shade-to-Bip}.

Moreover, we have $\mathbf{P}=\left\{  \left[  \alpha\left(  F\right)
,\ \tau\left(  F\right)  \right]  \ \mid\ F\in\mathcal{P}\left(  E\right)
\right\}  $ (by Theorem \ref{thm.bip.Bip-to-Shade} \textbf{(c)}). Thus, our
Boolean interval partition $\mathbf{P}$ is precisely the Boolean interval
partition $\mathbf{P}$ defined in Theorem \ref{thm.bip.Shade-to-Bip}
\textbf{(c)} (because we already know that our maps $\alpha$ and $\tau$ are
precisely the maps $\alpha$ and $\tau$ defined in Theorem
\ref{thm.bip.Shade-to-Bip}). However, the definition of $\mathbf{B}$ shows
that $\mathbf{B}\left(  \operatorname*{Shade}\right)  $ is the Boolean
interval partition $\mathbf{P}$ defined in Theorem \ref{thm.bip.Shade-to-Bip}
\textbf{(c)}. Thus, $\mathbf{B}\left(  \operatorname*{Shade}\right)
=\mathbf{P}$ (since our Boolean interval partition $\mathbf{P}$ is precisely
the Boolean interval partition $\mathbf{P}$ defined in Theorem
\ref{thm.bip.Shade-to-Bip} \textbf{(c)}). Now,%
\[
\left(  \mathbf{B}\circ\mathbf{S}\right)  \left(  \mathbf{P}\right)
=\mathbf{B}\left(  \underbrace{\mathbf{S}\left(  \mathbf{P}\right)
}_{=\operatorname*{Shade}}\right)  =\mathbf{B}\left(  \operatorname*{Shade}%
\right)  =\mathbf{P}=\operatorname*{id}\left(  \mathbf{P}\right)  .
\]


Now, forget that we fixed $\mathbf{P}$. We thus have shown that $\left(
\mathbf{B}\circ\mathbf{S}\right)  \left(  \mathbf{P}\right)
=\operatorname*{id}\left(  \mathbf{P}\right)  $ for any Boolean interval
partition $\mathbf{P}$ of $\mathcal{P}\left(  E\right)  $. In other words,
$\mathbf{B}\circ\mathbf{S}=\operatorname*{id}$. This proves Claim 1.]

\begin{statement}
\textit{Claim 2:} We have $\mathbf{S}\circ\mathbf{B}=\operatorname*{id}$.
\end{statement}

[\textit{Proof of Claim 2:} Let $\operatorname*{Shade}:\mathcal{P}\left(
E\right)  \rightarrow\mathcal{P}\left(  E\right)  $ be a shade map on $E$.
Define two maps $\alpha$ and $\tau$ as in Theorem \ref{thm.bip.Shade-to-Bip}.
Define a Boolean interval partition $\mathbf{P}$ as in Theorem
\ref{thm.bip.Shade-to-Bip} \textbf{(c)}. Then, $\mathbf{B}\left(
\operatorname*{Shade}\right)  =\mathbf{P}$ (by the definition of $\mathbf{B}$).

Let $F\in\mathcal{P}\left(  E\right)  $. Then, $\left[  \alpha\left(
F\right)  ,\ \tau\left(  F\right)  \right]  \in\mathbf{P}$ (by
(\ref{eq.thm.bip.Shade-to-Bip.P=})). In other words, $\left[  \alpha\left(
F\right)  ,\ \tau\left(  F\right)  \right]  $ is a block of the Boolean
interval partition $\mathbf{P}$. Furthermore, this block $\left[
\alpha\left(  F\right)  ,\ \tau\left(  F\right)  \right]  $ contains $F$
(since Theorem \ref{thm.bip.Shade-to-Bip} \textbf{(a)} yields $F\in\left[
\alpha\left(  F\right)  ,\ \tau\left(  F\right)  \right]  $).

However, there is a unique block of $\mathbf{P}$ that contains $F$ (since
$\mathbf{P}$ is a Boolean interval partition of $\mathcal{P}\left(  E\right)
$). This unique block must be $\left[  \alpha\left(  F\right)  ,\ \tau\left(
F\right)  \right]  $ (since $\left[  \alpha\left(  F\right)  ,\ \tau\left(
F\right)  \right]  $ is a block of the Boolean interval partition $\mathbf{P}$
and contains $F$). Thus, $\left[  \alpha\left(  F\right)  ,\ \tau\left(
F\right)  \right]  $ is the (unique) block of $\mathbf{P}$ that contains $F$.

Forget that we fixed $F$. We thus have shown that if $F\in\mathcal{P}\left(
E\right)  $, then $\left[  \alpha\left(  F\right)  ,\ \tau\left(  F\right)
\right]  $ is the (unique) block of $\mathbf{P}$ that contains $F$. In other
words, our definition of $\alpha\left(  F\right)  $ and $\tau\left(  F\right)
$ agrees with the definition of $\alpha\left(  F\right)  $ and $\tau\left(
F\right)  $ given in Theorem \ref{thm.bip.Bip-to-Shade} (because the latter
definition defines $\alpha\left(  F\right)  $ and $\tau\left(  F\right)  $ by
requiring that $\left[  \alpha\left(  F\right)  ,\ \tau\left(  F\right)
\right]  $ be the (unique) block of $\mathbf{P}$ that contains $F$). Hence,
everything said about $\alpha\left(  F\right)  $ and $\tau\left(  F\right)  $
in Theorem \ref{thm.bip.Bip-to-Shade} can be applied to our $\alpha\left(
F\right)  $ and $\tau\left(  F\right)  $ as well.

Moreover, Theorem \ref{thm.bip.Shade-to-Bip} \textbf{(d)} shows that we have
$\operatorname*{Shade}F=E\setminus\left(  \tau\left(  F\right)  \setminus
\alpha\left(  F\right)  \right)  $ for any $F\in\mathcal{P}\left(  E\right)
$. Hence, our map $\operatorname*{Shade}$ is precisely the map
$\operatorname*{Shade}$ defined in Theorem \ref{thm.bip.Bip-to-Shade} (because
we already know that our definition of $\alpha\left(  F\right)  $ and
$\tau\left(  F\right)  $ agrees with the definition of $\alpha\left(
F\right)  $ and $\tau\left(  F\right)  $ given in Theorem
\ref{thm.bip.Bip-to-Shade}). However, the definition of $\mathbf{S}$ shows
that $\mathbf{S}\left(  \mathbf{P}\right)  $ is the map $\operatorname*{Shade}%
$ defined in Theorem \ref{thm.bip.Bip-to-Shade}. Thus, $\mathbf{S}\left(
\mathbf{P}\right)  =\operatorname*{Shade}$ (since our map
$\operatorname*{Shade}$ is precisely the map $\operatorname*{Shade}$ defined
in Theorem \ref{thm.bip.Bip-to-Shade}). Now,%
\[
\left(  \mathbf{S}\circ\mathbf{B}\right)  \left(  \operatorname*{Shade}%
\right)  =\mathbf{S}\left(  \underbrace{\mathbf{B}\left(
\operatorname*{Shade}\right)  }_{=\mathbf{P}}\right)  =\mathbf{S}\left(
\mathbf{P}\right)  =\operatorname*{Shade}=\operatorname*{id}\left(
\operatorname*{Shade}\right)  .
\]


Now, forget that we fixed $\operatorname*{Shade}$. We thus have shown that
$\left(  \mathbf{S}\circ\mathbf{B}\right)  \left(  \operatorname*{Shade}%
\right)  =\operatorname*{id}\left(  \operatorname*{Shade}\right)  $ for any
shade map $\operatorname*{Shade}:\mathcal{P}\left(  E\right)  \rightarrow
\mathcal{P}\left(  E\right)  $. In other words, $\mathbf{S}\circ
\mathbf{B}=\operatorname*{id}$. This proves Claim 2.]

Combining Claim 1 with Claim 2, we conclude that the maps $\mathbf{B}$ and
$\mathbf{S}$ are mutually inverse. Hence, these two maps $\mathbf{B}$ and
$\mathbf{S}$ are bijections. Thus, there is a bijection from the set%
\[
\left\{  \text{shade maps }\operatorname*{Shade}:\mathcal{P}\left(  E\right)
\rightarrow\mathcal{P}\left(  E\right)  \right\}
\]
to the set%
\[
\left\{  \text{Boolean interval partitions of }\mathcal{P}\left(  E\right)
\right\}
\]
(namely, the map $\mathbf{B}$), and this bijection sends each map
$\operatorname*{Shade}$ to the Boolean interval partition $\mathbf{P}$ defined
in Theorem \ref{thm.bip.Shade-to-Bip}\textbf{ (c)} (because this is how
$\mathbf{B}$ was defined). Moreover, its inverse map is the map $\mathbf{S}$
(because $\mathbf{B}$ and $\mathbf{S}$ are mutually inverse), and therefore
sends each Boolean interval partition $\mathbf{P}$ to the map
$\operatorname*{Shade}$ defined in Theorem \ref{thm.bip.Bip-to-Shade} (because
this is how $\mathbf{S}$ was defined). Thus, the proof of Theorem
\ref{thm.bip.cryptomor} is complete.
\end{proof}
\end{verlong}

\begin{question}
According to Theorem \ref{thm.antimat.Shade}, any antimatroid gives rise to an
inclusion-reversing shade map, which in turn gives rise to a Boolean interval
partition by Theorem \ref{thm.bip.Shade-to-Bip}. Is this construction
equivalent to the construction of a Boolean interval partition from an
antimatroid described by Gordon and McMahon in \cite[Theorem 2.5]{GorMah97}?
\end{question}

\section{The topological viewpoint}

Now we return to the setting of Section \ref{sec.elser}. We aim to reinterpret
Theorem \ref{thm.elser-Shade0} in terms of combinatorial topology
(specifically, finite simplicial complexes) and strengthen it. We recall the
definition of a \emph{simplicial complex}:\footnote{We forget all the
conventions we have introduced so far. (Thus, for example, $E$ no longer means
the edge set of a graph $\Gamma$.)}

\begin{definition}
Let $E$ be a finite set. A \emph{simplicial complex} on ground set $E$ means a
subset $\mathcal{A}$ of the power set of $E$ with the following property:

\begin{statement}
If $P\in\mathcal{A}$ and $Q\subseteq P$, then $Q\in\mathcal{A}$.
\end{statement}
\end{definition}

Thus, in terms of posets, a simplicial complex on ground set $E$ means a
down-closed subset of the Boolean lattice on $E$. Note that a simplicial
complex contains the empty set $\varnothing$ unless it is empty itself.

We refer to \cite{Kozlov20} for context and theory about simplicial complexes.
We shall restrict ourselves to the few definitions relevant to what we will
prove. The following is fairly simple:

\begin{proposition}
\label{prop.elser-sc}Let us use the notations from Section \ref{sec.elser} as
well as Definition \ref{def.Shade}. Let $G$ be any subset of $E$. Let%
\begin{equation}
\mathcal{A}=\left\{  F\subseteq E\ \mid\ G\not \subseteq \operatorname*{Shade}%
F\right\}  . \label{eq.prop.elser-sc.A=}%
\end{equation}
Then, $\mathcal{A}$ is a simplicial complex on ground set $E$.
\end{proposition}

\begin{proof}
[Proof of Proposition \ref{prop.elser-sc}.]Clearly, $\mathcal{A}$ is a subset
of the power set of $E$. Thus, we only need to verify the following claim:

\begin{statement}
\textit{Claim 1:} If $P\in\mathcal{A}$ and $Q\subseteq P$, then $Q\in
\mathcal{A}$.
\end{statement}

\begin{vershort}
[\textit{Proof of Claim 1:} Let $P\in\mathcal{A}$ and let $Q\subseteq P$. We
must show that $Q\in\mathcal{A}$.

We have $P\in\mathcal{A}=\left\{  F\subseteq E\ \mid\ G\not \subseteq
\operatorname*{Shade}F\right\}  $. In other words, $P\subseteq E$ and
$G\not \subseteq \operatorname*{Shade}P$. But $Q\subseteq P$ and thus
$\operatorname*{Shade}Q\subseteq\operatorname*{Shade}P$ (by Lemma
\ref{lem.Shade-monoton}). Hence, from $G\not \subseteq \operatorname*{Shade}%
P$, we obtain $G\not \subseteq \operatorname*{Shade}Q$. Thus, $Q\in\left\{
F\subseteq E\ \mid\ G\not \subseteq \operatorname*{Shade}F\right\}  $. This
can be rewritten as $Q\in\mathcal{A}$ (by (\ref{eq.prop.elser-sc.A=})). Thus,
Claim 1 is proved.] \qedhere

\end{vershort}

\begin{verlong}
[\textit{Proof of Claim 1:} Let $P\in\mathcal{A}$ and let $Q\subseteq P$. We
must show that $Q\in\mathcal{A}$.

We have $P\in\mathcal{A}=\left\{  F\subseteq E\ \mid\ G\not \subseteq
\operatorname*{Shade}F\right\}  $. In other words, $P$ is a subset $F$ of $E$
satisfying $G\not \subseteq \operatorname*{Shade}F$. In other words, $P$ is a
subset of $E$ and satisfies $G\not \subseteq \operatorname*{Shade}P$. Now, $Q$
is a subset of $E$ (since $Q\subseteq P\subseteq E$). Also, from $Q\subseteq
P$, we obtain $\operatorname*{Shade}Q\subseteq\operatorname*{Shade}P$ (by
Lemma \ref{lem.Shade-monoton}, applied to $A=Q$ and $B=P$). Hence, from
$G\not \subseteq \operatorname*{Shade}P$, we obtain $G\not \subseteq
\operatorname*{Shade}Q$. Thus, $Q$ is a subset $F$ of $E$ satisfying
$G\not \subseteq \operatorname*{Shade}F$. In other words, $Q\in\left\{
F\subseteq E\ \mid\ G\not \subseteq \operatorname*{Shade}F\right\}  $. This
can be rewritten as $Q\in\mathcal{A}$ (by (\ref{eq.prop.elser-sc.A=})). Thus,
Claim 1 is proved.]

Claim 1 shows that $\mathcal{A}$ is a simplicial complex on ground set $E$.
This proves Proposition \ref{prop.elser-sc}.
\end{verlong}
\end{proof}

To state the main result of this section, we need a few more notions:

\begin{definition}
\label{def.sets.prec}Let $A$ and $B$ be two sets. Then, we say that $A\prec B$
if we have $B=A\cup\left\{  b\right\}  $ for some $b\in B\setminus A$.
\end{definition}

Equivalently, two sets $A$ and $B$ satisfy $A\prec B$ if and only if
$A\subseteq B$ and $\left\vert B\setminus A\right\vert =1$.

\begin{definition}
\label{def.matchings}Let $E$ be a finite set. Let $\mathcal{A}$ be a
simplicial complex on ground set $E$.

\textbf{(a)} A \emph{complete matching} of $\mathcal{A}$ means a triple
$\left(  \mathcal{A}_{-},\mathcal{A}_{+},\Phi\right)  $, where $\mathcal{A}%
_{-}$ and $\mathcal{A}_{+}$ are two disjoint subsets of $\mathcal{A}$
satisfying $\mathcal{A}_{-}\cup\mathcal{A}_{+}=\mathcal{A}$, and where
$\Phi:\mathcal{A}_{+}\rightarrow\mathcal{A}_{-}$ is a bijection with the
property that
\begin{equation}
\text{each }F\in\mathcal{A}_{+}\text{ satisfies }\Phi\left(  F\right)  \prec
F. \label{eq.def.matchings.a.cond}%
\end{equation}


\textbf{(b)} A complete matching $\left(  \mathcal{A}_{-},\mathcal{A}_{+}%
,\Phi\right)  $ of $\mathcal{A}$ is said to be \emph{acyclic} if there exists
no tuple $\left(  B_{1},B_{2},\ldots,B_{n}\right)  $ of distinct sets
$B_{1},B_{2},\ldots,B_{n}\in\mathcal{A}_{+}$ with the property that $n\geq2$
and that%
\[
\Phi\left(  B_{i}\right)  \prec B_{i+1}\ \ \ \ \ \ \ \ \ \ \text{for each
}i\in\left\{  1,2,\ldots,n-1\right\}
\]
and%
\[
\Phi\left(  B_{n}\right)  \prec B_{1}.
\]


\textbf{(c)} The simplicial complex $\mathcal{A}$ is said to be
\emph{collapsible} if it has an acyclic complete matching.
\end{definition}

These definitions are essentially equivalent to the definitions in
\cite{Kozlov20}, although it takes a bit of work to match them up precisely.
Our notion of a \textquotedblleft complete matching\textquotedblright\ as
defined in Definition \ref{def.matchings} \textbf{(a)} is a particular case of
the notion introduced in \cite[Chapter 10]{Kozlov20}, as we restrict ourselves
to simplicial complexes (i.e., down-closed sets of Boolean lattices) instead
of arbitrary posets. To be fully precise, our complete matchings are triples
$\left(  \mathcal{A}_{-},\mathcal{A}_{+},\Phi\right)  $, whereas the complete
matchings of \cite[Chapter 10]{Kozlov20} are certain fixed-point-free
involutions\footnote{A \emph{fixed-point-free involution} means an involution
(i.e., a map that is inverse to itself) that has no fixed point.}
$\mu:\mathcal{A}\rightarrow\mathcal{A}$; the equivalence between these two
objects is fairly easy to see (in particular, if $\left(  \mathcal{A}%
_{-},\mathcal{A}_{+},\Phi\right)  $ is a complete matching in our sense, then
the corresponding complete matching $\mu:\mathcal{A}\rightarrow\mathcal{A}$ in
the sense of \cite[Chapter 10]{Kozlov20} is the map that sends each
$B\in\mathcal{A}_{+}$ to $\Phi\left(  B\right)  \in\mathcal{A}_{-}$ and sends
each $B\in\mathcal{A}_{-}$ to $\Phi^{-1}\left(  B\right)  \in\mathcal{A}_{+}%
$). Our notion of \textquotedblleft collapsible\textquotedblright\ as defined
in Definition \ref{def.matchings} \textbf{(c)} is equivalent to the classical
notion of \textquotedblleft collapsible\textquotedblright\ (even though the
latter is usually defined differently) because of \cite[Theorem 10.9]%
{Kozlov20}.

We now claim:

\begin{theorem}
\label{thm.elser-coll}Let us use the notations from Section \ref{sec.elser} as
well as Definition \ref{def.Shade}. Let $G$ be any subset of $E$. Define
$\mathcal{A}$ as in (\ref{eq.prop.elser-sc.A=}). Then, the simplicial complex
$\mathcal{A}$ is collapsible.
\end{theorem}

Collapsible simplicial complexes are well-behaved in various ways -- in
particular, they are contractible (\cite[Corollary 9.19]{Kozlov20}), and thus
have trivial homotopy and homology groups (in positive degrees). Moreover, the
reduced Euler characteristic of any collapsible simplicial complex is $0$ (for
obvious reasons: having a complete matching suffices, even if it is not
acyclic); thus, Theorem \ref{thm.elser-Shade0} follows from Theorem
\ref{thm.elser-coll}.

Our proof of Theorem \ref{thm.elser-coll} will rely on the following simple lemma:

\begin{lemma}
\label{lem.elser-coll.lem}Let $X$ and $Y$ be two sets. Let $u\in X\cap Y$. If
$X\setminus\left\{  u\right\}  \prec Y$, then $X=Y$.
\end{lemma}

\begin{proof}
[Proof of Lemma \ref{lem.elser-coll.lem}.]Set $X^{\prime}=X\setminus\left\{
u\right\}  $. Thus, $\underbrace{X^{\prime}}_{=X\setminus\left\{  u\right\}
}\cup\left\{  u\right\}  =\left(  X\setminus\left\{  u\right\}  \right)
\cup\left\{  u\right\}  =X$ (since $u\in X\cap Y\subseteq X$). Also, clearly,
$u\notin X\setminus\left\{  u\right\}  $ (since $u\notin P\setminus\left\{
u\right\}  $ for any set $P$). In other words, $u\notin X^{\prime}$ (since
$X^{\prime}=X\setminus\left\{  u\right\}  $).

Now, assume that $X\setminus\left\{  u\right\}  \prec Y$. In other words,
$X^{\prime}\prec Y$ (since $X^{\prime}=X\setminus\left\{  u\right\}  $). In
other words, we have
\begin{equation}
Y=X^{\prime}\cup\left\{  b\right\}  \label{pf.lem.elser-coll.lem.1}%
\end{equation}
for some $b\in Y\setminus X^{\prime}$ (by Definition \ref{def.sets.prec}).
Consider this $b$. Combining $u\in X\cap Y\subseteq Y=X^{\prime}\cup\left\{
b\right\}  $ with $u\notin X^{\prime}$, we obtain%
\[
u\in\left(  X^{\prime}\cup\left\{  b\right\}  \right)  \setminus X^{\prime
}\subseteq\left\{  b\right\}  .
\]
In other words, $u=b$. Thus, (\ref{pf.lem.elser-coll.lem.1}) can be rewritten
as $Y=X^{\prime}\cup\left\{  u\right\}  $. Hence, $Y=X^{\prime}\cup\left\{
u\right\}  =X$, so that $X=Y$. This proves Lemma \ref{lem.elser-coll.lem}.
\end{proof}

\begin{proof}
[Proof of Theorem \ref{thm.elser-coll}.]We know from Proposition
\ref{prop.elser-sc} that $\mathcal{A}$ is a simplicial complex. It remains to
show that $\mathcal{A}$ is collapsible.

We have%
\[
\mathcal{A}=\left\{  F\subseteq E\ \mid\ G\not \subseteq \operatorname*{Shade}%
F\right\}  =\left\{  P\subseteq E\ \mid\ G\not \subseteq \operatorname*{Shade}%
P\right\}
\]
(here, we have renamed the index $F$ as $P$). Thus, our set $\mathcal{A}$ is
precisely the set $\mathcal{A}$ defined in the proof of Theorem
\ref{thm.elser-Shade0} above.

We equip the finite set $E$ with a total order (chosen arbitrarily, but fixed henceforth).

If $F\in\mathcal{A}$, then we define the edge $\varepsilon\left(  F\right)
\in G\setminus\operatorname*{Shade}F$ as in the proof of Theorem
\ref{thm.elser-Shade0}. That is, we define $\varepsilon\left(  F\right)  $ as
the \textbf{smallest} edge $e\in G\setminus\operatorname*{Shade}F$.

If two sets $F_{1}\in\mathcal{A}$ and $F_{2}\in\mathcal{A}$ satisfy
$\operatorname*{Shade}\left(  F_{1}\right)  =\operatorname*{Shade}\left(
F_{2}\right)  $, then
\begin{equation}
\varepsilon\left(  F_{1}\right)  =\varepsilon\left(  F_{2}\right)  .
\label{pf.thm.elser-coll.eps=eps}%
\end{equation}


\begin{vershort}
\noindent(Indeed, this is precisely the equality
(\ref{pf.thm.elser-Shade0.short.eps=eps}) from the above proof of Theorem
\ref{thm.elser-Shade0}.)
\end{vershort}

\begin{verlong}
\noindent(Indeed, this is precisely the equality
(\ref{pf.thm.elser-Shade0.eps=eps}) from the above proof of Theorem
\ref{thm.elser-Shade0}.)
\end{verlong}

We define two subsets $\mathcal{A}_{+}$ and $\mathcal{A}_{-}$ of $\mathcal{A}$
as in the proof of Theorem \ref{thm.elser-Shade0}. That is, we set%
\[
\mathcal{A}_{+}=\left\{  P\in\mathcal{A}\ \mid\ \varepsilon\left(  P\right)
\in P\right\}  \ \ \ \ \ \ \ \ \ \ \text{and}\ \ \ \ \ \ \ \ \ \ \mathcal{A}%
_{-}=\left\{  P\in\mathcal{A}\ \mid\ \varepsilon\left(  P\right)  \notin
P\right\}  .
\]
Thus, each $P\in\mathcal{A}$ satisfies either $P\in\mathcal{A}_{-}$ or
$P\in\mathcal{A}_{+}$ but not both at the same time (since it satisfies either
$\varepsilon\left(  P\right)  \notin P$ or $\varepsilon\left(  P\right)  \in
P$ but not both at the same time). Hence, $\mathcal{A}_{-}$ and $\mathcal{A}%
_{+}$ are two disjoint subsets of $\mathcal{A}$ satisfying $\mathcal{A}%
_{-}\cup\mathcal{A}_{+}=\mathcal{A}$.

We define a map $\Phi:\mathcal{A}_{+}\rightarrow\mathcal{A}_{-}$ as in the
proof of Theorem \ref{thm.elser-Shade0}. That is, we set%
\[
\Phi\left(  F\right)  =F\setminus\left\{  \varepsilon\left(  F\right)
\right\}  \ \ \ \ \ \ \ \ \ \ \text{for each }F\in\mathcal{A}_{+}.
\]
We know (from the proof of Theorem \ref{thm.elser-Shade0}) that the map $\Phi$
is a bijection. Moreover, it is clear that each $F\in\mathcal{A}_{+}$
satisfies $\Phi\left(  F\right)  \prec F$\ \ \ \ \footnote{\textit{Proof.} Let
$F\in\mathcal{A}_{+}$. Thus, $F\in\mathcal{A}_{+}=\left\{  P\in\mathcal{A}%
\ \mid\ \varepsilon\left(  P\right)  \in P\right\}  $. In other words, $F$ is
a $P\in\mathcal{A}$ satisfying $\varepsilon\left(  P\right)  \in P$. In other
words, $F\in\mathcal{A}$ and $\varepsilon\left(  F\right)  \in F$. From
$\varepsilon\left(  F\right)  \in F$, we obtain $F=\left(  F\setminus\left\{
\varepsilon\left(  F\right)  \right\}  \right)  \cup\left\{  \varepsilon
\left(  F\right)  \right\}  $ and $\varepsilon\left(  F\right)  \in
F\setminus\left(  F\setminus\left\{  \varepsilon\left(  F\right)  \right\}
\right)  $. Hence, $F\setminus\left\{  \varepsilon\left(  F\right)  \right\}
\prec F$ (by Definition \ref{def.sets.prec}). In other words, $\Phi\left(
F\right)  \prec F$ (since the definition of $\Phi$ yields $\Phi\left(
F\right)  =F\setminus\left\{  \varepsilon\left(  F\right)  \right\}  $).
Qed.}. Hence, the triple $\left(  \mathcal{A}_{-},\mathcal{A}_{+},\Phi\right)
$ is a complete matching of $\mathcal{A}$.

We shall now prove that this complete matching $\left(  \mathcal{A}%
_{-},\mathcal{A}_{+},\Phi\right)  $ is acyclic. Indeed, let $\left(
B_{1},B_{2},\ldots,B_{n}\right)  $ be a tuple of distinct sets $B_{1}%
,B_{2},\ldots,B_{n}\in\mathcal{A}_{+}$ with the property that $n\geq2$ and
that%
\begin{equation}
\Phi\left(  B_{i}\right)  \prec B_{i+1}\ \ \ \ \ \ \ \ \ \ \text{for each
}i\in\left\{  1,2,\ldots,n-1\right\}  \label{pf.thm.elser-coll.cyc1}%
\end{equation}
and%
\begin{equation}
\Phi\left(  B_{n}\right)  \prec B_{1}. \label{pf.thm.elser-coll.cyc2}%
\end{equation}
We shall derive a contradiction.

Set $B_{n+1}=B_{1}$. Then, (\ref{pf.thm.elser-coll.cyc2}) can be rewritten as
$\Phi\left(  B_{n}\right)  \prec B_{n+1}$. In other words, we have
$\Phi\left(  B_{i}\right)  \prec B_{i+1}$ for $i=n$. Combining this with
(\ref{pf.thm.elser-coll.cyc1}), we conclude that%
\begin{equation}
\Phi\left(  B_{i}\right)  \prec B_{i+1}\ \ \ \ \ \ \ \ \ \ \text{for each
}i\in\left\{  1,2,\ldots,n\right\}  . \label{pf.thm.elser-coll.cyc12}%
\end{equation}


Now, set $A_{i}=\operatorname*{Shade}\left(  B_{i}\right)  $ for each
$i\in\left\{  1,2,\ldots,n+1\right\}  $. Then, $A_{n+1}=A_{1}$ (since
$B_{n+1}=B_{1}$).

We now claim the following:

\begin{statement}
\textit{Claim 1:} We have $A_{i}\subseteq A_{i+1}$ for each $i\in\left\{
1,2,\ldots,n\right\}  $.
\end{statement}

[\textit{Proof of Claim 1:} Let $i\in\left\{  1,2,\ldots,n\right\}  $. Then,
the definition of $A_{i}$ yields $A_{i}=\operatorname*{Shade}\left(
B_{i}\right)  $. Likewise, $A_{i+1}=\operatorname*{Shade}\left(
B_{i+1}\right)  $.

We have $B_{i}\in\mathcal{A}_{+}=\left\{  P\in\mathcal{A}\ \mid\ \varepsilon
\left(  P\right)  \in P\right\}  $. In other words, $B_{i}$ is a
$P\in\mathcal{A}$ satisfying $\varepsilon\left(  P\right)  \in P$. In other
words, $B_{i}$ is an element of $\mathcal{A}$ and satisfies $\varepsilon
\left(  B_{i}\right)  \in B_{i}$.

We set $u=\varepsilon\left(  B_{i}\right)  $. The definition of $\Phi$ yields
$\Phi\left(  B_{i}\right)  =B_{i}\setminus\left\{  \varepsilon\left(
B_{i}\right)  \right\}  =B_{i}\setminus\left\{  u\right\}  $ (since
$\varepsilon\left(  B_{i}\right)  =u$).

Recall that $\varepsilon\left(  B_{i}\right)  $ is the smallest edge $e\in
G\setminus\operatorname*{Shade}\left(  B_{i}\right)  $ (by the definition of
$\varepsilon\left(  B_{i}\right)  $). Hence, $\varepsilon\left(  B_{i}\right)
\in G\setminus\operatorname*{Shade}\left(  B_{i}\right)  $. In other words,
$u\in G\setminus\operatorname*{Shade}\left(  B_{i}\right)  $ (since
$u=\varepsilon\left(  B_{i}\right)  $). In other words, $u\in G$ and
$u\notin\operatorname*{Shade}\left(  B_{i}\right)  $. Thus, $u\in G\subseteq
E$ and $u\notin\operatorname*{Shade}\left(  B_{i}\right)  $. Therefore,
(\ref{eq.lem.Shade-tog.diff}) (applied to $F=B_{i}$) yields
$\operatorname*{Shade}\left(  B_{i}\setminus\left\{  u\right\}  \right)
=\operatorname*{Shade}\left(  B_{i}\right)  $. This can be rewritten as
$\operatorname*{Shade}\left(  \Phi\left(  B_{i}\right)  \right)  =A_{i}$
(since $\Phi\left(  B_{i}\right)  =B_{i}\setminus\left\{  u\right\}  $ and
$A_{i}=\operatorname*{Shade}\left(  B_{i}\right)  $).

But (\ref{pf.thm.elser-coll.cyc12}) yields $\Phi\left(  B_{i}\right)  \prec
B_{i+1}$, so that $\Phi\left(  B_{i}\right)  \subseteq B_{i+1}$ and thus
$\operatorname*{Shade}\left(  \Phi\left(  B_{i}\right)  \right)
\subseteq\operatorname*{Shade}\left(  B_{i+1}\right)  $ (by Lemma
\ref{lem.Shade-monoton}, applied to $A=\Phi\left(  B_{i}\right)  $ and
$B=B_{i+1}$). In view of $\operatorname*{Shade}\left(  \Phi\left(
B_{i}\right)  \right)  =A_{i}$ and $\operatorname*{Shade}\left(
B_{i+1}\right)  =A_{i+1}$, this can be rewritten as $A_{i}\subseteq A_{i+1}$.
This proves Claim 1.]

Claim 1 shows that $A_{i}\subseteq A_{i+1}$ for each $i\in\left\{
1,2,\ldots,n\right\}  $. In other words,%
\[
A_{1}\subseteq A_{2}\subseteq\cdots\subseteq A_{n}\subseteq A_{n+1}.
\]
This is a chain of inclusions, but its last entry equals the first: indeed,
$A_{n+1}=A_{1}$. Thus, all inclusions in this chain must be equalities. That
is, we have%
\begin{equation}
A_{1}=A_{2}=\cdots=A_{n}=A_{n+1}. \label{pf.thm.elser-coll.c1eqs}%
\end{equation}
Hence, in particular, $A_{n}=A_{1}$. However, $A_{n}=\operatorname*{Shade}%
\left(  B_{n}\right)  $ (by the definition of $A_{n}$) and $A_{1}%
=\operatorname*{Shade}\left(  B_{1}\right)  $ (by the definition of $A_{1}$).
Hence, $\operatorname*{Shade}\left(  B_{n}\right)  =A_{n}=A_{1}%
=\operatorname*{Shade}\left(  B_{1}\right)  $. Thus,
(\ref{pf.thm.elser-coll.eps=eps}) (applied to $F_{1}=B_{n}$ and $F_{2}=B_{1}$)
yields $\varepsilon\left(  B_{n}\right)  =\varepsilon\left(  B_{1}\right)  $
(since $B_{n}\in\mathcal{A}_{+}\subseteq\mathcal{A}$ and $B_{1}\in
\mathcal{A}_{+}\subseteq\mathcal{A}$).

Set $u=\varepsilon\left(  B_{n}\right)  $. Thus, $u=\varepsilon\left(
B_{n}\right)  =\varepsilon\left(  B_{1}\right)  $.

Recall that the sets $B_{1},B_{2},\ldots,B_{n}$ are distinct. Hence,
$B_{n}\neq B_{1}$ (since $n\geq2$).

The definition of $\Phi$ yields $\Phi\left(  B_{n}\right)  =B_{n}%
\setminus\left\{  \varepsilon\left(  B_{n}\right)  \right\}  =B_{n}%
\setminus\left\{  u\right\}  $ (since $\varepsilon\left(  B_{n}\right)  =u$).
However, (\ref{pf.thm.elser-coll.cyc2}) says that $\Phi\left(  B_{n}\right)
\prec B_{1}$. In other words, $B_{n}\setminus\left\{  u\right\}  \prec B_{1}$
(since $\Phi\left(  B_{n}\right)  =B_{n}\setminus\left\{  u\right\}  $).

On the other hand, $B_{1}\in\mathcal{A}_{+}=\left\{  P\in\mathcal{A}%
\ \mid\ \varepsilon\left(  P\right)  \in P\right\}  $. In other words, $B_{1}$
is a $P\in\mathcal{A}$ satisfying $\varepsilon\left(  P\right)  \in P$. In
other words, $B_{1}$ is an element of $\mathcal{A}$ and satisfies
$\varepsilon\left(  B_{1}\right)  \in B_{1}$. Now, $u=\varepsilon\left(
B_{1}\right)  \in B_{1}$. The same argument (applied to $B_{n}$ instead of
$B_{1}$) yields $u\in B_{n}$ (since $u=\varepsilon\left(  B_{n}\right)  $).
Combining $u\in B_{n}$ with $u\in B_{1}$, we obtain $u\in B_{n}\cap B_{1}$.

Thus, Lemma \ref{lem.elser-coll.lem} (applied to $X=B_{n}$ and $Y=B_{1}$)
yields $B_{n}=B_{1}$ (since $B_{n}\setminus\left\{  u\right\}  \prec B_{1}$).
This contradicts $B_{n}\neq B_{1}$.

Forget that we fixed $\left(  B_{1},B_{2},\ldots,B_{n}\right)  $. We thus have
found a contradiction whenever $\left(  B_{1},B_{2},\ldots,B_{n}\right)  $ is
a tuple of distinct sets $B_{1},B_{2},\ldots,B_{n}\in\mathcal{A}_{+}$ with the
property that $n\geq2$ and that (\ref{pf.thm.elser-coll.cyc1}) and
(\ref{pf.thm.elser-coll.cyc2}) hold. Hence, there exists no such tuple. In other
words, the complete matching $\left(  \mathcal{A}_{-},\mathcal{A}_{+}%
,\Phi\right)  $ is acyclic. Therefore, the simplicial complex $\mathcal{A}$
has an acyclic complete matching, and thus is collapsible (by Definition
\ref{def.matchings} \textbf{(c)}). This finishes the proof of Theorem
\ref{thm.elser-coll}.
\end{proof}

The analogues of Proposition \ref{prop.elser-sc} and of Theorem
\ref{thm.elser-coll} for vertex-infection (instead of usual infection) also
hold (with the same proofs). More generally, Proposition \ref{prop.elser-sc}
and Theorem \ref{thm.elser-coll} can be generalized to any
inclusion-preserving shade map:

\begin{theorem}
\label{thm.Shade.coll}Let $E$ be any finite set. Let $\operatorname*{Shade}%
:\mathcal{P}\left(  E\right)  \rightarrow\mathcal{P}\left(  E\right)  $ be an
inclusion-preserving shade map on $E$. Let $G$ be any subset of $E$. Let%
\[
\mathcal{A}=\left\{  F\subseteq E\ \mid\ G\not \subseteq \operatorname*{Shade}%
F\right\}  .
\]
Then:

\textbf{(a)} This $\mathcal{A}$ is a simplicial complex on ground set $E$.

\textbf{(b)} This simplicial complex $\mathcal{A}$ is collapsible.
\end{theorem}

\begin{proof}
[Proof sketch.]Part \textbf{(a)} is a straightforward generalization of
Proposition \ref{prop.elser-sc}, while part \textbf{(b)} is a straightforward
generalization of Theorem \ref{thm.elser-coll}. The proofs we gave above
generalize (mutatis mutandis).
\end{proof}

However, Theorem \ref{thm.Shade.coll} cannot be lifted to the full generality
of arbitrary shade maps, since $\mathcal{A}$ will generally not be a
simplicial complex unless the shade map is inclusion-preserving. (However, for
inclusion-reversing shade maps, we can obtain a variant of Theorem
\ref{thm.Shade.coll} by applying Theorem \ref{thm.Shade.coll} to the dual
shade map $\operatorname*{Shade}\nolimits^{\prime}$ from Proposition
\ref{prop.Shade.dual}.)

\section{Open questions}

I shall now comment on two natural directions of research so far unexplored.

\subsection{The Alexander dual}

Any simplicial complex has an \emph{Alexander dual}, which is defined as follows:

\begin{definition}
Let $E$ be a finite set. Let $\mathcal{A}$ be a simplicial complex on ground
set $E$. Then, we define a new simplicial complex $\mathcal{A}^{\vee}$ on
ground set $E$ by%
\[
\mathcal{A}^{\vee}=\left\{  F\subseteq E\ \mid\ E\setminus F\notin%
\mathcal{A}\right\}  .
\]
(That is, $\mathcal{A}^{\vee}$ consists of those subsets of $E$ whose
complements don't belong to $\mathcal{A}$.) This simplicial complex
$\mathcal{A}^{\vee}$ is called the \emph{Alexander dual} of $\mathcal{A}$.
\end{definition}

It is well-known that a simplicial complex $\mathcal{A}$ and its Alexander
dual $\mathcal{A}^{\vee}$ share many properties; in particular, the reduced
homology of $\mathcal{A}$ is isomorphic to the reduced cohomology of
$\mathcal{A}^{\vee}$ in opposite degrees
(see, e.g., \cite[Theorem 1.1]{BjoTan09}). However, the
collapsibility and the homotopy types of $\mathcal{A}$ and $\mathcal{A}^{\vee
}$ are not always related. Thus, the following question is suggested but not
answered by Theorem \ref{thm.elser-coll}:

\begin{question}
\label{quest.elser-dual}Let us use the notations from Section \ref{sec.elser}
as well as Definition \ref{def.Shade}. Let $G$ be any subset of $E$. Define
$\mathcal{A}$ as in (\ref{eq.prop.elser-sc.A=}). Is the simplicial complex
\[
\mathcal{A}^{\vee}=\left\{  F\subseteq E\ \mid\ G\subseteq
\operatorname*{Shade}\left(  E\setminus F\right)  \right\}
\]
collapsible? Is it contractible?
\end{question}

\subsection{Several vertices $v$}

Elser's nuclei-based viewpoint in \cite{Elser84} (and \cite[Conjecture
9.1]{DHLetc19}) suggests yet another question.

Our definition of $\operatorname*{Shade}F$ (Definition \ref{def.Shade}), and
the underlying notion of \textquotedblleft infecting\textquotedblright\ an
edge, implicitly relied on the choice of vertex $v$. It thus is advisable to
rename the set $\operatorname*{Shade}F$ as $\operatorname*{Shade}%
\nolimits_{v}F$ and combine such sets for different values of $v$. In
particular, we can define%
\[
\mathcal{A}_{U}=\left\{  F\subseteq E\ \mid\ G\not \subseteq
\operatorname*{Shade}\nolimits_{v}F\text{ for some }v\in U\right\}
\]
for any subset $U$ of $V$. This $\mathcal{A}_{U}$ is a simplicial complex
(being the union of a family of simplicial complexes), and thus we can ask the
same questions about it as we did about $\mathcal{A}$:

\begin{question}
\label{quest.elser-multi}What can we say about the homotopy and discrete Morse
theory of $\mathcal{A}_{U}$ ? What about its Alexander dual?
\end{question}

For $G=E$ and $\left\vert U\right\vert >0$, this simplicial complex
$\mathcal{A}_{U}$ is the Alexander dual of the \textquotedblleft$U$-nucleus
complex\textquotedblright\ $\Delta_{U}^{G}$ from \cite[Definition
3.2]{DHLetc19} (when $G$ is connected). If \cite[Conjecture 9.1 for
$\left\vert U\right\vert >1$]{DHLetc19} is correct, then the homology of
$\mathcal{A}_{U}$ with real coefficients should be concentrated in a single
degree; this suggests the possible existence of an acyclic partial matching
with all critical faces in one degree.

\subsection*{Acknowledgments}

I thank Anders Bj\"{o}rner, Galen Dorpalen-Barry, Dmitry Feichtner-Kozlov,
Patricia Hersh, Sam Hopkins, Vic Reiner, Tom Roby and Richard Stanley for
insightful conversations. A referee has greatly simplified the proof of
Theorem \ref{thm.elser-coll}. This research was supported through the
programme \textquotedblleft Oberwolfach Leibniz Fellows\textquotedblright\ by
the Mathematisches Forschungsinstitut Oberwolfach in 2020. I am deeply
grateful to the Institute for its hospitality during a none-too-hospitable time.

\begin{thebibliography}{99999999}                                                                                         %


\bibitem[BenQui08]{BenQui08}%
\href{https://scholarship.claremont.edu/cgi/viewcontent.cgi?referer=&httpsredir=1&article=1581&context=hmc_fac_pub}{Arthur
T. Benjamin and Jennifer J. Quinn, \textit{An Alternate Approach to
Alternating Sums: A Method to DIE for}, The College Mathematics Journal,
Volume 39, Number 3, May 2008, pp. 191-202(12).}

\bibitem[BjoTan09]{BjoTan09}%
\href{https://doi.org/10.1007/s00454-008-9102-x}{Anders Bj\"{o}rner, Martin
Tancer, \textit{Combinatorial Alexander Duality -- A Short and Elementary
Proof}, Discrete Comput Geom (2009) 42, pp. 586--593.}

\bibitem[BrOlNo09]{BrOlNo09}%
\href{https://doi.org/10.1016/j.ejc.2008.12.010}{Josep M. Brunat, Ant\'{o}nio
Guedes de Oliveira, Marc Noy, \textit{Partitions of a finite Boolean lattice
into intervals}, European Journal of Combinatorics \textbf{30} (2009), pp.
1801--1809}.

\bibitem[Dawson80]{Dawson80}\href{https://doi.org/10.1007/BFb0091815}{J. E.
Dawson, \textit{A construction for a family of sets and its application to
matroids}, in: Combinatorial Mathematics VIII (Geelong, 1980), in: L.L.
MacAvaney (Ed.), Lect. Notes Math., vol. \textbf{884}, Springer, Berlin-New
York, 1981, pp. 136--147}.

\bibitem[DedTit20]{DedTit20}\href{https://arxiv.org/abs/2008.02006v2}{Kristina
Dedndreaj, Peter Tittmann, \textit{On the Activities and Partitions of the
Vertex Subsets of Graphs}, arXiv:2008.02006v2}.

\bibitem[DHLetc19]{DHLetc19}\href{https://arxiv.org/abs/1905.11330v5}{Galen
Dorpalen-Barry, Cyrus Hettle, David C. Livingston, Jeremy L. Martin, George
Nasr, Julianne Vega, Hays Whitlatch, \textit{A positivity phenomenon in
Elser's Gaussian-cluster percolation model}, arXiv:1905.11330v5},
\href{https://doi.org/10.1016/j.jcta.2020.105364}{Journal of Combinatorial
Theory, Series A, \textbf{179}:105364, April 2021,
doi:10.1016/j.jcta.2020.105364}.

\bibitem[EdeJam85]{EdeJam85}\href{https://doi.org/10.1007/BF00149365}{Paul H.
Edelman, Robert E. Jamison, \textit{The theory of convex geometries},
Geometriae Dedicata \textbf{19} (1985), pp. 247--270}.

\bibitem[Elser84]{Elser84}%
\href{https://iopscience.iop.org/article/10.1088/0305-4470/17/7/019/pdf}{Veit
Elser, \textit{Gaussian-cluster models of percolation and self-avoiding
walks}, J. Phys. A: Math. Gen. \textbf{17} (1984), pp. 1515--1523}.

\bibitem[GorMah97]{GorMah97}\href{https://doi.org/10.1006/aama.1996.0495}{Gary
Gordon, Elizabeth McMahon, \textit{Interval Partitions and Activities for the
Greedoid Tutte Polynomial}, Advances in Applied Mathematics \textbf{18}, Issue
1, January 1997, pp. 33--49}.

\begin{vershort}


\bibitem[Grinbe20]{verlong}%
\href{http://www.cip.ifi.lmu.de/~grinberg/algebra/elsersum-long.pdf}{Darij
Grinberg, \textit{The Elser nuclei sum revisited}, detailed version of the
present paper}. Also available as an ancillary file at
\href{http://arxiv.org/abs/2009.11527v8}{arXiv:2009.11527v8}.
\end{vershort}

\begin{verlong}


\bibitem[Grinbe20]{vershort}%
\href{http://www.cip.ifi.lmu.de/~grinberg/algebra/elsersum.pdf}{Darij
Grinberg, \textit{The Elser nuclei sum revisited}, standard version of the
present paper}. Also available as
\href{http://arxiv.org/abs/2009.11527v8}{arXiv:2009.11527v8}.
\end{verlong}

\bibitem[KoLoSc91]{KoLoSc91}Bernhard Korte, L\'{a}szl\'{o} Lov\'{a}sz, Rainer
Schrader, \textit{Greedoids}, Algorithms and Combinatorics \#4, Springer
1991.\newline\url{https://doi.org/10.1007/978-3-642-58191-5}

\bibitem[Kozlov20]{Kozlov20}\href{https://bookstore.ams.org/gsm-207}{Dmitry N.
Kozlov, \textit{Organized Collapse: An Introduction to Discrete Morse Theory},
Graduate Studies in Mathematics \textbf{207}, AMS 2020}. \\
See \url{https://dfklab.com/wp-content/uploads/2021/01/dmt.pdf} for a preprint.

\bibitem[Sagan20]{Sagan20}\href{https://bookstore.ams.org/gsm-210/}{Bruce
Sagan, \textit{Combinatorics: The Art of Counting}, Graduate Studies in
Mathematics \textbf{210}, AMS 2020}.\newline See
\url{https://users.math.msu.edu/users/bsagan/Books/Aoc/final.pdf} for a preprint.

\bibitem[Stanle11]{Stanley-EC1}Richard P. Stanley, \textit{Enumerative
Combinatorics, volume 1}, 2nd edition, version of 15 July 2011.\newline%
\url{http://math.mit.edu/~rstan/ec/} .\newline See
\url{http://math.mit.edu/~rstan/ec/} for errata.
\end{thebibliography}


\end{document}
